% ============================================================
%  CUADERNO DE ÁLGEBRA PREUNIVERSITARIA
%  Transcripción a LaTeX para Overleaf
%  Compilador: pdfLaTeX (el predeterminado de Overleaf)
% ============================================================
\documentclass[11pt,a4paper]{article}

% ---------- Idioma y codificación ----------
\usepackage[utf8]{inputenc}
\usepackage[T1]{fontenc}
\usepackage[spanish,es-noshorthands,es-nodecimaldot]{babel}

% ---------- Página ----------
\usepackage[margin=1.8cm]{geometry}
\setlength{\parindent}{0pt}
\setlength{\parskip}{4pt}

% ---------- Matemáticas ----------
\usepackage{amsmath,amssymb}
\allowdisplaybreaks

% ---------- Listas, columnas, cajas ----------
\usepackage{enumitem}
\usepackage{multicol}
\usepackage{fancybox}      % \ovalbox (cajas ovaladas como en el cuaderno)

% ---------- Colores ----------
\usepackage[dvipsnames]{xcolor}
\definecolor{rojo}{RGB}{200,30,45}
\definecolor{azul}{RGB}{30,70,180}
\definecolor{verde}{RGB}{20,130,60}
\definecolor{naranja}{RGB}{225,115,0}
\definecolor{graf}{RGB}{120,40,160}      % color de las gráficas
\definecolor{resaltado}{RGB}{160,0,120}  % lo que estaba resaltado en el cuaderno

% ---------- Gráficos ----------
\usepackage{tikz}
\usetikzlibrary{arrows.meta,patterns}

% ---------- Títulos ----------
\usepackage{titlesec}
\titleformat{\section}{\Large\bfseries\color{rojo}}{\thesection.}{0.6em}{}
\titleformat{\subsection}{\large\bfseries\color{azul}}{\thesubsection}{0.5em}{}



% ============================================================
%  MACROS PROPIAS
% ============================================================
% Número encerrado en círculo: \cn{1}
\newcommand{\cn}[1]{\raisebox{.5pt}{\textcircled{\raisebox{-.9pt}{\scriptsize #1}}}}

% Caja de color para fórmulas importantes
\newcommand{\caja}[1]{\fcolorbox{azul}{white}{#1}}

% Texto que estaba resaltado en el cuaderno
\newcommand{\hl}[1]{\textcolor{resaltado}{#1}}

% Aspa simple:  \aspasimple{a}{b}{c}{d}
\newcommand{\aspasimple}[4]{%
\begin{tikzpicture}[baseline=-3pt]
  \node (a) at (0,0.35) {$#1$}; \node (b) at (1.6,0.35) {$#2$};
  \node (c) at (0,-0.35) {$#3$}; \node (d) at (1.6,-0.35) {$#4$};
  \draw[rojo] (a.east) -- (d.west); \draw[rojo] (c.east) -- (b.west);
\end{tikzpicture}}

% Aspa doble: \aspadoble{a1}{b1}{c1}{a2}{b2}{c2}
\newcommand{\aspadoble}[6]{%
\begin{tikzpicture}[baseline=-3pt]
  \node (a1) at (0,0.35) {$#1$}; \node (b1) at (1.6,0.35) {$#2$}; \node (c1) at (3.2,0.35) {$#3$};
  \node (a2) at (0,-0.35) {$#4$}; \node (b2) at (1.6,-0.35) {$#5$}; \node (c2) at (3.2,-0.35) {$#6$};
  \draw[rojo] (a1.east) -- (b2.west); \draw[rojo] (a2.east) -- (b1.west);
  \draw[azul] (b1.east) -- (c2.west); \draw[azul] (b2.east) -- (c1.west);
\end{tikzpicture}}

% Intervalos: \intervalo{o|c}{o|c}{f|i}  (o = abierto, c = cerrado, f = finito, i = infinito)
\newcommand{\intervalo}[3]{%
\begin{tikzpicture}[baseline=-3pt]
  \draw[<->] (-0.4,0) -- (3.4,0);
  \if i#3
    \draw[<-,thick,graf] (3.3,0.25) -- (0.8,0.25) -- (0.8,0);
    \node[below] at (0.8,-0.05) {\small $a$};
    \node[below] at (3.2,-0.05) {\small $+\infty$};
  \else
    \draw[thick,graf] (0.8,0) -- (0.8,0.25) -- (2.2,0.25) -- (2.2,0);
    \node[below] at (0.8,-0.05) {\small $a$}; \node[below] at (2.2,-0.05) {\small $b$};
    \if c#2 \fill[graf] (2.2,0) circle (2pt); \else \draw[graf,fill=white] (2.2,0) circle (2pt); \fi
  \fi
  \if c#1 \fill[graf] (0.8,0) circle (2pt); \else \draw[graf,fill=white] (0.8,0) circle (2pt); \fi
\end{tikzpicture}}

% Recta de signos: \rectasignos{p1/o,p2/c,...}{+,-,+,...}{número de zonas}
\newcommand{\rectasignos}[3]{%
\begin{tikzpicture}[baseline=-3pt]
  \draw[<->] (0,0) -- ({1.2*#3},0);
  \foreach [count=\i] \p/\t in {#1} {
    \if c\t \fill (1.2*\i,0) circle (2pt); \else \draw[fill=white] (1.2*\i,0) circle (2pt); \fi
    \node[below] at (1.2*\i,-0.08) {\small $\p$};
  }
  \foreach [count=\i] \s in {#2} \node[above,rojo] at ({1.2*\i-0.6},0.02) {$\s$};
\end{tikzpicture}}

% Ejes cartesianos: \ejes{xmin}{xmax}{ymin}{ymax}
\newcommand{\ejes}[4]{%
  \draw[->] (#1,0) -- (#2,0) node[below] {$x$};
  \draw[->] (0,#3) -- (0,#4) node[left] {$y$};}

% Parábola pequeña para la tabla de discriminantes: \parabola{signo de a}{desplazamiento}{-}
\newcommand{\parabola}[3]{%
\begin{tikzpicture}[scale=0.45,baseline=0pt]
  \draw[->] (-2,0) -- (2.2,0); \draw[->] (0,-1.7) -- (0,1.9);
  \draw[thick,graf,domain=-1.4:1.4,samples=40] plot (\x,{(#1)*0.6*\x*\x+(#2)*0.6});
\end{tikzpicture}}

% Función "rampa" usada en desplazamientos: \rampa{x0}{y0}{estilo}
\newcommand{\rampa}[3]{%
  \draw[thick,#3] ({#1-0.6},{#2-0.6}) -- ({#1+1},{#2+1}) -- ({#1+2.2},{#2+1});}

% ============================================================
\begin{document}

\begin{center}
{\Huge\bfseries\color{rojo} Cuaderno de Álgebra}\\[6pt]
{\Large\color{azul} Preuniversitaria}
\end{center}

\tableofcontents
\newpage
% ============================================================
% Páginas 1-3
% ============================================================
\section{Teoría de exponentes}

Estudia todas las clases o tipos de exponentes y las diferentes relaciones que
existen entre ellos a través de leyes matemáticas.

\subsection*{1. Propiedades básicas}

\begin{multicols}{2}
\begin{itemize}[label=$\ast$]
  \item $a^0 = 1,\ \forall a \neq 0$
  \item $a^1 = a,\ \forall a \in \mathbb{R}$
  \item $a^{-n} = \left(\dfrac{1}{a}\right)^{n},\ \forall a \neq 0$
  \item $(a^m)^n = a^{m\cdot n}$
  \item $(a\cdot b)^n = a^n\cdot b^n$
  \item $a^m\cdot a^n = a^{m+n}$
  \item $\dfrac{a^m}{a^n} = a^{m-n}$
  \item $a^n = a^m \;\rightarrow\; n = m$
  \item $\left(\dfrac{a}{b}\right)^{n} = \dfrac{a^n}{b^n},\ \forall a\in\mathbb{R},\ \forall b\neq 0$
  \item $\sqrt[2n]{a^{2n}} = |a|,\ \forall a \in \mathbb{R}$
  \item $\sqrt[2n+1]{a^{2n+1}} = a,\ \forall a \in \mathbb{R}$
  \item $\sqrt[n]{a^m} = a^{\frac{m}{n}}$
  \item $\sqrt[n]{\dfrac{a}{b}} = \dfrac{\sqrt[n]{a}}{\sqrt[n]{b}},\ \forall a\in\mathbb{R},\ \forall b\neq 0$
  \item $\sqrt[n]{a}\cdot\sqrt[n]{b} = \sqrt[n]{a\cdot b}$
  \item $\sqrt[n]{\sqrt[m]{a}} = \sqrt[n\cdot m]{a}$
  \item $\sqrt[m]{a\sqrt[n]{b\sqrt[p]{c}}} = \sqrt[m]{a}\cdot\sqrt[m\cdot n]{b}\cdot\sqrt[m\cdot n\cdot p]{c}$
\end{itemize}
\end{multicols}

\subsection*{2. Propiedades adicionales}

\begin{itemize}[label=$\ast$]
  \item $\sqrt[m]{a^n\sqrt[p]{a^q\sqrt[r]{a^s}}} = \sqrt[m\cdot p\cdot r]{a^{(n\cdot p+q)\cdot r+s}}$
  \item $\sqrt[m]{a^n\div\sqrt[p]{a^q\div\sqrt[r]{a^s}}} = \sqrt[m\cdot p\cdot r]{a^{(n\cdot p-q)r+s}}$
  \item $\sqrt{a+\sqrt{a+\sqrt{a+\cdots\infty}}} = \dfrac{\sqrt{4a+1}+1}{2}$
  \item $\sqrt{a-\sqrt{a-\sqrt{a-\cdots\infty}}} = \dfrac{\sqrt{4a+1}-1}{2}$
  \item $\sqrt{a\sqrt{a\sqrt{a\cdots\infty}}} = a$
  \item $\sqrt[n]{a^m\cdot\sqrt[n]{a^m\cdots\infty}} = \sqrt[n-1]{a^m}$
  \item $\sqrt[n]{a^m\div\sqrt[n]{a^m\div\cdots\infty}} = \sqrt[n+1]{a^m}$
  \item $\underbrace{\sqrt[n]{a^{n-1}\sqrt[n]{a^{n-1}\cdots}}}_{\text{``$k$'' veces}} = \sqrt[n^k]{a^{n^k-1}}$
  \item $\underbrace{\sqrt[n]{a^{p}\sqrt[n]{a^{p}\cdots}}}_{\text{``$m$'' veces}} = \sqrt[n^m]{a^{\frac{p(n^m-1)}{n-1}}}$
  \item Si ``$m$'' es un número par:
  \[\underbrace{\sqrt[n]{a^{p}\div\sqrt[n]{a^{p}\cdots}}}_{\text{``$m$'' veces}} = \sqrt[n^m]{a^{\frac{p(n^m-1)}{n+1}}}\]
  \item Si ``$m$'' es un número impar:
  \[\underbrace{\sqrt[n]{a^{p}\div\sqrt[n]{a^{p}\cdots}}}_{\text{``$m$'' veces}} = \sqrt[n^m]{a^{\frac{p(n^m+1)}{n+1}}}\]
\end{itemize}

\subsection*{3. Propiedades auxiliares}

\begin{itemize}[label=$\ast$]
  \item Si ``$a = m(m+1)$'', se cumple que:
  \begin{itemize}[label=\checkmark]
    \item $\sqrt{a+\sqrt{a+\cdots\infty}} = m+1$
    \item $\sqrt{a-\sqrt{a-\cdots\infty}} = m$
  \end{itemize}
  \item Si ``$n \in \mathbb{R}-\{0\}$'', se cumple que:
  \begin{itemize}[label=\checkmark]
    \item $x^{x^{x^{\cdot^{\cdot^{\cdot\infty}}}}} = n \;\rightarrow\; x = \sqrt[n]{n}$
    \item $x^{x^{x^{\cdot^{\cdot^{\cdot n}}}}} = n \;\rightarrow\; x = \sqrt[n]{n}$
    \item $x^{-n^{n^{-x}}} = n \;\rightarrow\; \sqrt[x]{x} = n^{-1}$
    \item $x^{\frac{x-1}{x}} = n^{n-1} \;\rightarrow\; x = n^{-1}$
    \item $\sqrt[n]{n}^{\sqrt[n]{n}^{\cdot^{\cdot^{\cdot\infty}}}} = n$; no se cumple para todo $n$.
  \end{itemize}
  \item Si ``$a, b$ y $c \neq 0$'', se cumple que:
  \begin{itemize}[label=\checkmark]
    \item $\sqrt[n]{\dfrac{(ab)^n+(bc)^n+(ac)^n}{a^{-n}+b^{-n}+c^{-n}}} = a\cdot b\cdot c$
    \item $\sqrt[n]{\dfrac{a^n+b^n+c^n}{(ab)^{-n}+(bc)^{-n}+(ac)^{-n}}} = a\cdot b\cdot c$
  \end{itemize}
\end{itemize}

\newpage
% ------------------------------------------------------------
\section{Polinomios}

\textbf{El polinomio:} son expresiones algebraicas racionales enteras.

\textbf{Ejemplos:}
\begin{enumerate}[label=\alph*)]
  \item $E(x;y) = 3x^{20}y^{7}$ \quad en $\mathbb{Z}$.
  \item $F(a;b;c) = -\dfrac{1}{2}a^{5} + 6b^{3}c^{2}$ \quad en $\mathbb{Q}$.
  \item $P(x) = \sqrt{5}\,x^{12} - 7x^{4} + 8x^{7}$ \quad en $\mathbb{R}$.
  \item $F(y) = (\sqrt{y})^{4} + 8y + 11$ \quad \textcolor{rojo}{no es polinomio}.
\end{enumerate}

\subsection*{Términos semejantes}
\begin{enumerate}[label=\alph*)]
  \item $6x^4;\ -\sqrt{5}\,x^4;\ 2x^4$
  \item $11a^2b^3;\ -\dfrac{1}{2}a^2b^3$
  \item Si $6x^4 + mx^a = 20x^4$, \ $am = \text{¿?}$
  \[ m = 14 \qquad a = 4 \qquad a\times m = 56 \]
\end{enumerate}

\subsection*{Valor numérico}

\textbf{\textcircled{1}} $P(x) = 4x^2 - 7x + 11$, \quad $x = -1$
\[
\begin{array}{c|cc|c}
       & 4 & -7  & 11 \\
x = -1 &   & -4  & 11 \\ \hline
       & 4 & -11 & \boxed{22}
\end{array}
\quad \text{V.N.}
\]

\textbf{\textcircled{2}} $P(x) = ax + b$
\begin{align*}
P(P(x)) &= a^2x + ab + b\\
P(P(P(x))) &= a^3x + a^2b + ab + b\\
\underbrace{P(P(\cdots(P(x))\cdots))}_{\text{``$n$'' veces}} &= a^nx + a^{n-1}b + a^{n-2}b + \cdots + ab + b
\end{align*}

\textbf{\textcircled{3}} $P(x-3) = 5x + 11$, \quad $P(x) = \text{¿?}$
\begin{align*}
x - 3 = y \;&\rightarrow\; x = y + 3\\
P(y) &= 5(y+3) + 11\\
P(y) &= 5y + 26
\end{align*}
\textcolor{azul}{Cambiamos $y$ por $x$:} \quad $P(x) = 5x + 26$


% ============================================================
% Páginas 4-6
% ============================================================
\subsection*{Teoremas}
Sea un polinomio $P(x)$, se cumple que:
\begin{enumerate}[label=\Roman*.]
  \item \textcolor{magenta}{$\sum \text{coeficientes } P(x) = P(1)$}
  \item \textcolor{magenta}{Término independiente $= P(0)$}
\end{enumerate}

\textbf{Ejemplos:}

\begin{minipage}[t]{0.485\textwidth}
a) $P(x) = (x-1)^{2020} + 7x + 1$
\[P(1) = (1-1)^{2020} + 7(1) + 1\]
\[\textstyle\sum \text{coef.} = 8\]
\end{minipage}\hfill
\begin{minipage}[t]{0.485\textwidth}
b) $Q(y-1) = 2y + 11$
\[Q(0) = 2(1) + 11\]
\[T_i = 13\]
\end{minipage}

\subsection*{Polinomio con una sola variable}
Sea: \quad $P(x) \equiv a_0x^n + a_1x^{n-1} + \cdots + a_n$
\begin{itemize}[label={}]
  \item Variable: $x$
  \item Grado: $n$
  \item Coeficiente principal: $a_0$
  \item Término independiente: $a_n$
\end{itemize}

\textbf{Ejemplos:}

\begin{minipage}[t]{0.320\textwidth}
$P(x) = 6x - 7$
\begin{itemize}[label={}]
  \item V: $x$
  \item G: $1$ (Lineal)
  \item CP: $6$
  \item Ti: $-7$
\end{itemize}
\end{minipage}\hfill
\begin{minipage}[t]{0.359\textwidth}
$F(x) = x^2 - 7x + 11$
\begin{itemize}[label={}]
  \item V: $x$
  \item G: $2$ (Cuadrático)
  \item CP: $1$ (Mónico)
  \item Ti: $11$
\end{itemize}
\end{minipage}\hfill
\begin{minipage}[t]{0.291\textwidth}
$M(x) \equiv 10$
\begin{itemize}[label={}]
  \item Grado cero
  \item $M(2) = 10$
  \item $M(-1) = 10$
\end{itemize}
\end{minipage}

\subsection*{Grado}

\begin{minipage}[t]{0.485\textwidth}
\textbf{I. En un término:}
\[M(x;y) = 2x^{11}y^{7}\]
\begin{itemize}[label={}]
  \item $GR_x: 11$
  \item $GR_y: 7$
  \item $GA: 18$
\end{itemize}
\end{minipage}\hfill
\begin{minipage}[t]{0.485\textwidth}
\textbf{II. En dos o más términos:}
\[P(x;y) = 3\overbrace{x^{4}y^{6}}^{10} - \overbrace{x^{2}y^{11}}^{13}\]
\begin{itemize}[label={}]
  \item $GR_x: 4$
  \item $GR_y: 11$
  \item $GA: 13$
\end{itemize}
\end{minipage}

\subsection*{Polinomios especiales}

\textbf{I. Homogéneo:}
\[P(x,y) = 2\underbrace{x^{13}y^{4}}_{17} - \underbrace{x^{10}y^{7}}_{17} + \underbrace{x^{17}}_{17}\]

\textbf{II. Completo:}

\begin{minipage}[t]{0.485\textwidth}
\[P(x) = 7x^2 - 11x + 4\]
\centering Completo y ordenado\\ de forma descendente
\end{minipage}\hfill
\begin{minipage}[t]{0.485\textwidth}
\[Q(y) = 2 - 6y + 8y^2 + y^3\]
\centering Completo y ordenado\\ de forma ascendente
\end{minipage}

\[\#\text{ de términos} = \text{grado} + 1\]

\begin{minipage}[t]{0.485\textwidth}
\textbf{III. Idénticos:}
\[ax + b \equiv mx + n\]
\[a = m \;\wedge\; b = n\]
\textcolor{rojo}{Además:}
\begin{align*}
&(x+1)(x+2) \equiv x(x+3) + 2\\
&x = 2 \;\rightarrow\; 3\times 4 = 2\times 5 + 2\\
&x = -1 \;\rightarrow\; 0\times 1 = -1\times 2 + 2
\end{align*}
\end{minipage}\hfill
\begin{minipage}[t]{0.485\textwidth}
\textbf{IV. Nulo:}
\[ax^2 + bx + c \equiv 0\]
\[a = 0,\quad b = 0,\quad c = 0\]
\end{minipage}

\medskip
\textcolor{naranja}{\textbf{Obs.:}} Si $P(x,y)$ es homogéneo de grado ``$n$''
$\;\longrightarrow\; \boxed{P(\alpha x, \alpha y) = \alpha^n P(x,y)}$

\newpage
% ------------------------------------------------------------
\section{Productos notables}

\begin{multicols}{2}
\textbf{\cn{1} Identidades:}
\begin{itemize}[label=\checkmark]
  \item $(a+b)^2 = a^2 + 2ab + b^2$
  \item $(a-b)^2 = a^2 - 2ab + b^2$
\end{itemize}

\textbf{\cn{2} Legendre:}
\begin{itemize}[label=\checkmark]
  \item $(a+b)^2 + (a-b)^2 = 2(a^2+b^2)$
  \item $(a+b)^2 - (a-b)^2 = 4ab$
\end{itemize}

\textbf{\cn{3}}
\begin{itemize}[label=\checkmark]
  \item $(a+b)(a-b) = a^2 - b^2$
\end{itemize}

\textbf{\cn{4} Stevin:}
\begin{itemize}[label=\checkmark]
  \item $(x+a)(x+b) = x^2 + (a+b)x + ab$
\end{itemize}

\textbf{\cn{5} Cauchy:}
\begin{itemize}[label=\checkmark]
  \item $(a+b)^3 = a^3 + b^3 + 3ab(a+b)$
  \item $(a-b)^3 = a^3 - b^3 - 3ab(a-b)$
\end{itemize}

\textbf{\cn{6}}
\begin{itemize}[label=\checkmark]
  \item $(a+b)(a^2 - ab + b^2) = a^3 + b^3$
  \item $(a-b)(a^2 + ab + b^2) = a^3 - b^3$
\end{itemize}
\end{multicols}

\textbf{\cn{7}} \checkmark\ $(a+b+c)^2 = a^2 + b^2 + c^2 + 2(ab + bc + ac)$, \ Lagrange

\textbf{\cn{8}} $(a^{2m} + a^mb^n + b^{2n})(a^{2m} - a^mb^n + b^{2n}) = a^{4m} + a^{2m}b^{2n} + b^{4n}$

\textbf{\cn{9}} $(a+b+c)^3 = a^3 + b^3 + c^3 + 3(a+b)(a+c)(b+c)$

\subsection*{Relaciones importantes}

\begin{minipage}[t]{0.320\textwidth}
\cn{1} Si: $(a+b)^2 = 4ab$
\[\rightarrow a = b\]
\end{minipage}\hfill
\begin{minipage}[t]{0.320\textwidth}
\cn{2} Si: $\dfrac{1}{a} + \dfrac{1}{b} = \dfrac{4}{a+b}$
\[\rightarrow a = b\]
\end{minipage}\hfill
\begin{minipage}[t]{0.330\textwidth}
\cn{3} $(a+b)^4 - (a-b)^4 = 8ab(a^2+b^2)$
\end{minipage}

\medskip
\begin{minipage}[t]{0.320\textwidth}
\cn{4} Si: $x^2 - 3x + 1 = 0$
\[\rightarrow x + \frac{1}{x} = 3\]
\end{minipage}\hfill
\begin{minipage}[t]{0.650\textwidth}
\cn{5} $\sqrt{A + B \pm 2\sqrt{AB}} = \sqrt{A} \pm \sqrt{B}, \quad A > B$
\end{minipage}

\subsection*{Identidades auxiliares}
\begin{enumerate}[label=\cn{\arabic*}]
  \item $(x+a)(x+b)(x+c) = x^3 + x^2(a+b+c) + x(ab+bc+ac) + abc$
  \item $a^3 + b^3 + c^3 - 3abc = (a+b+c)(a^2+b^2+c^2-ab-bc-ac)$, \ Gauss
  \item $(x^{2n} + x^n + 1)(x^{2n} - x^n + 1) = x^{4n} + x^{2n} + 1$
\end{enumerate}


% ============================================================
% Páginas 7-9
% ============================================================
\subsection*{Identidades condicionales}

\begin{minipage}[t]{0.436\textwidth}
\cn{1} Si: $a + b + c = 0$
\begin{align*}
&\rightarrow a^2 + b^2 + c^2 = -2(ab + bc + ac)\\
&\rightarrow a^3 + b^3 + c^3 = 3abc
\end{align*}
\end{minipage}\hfill
\begin{minipage}[t]{0.533\textwidth}
\cn{2} Sea: $\{a, b, c\} \subset \mathbb{R}$

$\ast$ Si: $a^2 + b^2 + c^2 = ab + bc + ac$
\[\rightarrow a = b = c\]
$\ast$ Si: $a^2 + b^2 + c^2 = 0$
\[\rightarrow a = 0 \wedge b = 0 \wedge c = 0\]
\end{minipage}

\medskip
\cn{3} Si: $a^3 + b^3 + c^3 = 3abc \;\rightarrow\; a = b = c \;\vee\; a + b + c = 0$

\subsection*{Más identidades}
\begin{itemize}[label=$\ast$]
  \item $(a+b+c)^3 = a^3 + b^3 + c^3 - 3abc + 3(a+b+c)(ab+bc+ac)$
  \item $(a+b+c)^3 = 3(a+b+c)(a^2+b^2+c^2) - 2(a^3+b^3+c^3) + 6abc$
  \item $(a+b)^4 + (a-b)^4 = 2(a^4 + 6a^2b^2 + b^4)$
\end{itemize}

\begin{minipage}[t]{0.485\textwidth}
$\ast$ $x - x^{-1} = k,\ \forall k \in \mathbb{R}$
\begin{align*}
x^2 + x^{-2} &= k^2 + 2\\
x^3 - x^{-3} &= k^3 + 3k\\
x^4 + x^{-4} &= (k^2+2)^2 - 2\\
x^6 + x^{-6} &= (k^3+3k)^2 + 2
\end{align*}
\end{minipage}\hfill
\begin{minipage}[t]{0.485\textwidth}
$\ast$ $m + m^{-1} = k,\ \forall k \in \mathbb{R}$
\begin{align*}
m^2 + m^{-2} &= k^2 - 2\\
m^3 + m^{-3} &= k^3 - 3k\\
m^4 + m^{-4} &= (k^2-2)^2 - 2\\
m^6 + m^{-6} &= (k^3-3k)^2 - 2
\end{align*}
\end{minipage}

\newpage
% ------------------------------------------------------------
\section{Factorización}

Factorizar es el proceso que consiste en transformar un polinomio entero y de
coeficientes racionales, en un producto indicado de potencias de sus factores
primos también racionales enteros y de coeficientes racionales.

\begin{enumerate}
  \item \textbf{Factor:} El factor de una expresión es aquel que la divide exactamente.
  \[(x)(x+2) = x^2 + 2x \;\longrightarrow\; (x) \text{ y } (x+2) \text{ son factores de } x^2 + 2x.\]

  \item \textbf{Factor común:} Consiste en identificar el factor que se repite para luego
  retirarlo y anotar en un paréntesis los cocientes de dividir cada término entre dicho factor.
  \[6x^3 + 18x^2 = 6x^2(x+3) \;\longrightarrow\; (6x^2) \text{ es el factor común.}\]

  \item \textbf{Factor primo:} Es aquel que sólo acepta descomponerse en dos factores
  (divisores): la unidad y el propio factor.
  \[(x)(y^2+2y) \;\longrightarrow\; (x) \text{ es primo y } (y^2+2y) \text{ es no primo.}\]
\end{enumerate}

\subsection*{Criterios de factorización}

\textbf{1. Aspa simple:} Consiste en descomponer los términos extremos, de tal manera que
al multiplicar en aspa y sumar estos productos se obtenga el término central.

\textcolor{naranja}{Sirve para:} \
\fbox{$ax^{2n} + bx^n + c$} \quad \fbox{$ax^{2n} + bx^ny^m + cy^{2m}$}

\medskip
$\bullet\;\; x^2 - 8x + 15$ \qquad \aspasimple{x}{-5}{x}{-3}
\qquad $\therefore\; x^2 - 8x + 15 = (x-5)(x-3)$

\bigskip
\textbf{2. Aspa doble:} Consiste en descomponer los términos que poseen grado par y/o el
término independiente, a fin de generar tres aspas simples que verificarán los términos
restantes.

\textcolor{naranja}{Sirve para:}
\begin{center}
\fbox{$ax^{2n} + bx^ny^m + cy^{2m} + dx^n + ey^m + f$}\\[4pt]
\fbox{$ax^{2n} + bx^ny^m + cy^{2m} + dx^nz^p + ey^mz^p + fz^{2p}$}
\end{center}

$\bullet\;\; x^2 + 5xy + 4y^2 + 3x + 6y + 2$
\begin{center}
\aspadoble{x}{4y}{2}{x}{y}{1}
\end{center}
\[\therefore\; x^2 + 5xy + 4y^2 + 3x + 6y + 2 = (x + 4y + 2)(x + y + 1)\]

\bigskip
\textbf{3. Aspa doble especial:} Consiste en descomponer los términos extremos para que a
través de un aspa simple generemos un resultado. Este restará al término central y el
valor final de esta operación se va a descomponer debajo del mismo a fin de generar dos
aspas simples que verificarán los términos restantes.

\textcolor{naranja}{Sirve para:}
\begin{center}
\fbox{$ax^{4n} + bx^{3n} + cx^{2n} + dx^n + e$} \quad
\fbox{$ax^{4n} + bx^{3n}y^m + cx^{2n}y^{2m} + dx^ny^{3m} + ey^{4m}$}
\end{center}

$\bullet\;\; 6x^4 - 13x^3 + 7x^2 + 6x - 8$

\begin{center}
\aspasimple{3x^2}{4}{2x^2}{-2}
\qquad$\longrightarrow$\qquad
\aspadoble{3x^2}{-5x}{4}{2x^2}{-x}{-2}
\end{center}

Verificaciones: $-6x^2 + 8x^2 = 2x^2$\\
Entonces: $7x^2 - 2x^2 = 5x^2$
\[\therefore\; 6x^4 - 13x^3 + 7x^2 + 6x - 8 = (3x^2 - 5x + 4)(2x^2 - x - 2)\]

\bigskip
\textbf{4. Ruffini:} Método usado para factorizar \hl{polinomios completos y ordenados de una
variable, grado impar y que presenten por lo menos una solución real.} Consiste en hallar
los divisores del término independiente, de tal manera que al elegir uno de ellos y
reemplazarlo en el polinomio a factorizar, el resultado sea $0$. Luego de encontrar ese
valor, se utilizará para empezar la factorización como se muestra a continuación.

\begin{minipage}[t]{0.485\textwidth}
\[P(x) = x^3 + 6x^2 + 3x - 10\]
Calculamos todos los divisores del $T_i$:
\[\#D(10) = \pm(1, 2, 5, 10)\]
Elegimos el valor correspondiente:
\[x = 1 \longrightarrow P(1) = 1^3 + 6\cdot1^2 + 3\cdot1 - 10 = 0\]
\end{minipage}\hfill
\begin{minipage}[t]{0.485\textwidth}
\[
\begin{array}{c|ccc|c}
x - 1 = 0 & 1 & 6 & 3  & -10\\
x = 1     &   & 1 & 7  & 10\\ \hline
          & 1 & 7 & 10 & \boxed{0}
\end{array}
\]
\[Q(x) = x^2 + 7x + 10\]
Aspa simple: $Q(x) = (x+5)(x+2)$
\[\therefore\; x^3 + 6x^2 + 3x - 10 = (x+2)(x+5)(x-1)\]
\end{minipage}

\bigskip
\textbf{5. Cambio de variable:} Consiste en identificar términos comunes dentro del
polinomio, luego en conjunto, éstos serán reemplazados por una variable a fin de reducir
el tamaño del polinomio y hacer más sencilla su factorización.
\begin{align*}
P(x) &= x(x+1)(x-2)(x-1) + 1\\
     &= (x^2 - x)(x^2 - x - 2) + 1\\
x^2 - x = a \;&\rightarrow\; a(a-2) + 1 = a^2 - 2a + 1 = (a-1)^2\\
&\therefore\; (x^2 - x - 1)^2
\end{align*}


% ============================================================
% Páginas 10-12
% ============================================================
\textbf{6. Productos notables:} Consiste en utilizar las identidades de los productos
notables como lo es, \hl{la diferencia de cuadrados, la suma y diferencia de cubos, etc.}
para factorizar polinomios.
\[\therefore\; 8x^3 + 27 = (2x)^3 + 3^3 = (2x+3)(4x^2 - 6x + 9)\]

\textbf{7. Quita y pon:} Consiste en identificar el término faltante en el polinomio, para
luego, agregarlo y a su vez quitarlo con el propósito de reagrupar los términos del
polinomio y finalmente factorizarlo usando los productos notables.

\textbf{Nota:} los principales productos notables que se usan en este método son:
\hl{la suma y diferencia de cubos y la identidad de Argand.}

\begin{minipage}[t]{0.436\textwidth}
\[P(x) = x^5 + x + 1\]
Agregando el término que falta: $x^2$
\[x^5 + x + 1 + x^2 - x^2\]
\end{minipage}\hfill
\begin{minipage}[t]{0.533\textwidth}
\begin{gather*}
(x^5 - x^2) + (x^2 + x + 1)\\
x^2(x^3 - 1) + (x^2 + x + 1)\\
x^2(x-1)(x^2+x+1) + (x^2+x+1)\\
(x^2+x+1)(x^3 - x^2 + 1)\\
\therefore\; x^5 + x + 1 = (x^2+x+1)(x^3 - x^2 + 1)
\end{gather*}
\end{minipage}

\subsection*{Observaciones adicionales}

$\ast$ \textbf{Teorema:} Todo polinomio de la forma:
$P(x) = Ax^2 + Bx + C;\ \{A, B, C\} \subset \mathbb{Z} \wedge \forall A \neq 0$
es factorizable en los racionales si: $B^2 - 4AC$ es un cuadrado perfecto.

\textbf{Corolario:} Todo polinomio cuadrático en una variable, si es factorizable, debe
admitir el criterio del aspa simple. Si no admite aspa simple, es porque no es
factorizable en $\mathbb{Q}$.

$\ast$ \textbf{Teorema:} Un polinomio tiene factores de primer grado de coeficientes
racionales si y solo si tiene raíces racionales.

\textbf{Nota:} Toda raíz racional de un polinomio pertenece, necesariamente, al conjunto
de los posibles ceros racionales.

\begin{minipage}[t]{0.339\textwidth}
\[P(x) = x^3 - 3x + 1\]
Sus posibles ceros racionales son: $1$ ó $-1$.
\end{minipage}\hfill
\begin{minipage}[t]{0.194\textwidth}
\[P(1) = -1\]
\[P(-1) = 3\]
\end{minipage}\hfill
\begin{minipage}[t]{0.436\textwidth}
Esto nos indica que no tiene ceros racionales, por lo tanto no tendrá factores lineales.
Indicando que $f(x)$ no será factorizable en los racionales.
\end{minipage}

\bigskip
$\ast$ \textbf{Teorema del factor:} Dado un polinomio $P(x)$, si el número ``$b$'' es un
cero de este polinomio, entonces $(x-b)$ será un factor de $P(x)$.

\begin{minipage}[t]{0.436\textwidth}
\[P(x) = x^3 + 5x + 6\]
\[\text{P.C.R.} = \{1, 2, 3, 6\}\]
Como $P(-1) = (-1)^3 + 5(-1) + 6 = 0$
\end{minipage}\hfill
\begin{minipage}[t]{0.533\textwidth}
$\rightarrow (x - (-1)) \equiv (x+1)$ será un factor de $P(x)$; en tal caso será posible
escribir:
\[\boxed{\therefore\; P(x) = (x+1)\,q(x)}\]
\end{minipage}

\subsection*{Clasificación de factores}

\begin{itemize}[label=\checkmark]
  \item \textbf{Factores primos:} Son todos aquellos factores no constantes \hl{(bases)} que
  presenta un polinomio. Se calcula contando los factores que se encuentran en la base de
  cada exponente.
  \item \textbf{Factores totales:} Es la cantidad total de factores que posee un polinomio.
  Se calcula sumándole 1 al exponente de cada factor primo y luego se le multiplican
  estos valores.
  \item \textbf{Factores algebraicos:} Son todos los factores que posee un polinomio a
  excepción de la unidad. Se calcula restándole 1 a la cantidad de factores totales.
  \item \textbf{Factores compuestos:} Son todos aquellos factores que se han formado de
  multiplicar los factores primos entre sí, teniendo en cuenta sus respectivos
  exponentes. Se calcula utilizando la siguiente fórmula:
  \[\boxed{\#F.T = \#F.P + \#F.C + 1}\]
\end{itemize}

\textcolor{naranja}{\textbf{Ejemplo:}}

\begin{minipage}[t]{0.533\textwidth}
\begin{gather*}
P(x) = x^4(x-2)^2(x-5)\\
\#F.P = 3\\
\#F.T = (4+1)(2+1)(1+1) = 30
\end{gather*}
\end{minipage}\hfill
\begin{minipage}[t]{0.436\textwidth}
\begin{gather*}
\#F.A = 30 - 1 = 29\\
\#F.C = 26
\end{gather*}
\end{minipage}

\newpage
% ------------------------------------------------------------
\section{División algebraica}

Es aquella operación matemática que se realiza entre dos expresiones algebraicas no
nulas llamadas ``dividendo'' y ``divisor'' (las que pueden ser monomios o polinomios) y
consiste en hallar otras dos expresiones algebraicas llamadas ``cociente'' y
``residuo o resto''.

\textbf{1. Notación:} \quad \fcolorbox{verde}{verde!15}{$D(x) = d(x)\cdot q(x) + R(x)$}
\begin{center}
\textcolor{rojo}{$D(x)$}: Dividendo; \quad \textcolor{rojo}{$d(x)$}: Divisor\\
\textcolor{rojo}{$q(x)$}: Cociente; \quad \textcolor{rojo}{$R(x)$}: Resto
\end{center}

\textbf{2. Propiedades:}
\begin{itemize}[label=\checkmark]
  \item El grado del dividendo es mayor o igual que el grado del divisor.
  \[\textcolor{magenta}{G(D(x)) \geq G(d(x))}\]
  \item El grado del cociente es igual al grado del dividendo menos el grado del divisor.
  \[\textcolor{magenta}{G(q(x)) = G(D(x)) - G(d(x))}\]
  \item El grado del resto es menor o igual que el grado del divisor disminuido en 1.
  \[\textcolor{magenta}{G(R(x)) \leq G(d(x)) - 1 \;\longrightarrow\; G(R(x))_{\text{máx}} = G(d(x)) - 1}\]
  \item Si la división es exacta, el resto es un polinomio idénticamente nulo.
  \[\textcolor{magenta}{D(x) = d(x)\cdot q(x) \;\longrightarrow\; R(x) = 0}\]
\end{itemize}

\textbf{3. Principales métodos de división de polinomios:} \hl{Para cualquier método de
división, el dividendo y el divisor deben estar completos y ordenados} respecto a una de
sus variables. \hl{En caso de que faltara uno o más términos, se completará con ceros.}

\textcolor{verde}{\textbf{Método clásico:}} Usado para dividir polinomios de cualquier grado.

\textcolor{verde}{\textbf{Pasos a seguir:}}
\begin{enumerate}
  \item Los términos del dividendo y del divisor deben estar ordenados en forma decreciente
  según el exponente de su variable.
  \item Los términos del cociente se empiezan a obtener de dividir el 1\textsuperscript{er}
  término del dividendo entre el 1\textsuperscript{er} término del divisor.

% ============================================================
% Páginas 13-15
% ============================================================
  \item Cada término del cociente se multiplica por cada término del divisor, se obtienen
  los resultados parciales y a los correspondientes términos del dividendo les vamos a
  restar estos productos parciales.
  \item La expresión obtenida al realizar la última suma, es el residuo de la división.
\end{enumerate}

\textcolor{naranja}{\textbf{Ejemplo:}}
\[
\textcolor{magenta}{D(x) = 5x^6 - 12x^5 + 10x^3 - 7x + 3} \qquad
\textcolor{magenta}{d(x) = x - 2}
\]

{\footnotesize
\[
\begin{array}{rrrrrrr|l}
5x^6 & -12x^5 & \textcolor{magenta}{+0x^4} & +10x^3 & \textcolor{magenta}{+0x^2} & -7x & +3 & \;x-2\\ \cline{8-8}
5x^6 & -10x^5 &        &        &       &     &    & \;5x^5-2x^4-4x^3+2x^2+4x+1\\ \cline{1-2}
     & -2x^5  & +0x^4  &        &       &     &    & \\
     & -2x^5  & +4x^4  &        &       &     &    & \\ \cline{2-3}
     &        & -4x^4  & +10x^3 &       &     &    & \\
     &        & -4x^4  & +8x^3  &       &     &    & \\ \cline{3-4}
     &        &        & 2x^3   & +0x^2 &     &    & \\
     &        &        & 2x^3   & -4x^2 &     &    & \\ \cline{4-5}
     &        &        &        & 4x^2  & -7x &    & \\
     &        &        &        & 4x^2  & -8x &    & \\ \cline{5-6}
     &        &        &        &       & x   & +3 & \\
     &        &        &        &       & x   & -2 & \\ \cline{6-7}
     &        &        &        &       &     & 5  &
\end{array}
\]
}
\[\therefore\; q(x) = 5x^5 - 2x^4 - 4x^3 + 2x^2 + 4x + 1\]
\[\therefore\; R(x) = 5\]

\bigskip
\textcolor{verde}{\textbf{Método de Ruffini:}} Usado para dividir polinomios \hl{cuando el
divisor es lineal.}

\textcolor{verde}{\textbf{Pasos a seguir:}}
\begin{enumerate}
  \item Los coeficientes del dividendo deben estar completos y ordenados en forma
  decreciente según el exponente de su variable.
  \item El valor de la variable se obtiene al igualar el divisor a cero.
  \item Los coeficientes del cociente se obtienen de sumar cada columna, luego que el
  coeficiente anterior se ha multiplicado por el valor de la variable y este resultado
  colocado en la columna siguiente.
  \item El resto se obtiene de sumar la última columna.
\end{enumerate}

\begin{center}
\begin{tikzpicture}[scale=0.8]
  \draw[thick] (0,0) -- (0,2.2);
  \draw[thick] (-1.2,0.6) -- (4.6,0.6);
  \draw[thick] (3.6,0) -- (3.6,2.2);
  \draw (0.2,1.5) rectangle (4.4,2.0); \node at (2.3,1.75) {\cn{1}};
  \node at (-0.6,0.9) {\cn{2}};
  \node at (1.8,0.25) {\cn{3}};
  \node at (4.1,0.25) {\cn{4}};
\end{tikzpicture}
\end{center}

\textcolor{naranja}{\textbf{Ejemplo:}}
\[
\textcolor{magenta}{D(x) = 3x^4 + 4x^3 - 8} \qquad \textcolor{magenta}{d(x) = x + 1}
\]
\begin{minipage}{0.485\textwidth}
\[
\begin{array}{c|cccc|c}
   & 3 & 4  & 0  & 0 & -8\\
-1 &   & -3 & -1 & 1 & -1\\ \hline
   & 3 & 1  & -1 & 1 & -9
\end{array}
\]
\end{minipage}\hfill
\begin{minipage}{0.485\textwidth}
\[\therefore\; q(x) = 3x^3 + x^2 - x + 1\]
\[\therefore\; R(x) = -9\]
\end{minipage}

\medskip
\textbf{Nota:} Si el divisor es de la forma ``$ax + b$'', con $a \neq 0$, los coeficientes
del cociente deben dividirse entre ``$a$'' para obtener el cociente correcto.

\bigskip
\textcolor{verde}{\textbf{Método de Horner:}} Usado para dividir polinomios cuando el divisor
es de un grado mayor o igual que 2.

\textcolor{verde}{\textbf{Pasos a seguir:}}
\begin{enumerate}
  \item Los coeficientes del dividendo deben estar completos y ordenados de forma
  decreciente según el exponente de su variable.
  \item Los coeficientes del divisor deben estar ordenados en forma decreciente y
  completos, según el exponente de su variable pero con signo contrario, salvo el primer
  coeficiente.
  \item Los coeficientes del cociente se obtienen al dividir la suma de los elementos de
  cada columna entre el 1\textsuperscript{er} coeficiente del divisor. Cada coeficiente del
  cociente se multiplica por los demás coeficientes del divisor y dichos resultados se
  ubican a partir de la columna siguiente en forma horizontal.
  \item El resto se obtiene de sumar las últimas columnas.
\end{enumerate}

\begin{minipage}{0.485\textwidth}
\centering
\begin{tikzpicture}[scale=0.8]
  \draw[thick] (0,0) -- (0,2.6);
  \draw[thick] (-1,0.6) -- (5.2,0.6);
  \draw[thick] (4,0) -- (4,2.6);
  \draw (0.2,2.0) rectangle (5.0,2.4); \node at (2.6,2.2) {\cn{1}};
  \draw (-0.6,0.8) rectangle (-0.2,2.4); \node at (-0.4,1.6) {\cn{2}};
  \draw (0.2,0.1) rectangle (3.8,0.5); \node at (2,0.3) {\cn{3}};
  \draw (4.2,0.1) rectangle (5.0,0.5); \node at (4.6,0.3) {\cn{4}};
  \draw[->] (3,-0.6) node[below left] {\small Línea divisora} -- (3.95,-0.05);
\end{tikzpicture}
\end{minipage}\hfill
\begin{minipage}{0.485\textwidth}
$\ast$ La línea divisora se coloca separando tantos términos del final del dividendo
según lo indique el grado del divisor.
\end{minipage}

\bigskip
\textcolor{naranja}{\textbf{Ejemplo:}}
\[
\textcolor{magenta}{D(x) = 6x^5 - 20x^4 - 13x^3 + 20x^2 - 12x + 7} \qquad
\textcolor{magenta}{d(x) = 3x^2 - x + 1}
\]
\begin{minipage}{0.533\textwidth}
\[
\begin{array}{c|cccc|cc}
   &   & -18 & -21 & 19 & & \\[-2pt]
3  & 6 & -20 & -13 & 20 & -12 & 7\\
1  &   & 2   & -2  &    &     & \\
-1 &   &     & -6  & 6  &     & \\
   &   &     &     & -7 & 7   & \\
   &   &     &     &    & \frac{19}{3} & -\frac{19}{3}\\ \hline
   & 2 & -6  & -7  & \frac{19}{3} & \frac{4}{3} & \frac{2}{3}
\end{array}
\]
\end{minipage}\hfill
\begin{minipage}{0.436\textwidth}
\[\therefore\; q(x) = 2x^3 - 6x^2 - 7x + \frac{19}{3}\]
\[\therefore\; R(x) = \frac{4}{3}x + \frac{2}{3}\]
\end{minipage}

\bigskip
\textbf{4. Teorema del resto:} Es una regla práctica que nos va a permitir encontrar el
residuo de una división cualquiera, sin necesidad de realizar la operación.

\textcolor{verde}{\textbf{Pasos a seguir:}}
\begin{enumerate}
  \item El divisor se iguala a cero con el propósito de obtener diversas ecuaciones.
  \item Estas ecuaciones son reemplazadas en el dividendo y la expresión final viene a
  ser el resto de la división.
\end{enumerate}

\textcolor{naranja}{\textbf{Ejemplo:}}
\[
\textcolor{magenta}{D(x) = (x-2)^5(x+5)^5 + x(x+3) + 15} \qquad
\textcolor{magenta}{d(x) = x^2 + 3x - 8}
\]
\begin{minipage}[t]{0.388\textwidth}
\begin{gather*}
x^2 + 3x - 8 = 0\\
x^2 + 3x = 8\\[6pt]
\therefore\; \textcolor{magenta}{R(x) = -9}
\end{gather*}
\end{minipage}\hfill
\begin{minipage}[t]{0.582\textwidth}
\textcolor{magenta}{Acomodando el dividendo:}
\begin{gather*}
(x^2 + 3x - 10)^5 + (x^2 + 3x) + 15\\
(8 - 10)^5 + (8) + 15\\
-32 + 23 = -9
\end{gather*}
\end{minipage}


% ============================================================
% Páginas 16-18
% ============================================================
\newpage
\section{Cocientes notables}

Son aquellos cocientes exactos que se pueden obtener en forma directa sin necesidad de
efectuar la operación de división.

\textbf{1. Notación:} \quad
\caja{$\dfrac{x^m \pm y^m}{x \pm y}$ \quad siendo $m \in \mathbb{Z}^+$ y $m \geq 2$}

\medskip
\textbf{2. Condición necesaria y suficiente:} Todo cociente notable debe poseer como número
de términos a un número que sea entero positivo, el cual se obtiene de la siguiente manera:
\[\caja{$\dfrac{x^p \pm y^q}{x^a \pm y^b}\;, \quad \#t = \dfrac{p}{a} = \dfrac{q}{b}$}\]

\begin{center}
\hl{$\ast$ El ``$\#t$'' suele representarse con la letra ``$n$''.}
\end{center}

\textbf{3. Casos posibles:}

\begin{minipage}[t]{0.320\textwidth}
\centering
1\textsuperscript{ero}: Para todo $n \geq 2$\\[4pt]
\caja{$\dfrac{x^p - y^q}{x^a - y^b} = + + + \cdots + +$}\\[4pt]
\small (Todos los términos de su desarrollo son positivos)
\end{minipage}\hfill
\begin{minipage}[t]{0.320\textwidth}
\centering
2\textsuperscript{do}: Para todo ``$n$'' impar\\[4pt]
\caja{$\dfrac{x^p + y^q}{x^a + y^b} = + - + \cdots - +$}
\end{minipage}\hfill
\begin{minipage}[t]{0.320\textwidth}
\centering
3\textsuperscript{ero}: Para todo ``$n$'' par\\[4pt]
\caja{$\dfrac{x^p - y^q}{x^a + y^b} = + - + \cdots + -$}
\end{minipage}

\begin{center}
\small (En los casos 2 y 3: los términos de posición impar son positivos y los de
posición par negativos)
\end{center}

\subsection*{Características del cociente}
\begin{enumerate}
  \item El resultado de un cociente notable es un polinomio homogéneo de grado de
  homogeneidad $(n-1)$; es un polinomio de ``$n$'' términos completo y ordenado respecto
  de sus dos variables.
  \item Si el divisor es ``$x + y$'' el polinomio resultante del cociente presenta signos
  alternados: $(+)(-)(+)(-)\cdots$ etc. Si el divisor es ``$x - y$'' el polinomio resultante
  del cociente presenta todos sus signos positivos.
  \item Si ``$n$'' es impar, existe un solo término central en el cual se cumple lo
  siguiente: $n = 2k - 1$.
  \item Si ``$n$'' es par, existen dos términos centrales de los cuales el 1\textsuperscript{ero}
  cumple que: $n = 2k$, mientras que el 2\textsuperscript{do} es sólo el término consecutivo.
\end{enumerate}
\[\#t = G.A.(\text{Término } k) + 1\]

\textbf{4. Término general:} \quad \caja{$T_k = \pm\, x^{n-k}\, y^{k-1}$}

\begin{minipage}[t]{0.436\textwidth}
\begin{itemize}[label={}]
  \item $k = 1; 2; 3; \cdots; n$
  \item $x$ = 1\textsuperscript{er} término del divisor
  \item $y$ = 2\textsuperscript{do} término del divisor
  \item $n$ = número de términos
\end{itemize}
\end{minipage}\hfill
\begin{minipage}[t]{0.533\textwidth}
\begin{itemize}[label=\textcolor{magenta}{$\ast$}]
  \item El ``$\pm$'' depende de si ``$k$'' es par o impar.
  \item El ``$\pm$'' es prescindible de la fórmula si se trata del primer caso.
\end{itemize}
\end{minipage}

\medskip
\textbf{5. Términos equidistantes:} Son aquellos términos que se encuentran a una misma
distancia o posición siendo contados desde los extremos, y presentan las siguientes
características:
\begin{itemize}[label=\textcolor{naranja}{\checkmark}]
  \item La suma de sus posiciones es igual al número de términos más la unidad.
  \[\caja{$k_1 + k_2 = n + 1$}\]
  \item Sus bases cambian de exponente.
  \[\caja{$\begin{aligned}
    T_{k_1} &= \pm\, x^{n-k_1}\cdot y^{k_1-1}\\
    T_{k_2} &= \pm\, x^{k_2-1}\cdot y^{n-k_2}
  \end{aligned}$}\]
\end{itemize}

\newpage
% ------------------------------------------------------------
\section{Factoriales y combinatorias}

\textbf{1. Factoriales:} El factorial de un número, es el producto de todos los números
naturales consecutivos, desde el 1 hasta el número dado.

\begin{minipage}[t]{0.436\textwidth}
\textbf{Notación:}
\[\fbox{$n! = 1\cdot2\cdot3\cdots(n-2)(n-1)n$}\]
\[\fbox{$n! = n(n-1)!$}\]

\textbf{Factoriales básicos:}
\[
\begin{array}{|rcl|rcl|}
\hline
0! &=& 1  & 5! &=& 120\\
1! &=& 1  & 6! &=& 720\\
2! &=& 2  & 7! &=& 5040\\
3! &=& 6  & 8! &=& 40320\\
4! &=& 24 & 9! &=& 362880\\ \hline
\end{array}
\]

\textbf{Propiedades básicas:}
\[
\begin{array}{|l|}
\hline
\textcolor{magenta}{\ast}\ (-n)! = \nexists\\
\textcolor{magenta}{\ast}\ (n/m)! = \nexists\\
\textcolor{magenta}{\ast}\ n! = n(n-1)! \text{ para } n \geq 2\\
\textcolor{magenta}{\ast}\ n! + (n+1)! = (n+2)\,n!\\
\textcolor{magenta}{\ast}\ n! = m! \rightarrow n = m\\ \hline
\end{array}
\]
\end{minipage}\hfill
\begin{minipage}[t]{0.533\textwidth}\small
\textbf{Propiedades adicionales:}
\[
\begin{array}{|l|}
\hline
n! + (n+1)! + (n+2)! = (n+2)^2\cdot n!\\[6pt]
\dfrac{n}{(n+1)!} = \dfrac{1}{n!} - \dfrac{1}{(n+1)!} \quad \text{para } n \geq 2\\[10pt]
1\cdot1! + 2\cdot2! + 3\cdot3! + \cdots + n\cdot n! = (n+1)! - 1\\[4pt]
2^2\cdot1! + 3^2\cdot2! + 4^2\cdot3! + \cdots + (n+1)^2 n! = (n+2)! - 2\\ \hline
\end{array}
\]
\[
\begin{array}{|l|}
\hline
\displaystyle\sum_{i=1}^{n}\frac{i}{(i+1)!} = \frac{1}{2!} + \frac{2}{3!} + \frac{3}{4!} + \cdots + \frac{n}{(n+1)!} = 1 - \frac{1}{(n+1)!}\\[14pt]
\displaystyle\sum_{i=0}^{\infty}\frac{x^i}{i!} = 1 + x + \frac{x^2}{2!} + \frac{x^3}{3!} + \cdots = e^x\\[14pt]
\displaystyle\sum_{i=0}^{\infty}\frac{i\,x^i}{i!} = x + \frac{2x^2}{2!} + \frac{3x^3}{3!} + \cdots = x\cdot e^x\\[10pt] \hline
\end{array}
\]
\end{minipage}

\bigskip
\textbf{Semifactorial:} Se aplica a un número natural par y viene a ser el producto de
todos los números naturales pares consecutivos, desde el 2 hasta el número dado.
\[\fbox{$n!! = 2\cdot4\cdot6\cdot8\cdots(n-2)\,n = 2^{\left(\frac{n}{2}\right)}\cdot\left(\dfrac{n}{2}\right)!$}\]

\textbf{Cofactorial:} Se aplica a un número natural impar y viene a ser el producto de
todos los números naturales impares consecutivos, desde el 1 hasta el número dado.
\[\fbox{$n!! = 1\cdot3\cdot5\cdot7\cdots(n-2)\,n = \dfrac{n!}{2^{\left(\frac{n-1}{2}\right)}\left(\dfrac{n-1}{2}\right)!}$}\]


% ============================================================
% Páginas 19-21
% ============================================================
\subsection*{Fórmulas de seno y coseno}
\[
\begin{array}{|l|}
\hline
\displaystyle\sum_{i=0}^{\infty}\frac{(-1)^i x^{2i+1}}{(2i+1)!} = x - \frac{x^3}{3!} + \frac{x^5}{5!} - \frac{x^7}{7!} + \cdots = \operatorname{sen} x\\[16pt]
\displaystyle\sum_{i=0}^{\infty}\frac{(-1)^i\cdot x^{2i}}{(2i)!} = 1 - \frac{x^2}{2!} + \frac{x^4}{4!} - \frac{x^6}{6!} + \cdots = \cos x\\[12pt] \hline
\end{array}
\]

\textbf{2. Número combinatorio:} Llamado también combinatoria y viene a ser el número de
formas en que se pueden extraer subconjuntos a partir de un conjunto dado.

\textbf{Notación:}
\[
\fbox{$\underbrace{\dbinom{n}{k}}_{\text{Coeficiente binomial}} = \underbrace{C^{n}_{k}}_{\text{Combinatoria}} = \dfrac{n!}{k!\,(n-k)!}\;; \quad n \in \mathbb{N},\ k \in \mathbb{N}_0,\ n \geq k$}
\]
\begin{center}
``$n$'': Base \qquad ``$k$'': Orden
\end{center}

\begin{minipage}[t]{0.466\textwidth}
\textbf{Propiedades básicas:}
\[
\begin{array}{|ll|}
\hline
C^{n}_{0} = 1 & C^{n}_{1} = n\\
C^{n}_{n} = 1 & C^{n}_{k} = C^{n}_{n-k}\\ \hline
\end{array}
\]
\[
\begin{array}{|l|}
\hline
C^{n}_{k} = \dfrac{n}{k}\,C^{n-1}_{k-1}\\[10pt]
C^{n}_{k} = \dfrac{n}{n-k}\,C^{n-1}_{k}\\[10pt]
C^{n}_{k} = \dfrac{n-k+1}{k}\,C^{n}_{k-1}\\[10pt]
C^{n}_{k} = C^{m}_{p} \rightarrow n = m \wedge (k = p \vee k + p = n)\\[6pt]
C^{n}_{k} + C^{n}_{k+1} = C^{n+1}_{k+1}\\[6pt]
C^{n}_{0} + C^{n}_{1} + C^{n}_{2} + \cdots + C^{n}_{n} = 2^n\\[6pt]
C^{n}_{0} + C^{n}_{2} + C^{n}_{4} + \cdots + C^{n}_{n} = 2^{n-1}\\[4pt] \hline
\end{array}
\]
\end{minipage}\hfill
\begin{minipage}[t]{0.504\textwidth}
\[
\fbox{$C^{n}_{1} + C^{n}_{3} + C^{n}_{5} + \cdots + C^{n}_{n-1} = 2^{n-1}$}
\]
\textbf{Propiedades auxiliares:}
\[
\begin{array}{|l|}
\hline
C^{n}_{0} + C^{n+1}_{1} + C^{n+2}_{2} + \cdots + C^{n+k}_{k} = C^{n+k+1}_{k}\\[6pt]
C^{n}_{1} + 2C^{n}_{2} + 3C^{n}_{3} + \cdots + nC^{n}_{n} = n\cdot 2^{n-1}\\[6pt]
C^{n}_{1} + 2^2C^{n}_{2} + 3^2C^{n}_{3} + \cdots + n^2C^{n}_{n} = \dfrac{n(n+1)}{2}\cdot 2^{n-1}\\[10pt]
(C^{n}_{0})^2 + (C^{n}_{1})^2 + (C^{n}_{2})^2 + \cdots + (C^{n}_{n})^2 = \dfrac{(2n)!}{(n!)^2}\\[8pt] \hline
\end{array}
\]
\end{minipage}

\newpage
% ------------------------------------------------------------
\section{Binomio de Newton}

La potencia de un binomio es un polinomio que se denomina binomial o de Newton.
Así tenemos:
\[(x+y)^n = \binom{n}{0}x^n + \binom{n}{1}x^{n-1}y + \cdots + \binom{n}{n-1}x\,y^{n-1} + \binom{n}{n}y^n\]

\subsection*{Propiedades del binomio de Newton}
\begin{enumerate}[label=\cn{\arabic*}]
  \item \hl{En el desarrollo existen $(n+1)$ términos siendo $n \in \mathbb{N}$.}
  \item Los coeficientes de los términos equidistantes de los extremos son iguales, lo
  cual es evidente por ser números combinatorios complementarios.
  \item Si los términos del binomio tienen signos contrarios, los términos del desarrollo
  serán alternativamente positivos y negativos.
  \begin{center}
  Siendo los de lugar impar \textcolor{rojo}{$(+)$}\\
  Siendo los de lugar par \textcolor{rojo}{$(-)$}
  \end{center}
  \item Si los dos términos del binomio son negativos, todos los términos del desarrollo
  serán positivos o negativos según que el exponente sea par o impar. En efecto, se tiene:
  \[(-x-a)^m = [-(x+a)]^m\]
  \item La suma de los coeficientes de los términos de lugar impar es igual a la suma de
  los de lugar par.
  \item Término central de un binomio con exponente par: $\dfrac{n+2}{2}$
  \item Términos centrales de un binomio con exponente impar: $\dfrac{n+1}{2}$ \ y \
  $\dfrac{n+1}{2} + 1$
  \item La suma de coeficientes de un binomio de Newton cuya fórmula es
  \textcolor{magenta}{$(ax^{\alpha} \pm by^{\beta})^n$} está dada por:
  \[\sum \text{coef.} = (a \pm b)^n \quad \text{siendo } x = y = 1\]
  \item La suma de exponentes (grados absolutos) de los términos del desarrollo de un
  binomio de Newton cuya forma más general es:
  \textcolor{magenta}{$(ax^{\alpha} \pm by^{\beta})^n$}
  \[\sum G.A. = \left[\frac{n(n+1)}{2}\right](\alpha + \beta)\]
  \item \hl{La fórmula del término general es: $t_{k+1} = C^{n}_{k}\cdot x^{n-k}\cdot y^{k}$}
  \item La suma de las posiciones de dos términos equidistantes es: $n + 2$
\end{enumerate}

\subsection*{Expansión de Leibnitz}
Llamado también desarrollo de Leibnitz. Es una generalización del binomio de Newton que
tiene por finalidad hacer más fácil el desarrollo de la potencia $n$-ésima de los
binomios que presentan más de 2 bases. Su forma es:
\[(\underbrace{a + b + c + d + \cdots}_{\text{``$k$'' bases}})^n\]

\begin{itemize}[label=\textcolor{naranja}{\checkmark}]
  \item El número de términos está dado por:
  \[\#t = C^{n+k-1}_{k-1}\]
  \item Su término general está dado por:
  \[
  T_{\text{General}} = \frac{n!}{\underbrace{\alpha!\,\beta!\,\theta!\,\lambda!\cdots}_{\text{Exponentes de las bases}}}\;
  \underbrace{a^{\alpha}\cdot b^{\beta}\cdot c^{\theta}\cdot d^{\lambda}\cdots}_{\text{Bases}}
  \qquad \text{obs: } \fbox{$\underbrace{\alpha + \beta + \theta + \lambda + \cdots}_{\text{``$k$'' sumandos}} = n$}
  \]
\end{itemize}

\textbf{Nota:} Cuando todas las bases de la expansión de Leibnitz presenten exponentes
iguales a la unidad, la suma de grados absolutos de todos sus términos se calculará
usando la siguiente fórmula deducida:
\[\sum G.A. = n\cdot C^{n+k-1}_{k-1} = n\cdot\#t\]

\subsection*{\textcolor{rojo}{Expansión o desarrollo del binomio de Newton, cuando el exponente es negativo y/o fraccionario}}
\[
B_{(x)} = (1+x)^n = \binom{n}{0} + \binom{n}{1}x + \binom{n}{2}x^2 + \binom{n}{3}x^3 + \binom{n}{4}x^4 + \cdots\infty
= \sum_{k=0}^{\infty}\binom{n}{k}x^k
\]
\begin{center}
\textcolor{azul}{Serie binomial o Serie binómica de Newton}
\end{center}

\textbf{Propiedades:}
\begin{itemize}[label=\textcolor{magenta}{\checkmark}]
  \item Cuando el exponente ``$n$'' no es un entero positivo, el número de términos es ilimitado.
  \item Su término general es: \ \fbox{$t_{k+1} = \dbinom{n}{k}\cdot x^k$}
\end{itemize}

\textbf{Obs:} \ \fbox{$\displaystyle\lim_{x\to 0}(1+x)^n = 1 + nx$}

\medskip
\textbf{Ejemplo:}
\[
(1-x)^{-1/3} = \binom{-\frac{1}{3}}{0} + \binom{-\frac{1}{3}}{1}(-x) + \binom{-\frac{1}{3}}{2}(-x)^2 + \cdots = 1 + \frac{x}{3} + \frac{2x^2}{9} + \cdots
\]


% ============================================================
% Páginas 22-24
% ============================================================
\newpage
\section{Racionalización de radicales}

Es el proceso que consiste en eliminar la raíz o raíces que están presentes en el
denominador de una fracción.

\textbf{1. Factor racionalizante (F.R.):} El factor racionalizante de una expresión
irracional es otra expresión irracional que, multiplicada por la primera, la convierte en
una expresión racional.
\[\sqrt[4]{5} \rightarrow \text{F.R.} = \sqrt[4]{5^3} \rightarrow \left(\sqrt[4]{5}\right)\left(\sqrt[4]{5^3}\right) = \sqrt[4]{5^4} = 5\]

\textbf{2. Tipos de radicales:}

\textbf{2.1. Homogéneos:} Son todos aquellos radicales que presentan el mismo índice.
\[\sqrt[4]{x}\,,\quad -2\sqrt[4]{y}\,,\quad \frac{1}{5}\sqrt[4]{xyz}\]

\textbf{2.2. Semejantes:} Son todos aquellos radicales que presentan el mismo índice y a su
vez el mismo radicando.
\[2\sqrt[5]{x}\,,\quad 5\sqrt[5]{x}\,,\quad \frac{-\sqrt{3}}{2}\sqrt[5]{x}\]

\textbf{3. Propiedades:} Las propiedades se presentan según la forma que tenga el denominador.

\begin{minipage}[t]{0.242\textwidth}
\centering
\textcolor{magenta}{1\textsuperscript{ERA} FORMA: $\sqrt{a}$}\\[6pt]
\fbox{$\dfrac{b}{\sqrt{a}}\left(\dfrac{\sqrt{a}}{\sqrt{a}}\right)$}
\end{minipage}\hfill
\begin{minipage}[t]{0.320\textwidth}
\centering
\textcolor{magenta}{2\textsuperscript{DA} FORMA: $\sqrt{a} \pm \sqrt{b}$}\\[6pt]
\fbox{$\dfrac{c}{\sqrt{a} \pm \sqrt{b}}\left(\dfrac{\sqrt{a} \mp \sqrt{b}}{\sqrt{a} \mp \sqrt{b}}\right)$}
\end{minipage}\hfill
\begin{minipage}[t]{0.407\textwidth}
\centering
\textcolor{magenta}{3\textsuperscript{ERA} FORMA: $\sqrt{a} + \sqrt{b} \pm \sqrt{c}$}\\[6pt]
\fbox{$\dfrac{d}{\sqrt{a} + \sqrt{b} \pm \sqrt{c}}\left(\dfrac{\sqrt{a} + \sqrt{b} \mp \sqrt{c}}{\sqrt{a} + \sqrt{b} \mp \sqrt{c}}\right)$}
\end{minipage}

\bigskip
\begin{minipage}[t]{0.388\textwidth}
\centering
\textcolor{magenta}{4\textsuperscript{TA} FORMA: $\sqrt[n]{a^m}$}\\[6pt]
\fbox{$\dfrac{b}{\sqrt[n]{a^m}}\left(\dfrac{\sqrt[n]{a^{n-m}}}{\sqrt[n]{a^{n-m}}}\right)$}
\end{minipage}\hfill
\begin{minipage}[t]{0.582\textwidth}
\centering
\textcolor{magenta}{5\textsuperscript{TA} FORMA: $\sqrt[n]{a} \pm \sqrt[n]{b}$} $\rightarrow n$ es par\\[6pt]
\fbox{$\dfrac{c}{\sqrt[n]{a} \pm \sqrt[n]{b}}\left(\dfrac{\sqrt[n]{a} \mp \sqrt[n]{b}}{\sqrt[n]{a} \mp \sqrt[n]{b}}\right)$}
\end{minipage}

\bigskip
\textcolor{magenta}{6\textsuperscript{TA} FORMA:} $n \in \mathbb{Z}^+ \wedge n \geq 2$
\begin{gather*}
\left(\sqrt[n]{A} - \sqrt[n]{B}\right)\text{F.R.} = A - B\\
\text{F.R.} = \sqrt[n]{A}^{\,n-1} + \sqrt[n]{A}^{\,n-2}\sqrt[n]{B} + \cdots + \sqrt[n]{B}^{\,n-1}
\end{gather*}

\textcolor{magenta}{7\textsuperscript{MA} FORMA:} $n \in \mathbb{Z}^+ \wedge n$ es impar
\begin{gather*}
\left(\sqrt[n]{A} + \sqrt[n]{B}\right)\text{F.R.} = A + B\\
\text{F.R.} = \sqrt[n]{A}^{\,n-1} - \sqrt[n]{A}^{\,n-2}\sqrt[n]{B} + \cdots + \sqrt[n]{B}^{\,n-1}
\end{gather*}

\textcolor{magenta}{8\textsuperscript{VA} FORMA:} $n \in \mathbb{Z}^+ \wedge n$ es par
\begin{gather*}
\left(\sqrt[n]{A} + \sqrt[n]{B}\right)\text{F.R.} = A - B\\
\text{F.R.} = \sqrt[n]{A}^{\,n-1} - \sqrt[n]{A}^{\,n-2}\sqrt[n]{B} + \cdots - \sqrt[n]{B}^{\,n-1}
\end{gather*}

$\ast$ En la racionalización se usa con mucha frecuencia algunas identidades de los
productos notables como lo son: la diferencia de cuadrados, la suma y diferencia de cubos,
la identidad de Lagrange, etc. Así también, se hace necesario usar la factorización y la
teoría de cocientes notables.

\textbf{4. Radicales dobles:}
\begin{itemize}[label=$\ast$]
  \item \caja{$\sqrt{A \pm \sqrt{B}} = \sqrt{\dfrac{A+C}{2}} \pm \sqrt{\dfrac{A-C}{2}}\,, \quad C = \sqrt{A^2 - B}$}
  \quad \textcolor{magenta}{$(C \in \mathbb{Q}^+)$}
  \item \caja{$\sqrt{A \pm 2\sqrt{B}} = \sqrt{m} \pm \sqrt{n}\,, \quad B = m\cdot n \text{ y } A = m + n$}
  \quad \textcolor{magenta}{$(m > n)$}
\end{itemize}

\textbf{5. Radicales triples:}
\begin{itemize}[label=$\ast$]
  \item \caja{$\sqrt{(a+b+c) + 2\sqrt{ab} + 2\sqrt{bc} + 2\sqrt{ac}} = \sqrt{a} + \sqrt{b} + \sqrt{c}$}
  \item \caja{$\sqrt{(a+b+c) + 2\sqrt{ab} - 2\sqrt{bc} - 2\sqrt{ac}} = \sqrt{a} + \sqrt{b} - \sqrt{c}$}
  \item \caja{$\sqrt{(a+b+c) - 2\sqrt{ab} - 2\sqrt{bc} + 2\sqrt{ac}} = \sqrt{a} - \sqrt{b} + \sqrt{c}$}
  \item \caja{$\sqrt{(a+b+c) - 2\sqrt{ab} + 2\sqrt{bc} - 2\sqrt{ac}} = -\sqrt{a} + \sqrt{b} + \sqrt{c}$}
\end{itemize}

\textbf{6. Raíz cúbica:}
\begin{itemize}[label=$\ast$]
  \item \fbox{$\sqrt[3]{A \pm \sqrt{B}} = n \pm \sqrt{m}$}
  \begin{itemize}[label=\checkmark]
    \item $C = \sqrt[3]{A^2 - B}$
    \item $A = 4n^3 - 3nC$
    \item $m = n^2 - C$
  \end{itemize}
\end{itemize}

\textbf{7. Extraer raíz cuadrada:} El polinomio tiene que ser completo y ordenado.

Extraer la raíz cuadrada de: \textcolor{magenta}{$25x^6 - 20x^5 + 34x^4 - 52x^3 + 24x^2 + 15$}

{\small
\[
\begin{array}{rrrrrrr|l}
\sqrt{25x^6} & -20x^5 & +34x^4 & -52x^3 & +24x^2 & \textcolor{naranja}{+0x} & +15 & \;5x^3 - 2x^2 + 3x - 4\\ \cline{8-8}
25x^6 &        &        &        &        &     &     & \;(10x^3 + (-2x^2))(-2x^2)\\ \cline{1-1}
0     & -20x^5 & +34x^4 &        &        &     &     & \;(10x^3 - 4x^2 + 3x)(3x)\\
      & -20x^5 & +4x^4  &        &        &     &     & \;(10x^3 - 4x^2 + 6x - 4)(-4)\\ \cline{2-3}
      &        & 30x^4  & -52x^3 & +24x^2 &     &     & \\
      &        & 30x^4  & -12x^3 & +9x^2  &     &     & \\ \cline{3-5}
      &        &        & -40x^3 & +15x^2 & +0x & +15 & \\
      &        &        & -40x^3 & +16x^2 & -24x & +16 & \\ \cline{4-7}
\multicolumn{3}{r}{\text{Residuo} \rightarrow} & & -x^2 & +24x & -1 &
\end{array}
\]
}

\[\therefore\; 25x^6 - 20x^5 + 34x^4 - 52x^3 + 24x^2 + 15 = (5x^3 - 2x^2 + 3x - 4)^2 + (-x^2 + 24x - 1)\]

\begin{tabular}{@{}l@{\;}l}
\textcolor{magenta}{Polinomio radicando de grado par} $\rightarrow$ & $P(x) = 25x^6 - 20x^5 + 34x^4 - 52x^3 + 24x^2 + 15$\\
\textcolor{magenta}{Polinomio raíz} $\rightarrow$ & $q(x) = 5x^3 - 2x^2 + 3x - 4$\\
\textcolor{magenta}{Polinomio residuo} $\rightarrow$ & $R(x) = -x^2 + 24x - 1$
\end{tabular}

\medskip
$\ast$ \textcolor{magenta}{Cuando el polinomio residuo es igual a cero es raíz cuadrada
exacta, caso contrario será raíz cuadrada inexacta.}


% ============================================================
% Páginas 25-27
% ============================================================
\newpage
\section{Matrices y determinantes}

\textbf{Definición:} Una matriz de orden $m\times n$ es un arreglo rectangular de elementos
(números o funciones), dispuestos en $m$ líneas horizontales, llamadas filas y en $n$
líneas verticales llamadas columnas, de la forma:

\begin{minipage}{0.436\textwidth}
\[
A = \begin{bmatrix}
a_{11} & a_{12} & \cdots & a_{1n}\\
a_{21} & a_{22} & \cdots & a_{2n}\\
a_{31} & a_{32} & \cdots & a_{3n}\\
\vdots & \vdots & \ddots & \vdots\\
a_{m1} & a_{m2} & \cdots & a_{mn}
\end{bmatrix}
\]
\end{minipage}\hfill
\begin{minipage}{0.533\textwidth}
$\ast$ Los subíndices indican la posición del elemento dentro de la matriz, el primero
denota la fila $i$ y el segundo la columna $j$. Por ejemplo, el elemento $a_{25}$ se
ubica en la segunda fila y quinta columna de la matriz.

$\ast$ La dimensión de una matriz es el número $m\times n$ de elementos que tiene la matriz.
\end{minipage}

\medskip
\textbf{Matrices iguales:} Dos matrices son iguales cuando tienen la misma dimensión y
cuando los elementos que ocupan los mismos lugares son iguales. Si
$A = (a_{ij})_{m\times n}$ y $B = (b_{ij})_{m\times n}$, entonces $A = B$ si y sólo si
$a_{ij} = b_{ij}$ para cada valor de $i, j$. Las siguientes matrices no son iguales:
\[
A = \begin{pmatrix} 1 & 2 & -3\\ 0 & 5 & 6 \end{pmatrix}
\qquad
B = \begin{pmatrix} 1 & 0\\ 2 & 5\\ -3 & 6 \end{pmatrix}
\]
\begin{center}
\textcolor{naranja}{Orden $2\times3$} \hspace{3cm} \textcolor{naranja}{Orden $3\times2$}
\end{center}

\subsection*{Algunos tipos de matrices}

\textbf{1. Matriz cuadrada:} Tiene el mismo número de filas que de columnas, es decir
\hl{$m = n$}, y se dice que la matriz cuadrada es de \hl{orden $n$}. La diagonal principal
de una matriz cuadrada es el conjunto formado por los elementos
$a_{11}, a_{22}, a_{33}, a_{44}, \ldots, a_{nn}$ y la traza de la matriz cuadrada es el
número dado por la suma de los elementos de la diagonal principal, es decir:
\[\fbox{$\operatorname{Traz}(A) = a_{11} + a_{22} + a_{33} + a_{44} + \cdots + a_{nn} = \displaystyle\sum_{i=1}^{n} a_{ii}$}\]

\textbf{2. Matriz rectangular:} Es toda matriz en la que $m \neq n$.

\begin{minipage}[t]{0.485\textwidth}
\textbf{3. Matriz fila:} Es una matriz de orden $1\times n$.
\[A = \begin{pmatrix} a_{11} & a_{12} & a_{13} & \cdots & a_{1n}\end{pmatrix}\]
\end{minipage}\hfill
\begin{minipage}[t]{0.485\textwidth}
\textbf{4. Matriz columna:} Es una matriz de orden $m\times 1$.
\[A = \begin{pmatrix} a_{11}\\ a_{21}\\ a_{31}\\ \vdots\\ a_{m1}\end{pmatrix}\]
\end{minipage}

\begin{minipage}[t]{0.485\textwidth}
\textbf{5. Matriz nula:} Es aquella que tiene todos sus elementos nulos.
\[\begin{pmatrix} 0 & 0 & 0\\ 0 & 0 & 0\end{pmatrix}\]

\textbf{7. Matriz escalar:} Es una matriz diagonal en la que todos los elementos de la
diagonal principal son iguales a una constante.
\[
B = \begin{bmatrix}
k & 0 & 0 & \cdots & 0\\
0 & k & 0 & \cdots & 0\\
0 & 0 & k & \cdots & 0\\
\vdots & \vdots & \vdots & \ddots & \vdots\\
0 & 0 & 0 & \cdots & k
\end{bmatrix}
\]
\end{minipage}\hfill
\begin{minipage}[t]{0.485\textwidth}
\textbf{6. Matriz diagonal:} Es una matriz cuadrada en la que todos los elementos que no
pertenecen a la diagonal principal son nulos.
\[
B = \begin{bmatrix}
a_{11} & 0 & 0 & \cdots & 0\\
0 & a_{22} & 0 & \cdots & 0\\
0 & 0 & a_{33} & \cdots & 0\\
\vdots & \vdots & \vdots & \ddots & \vdots\\
0 & 0 & 0 & \cdots & a_{nn}
\end{bmatrix}
\]

\textbf{8. Matriz unidad o identidad:} Es la matriz escalar en la que $k = 1$.
\[
B = \begin{bmatrix}
1 & 0 & 0 & \cdots & 0\\
0 & 1 & 0 & \cdots & 0\\
0 & 0 & 1 & \cdots & 0\\
\vdots & \vdots & \vdots & \ddots & \vdots\\
0 & 0 & 0 & \cdots & 1
\end{bmatrix}_{n\times n} = I_n
\]
\end{minipage}

\medskip
\begin{minipage}[t]{0.485\textwidth}
\textbf{9. Matriz triangular superior:} Es la matriz cuadrada que tiene todos sus elementos
que se encuentran por debajo de la diagonal principal nulos, es decir $a_{ij} = 0$ para
todo $i > j$.
\[
A = \begin{bmatrix}
a_{11} & a_{12} & a_{13} & a_{14}\\
0 & a_{22} & a_{23} & a_{24}\\
0 & 0 & a_{33} & a_{34}\\
0 & 0 & 0 & a_{44}
\end{bmatrix}
\]
\end{minipage}\hfill
\begin{minipage}[t]{0.485\textwidth}
\textbf{10. Matriz triangular inferior:} Es la matriz cuadrada que tiene todos sus
elementos que se encuentran por encima de la diagonal principal nulos, es decir
$a_{ij} = 0$ para todo $i < j$.
\[
A = \begin{bmatrix}
a_{11} & 0 & 0 & 0\\
a_{21} & a_{22} & 0 & 0\\
a_{31} & a_{32} & a_{33} & 0\\
a_{41} & a_{42} & a_{43} & a_{44}
\end{bmatrix}
\]
\end{minipage}

\medskip
\textbf{11. Matriz traspuesta:} Es la matriz que se obtiene de la matriz
$A = (a_{ij})_{m\times n}$ intercambiando las filas por columnas. Se denota
$A^t = (a_{ji})_{n\times m}$.
\[
A = \begin{bmatrix}
a_{11} & a_{12} & a_{13}\\
a_{21} & a_{22} & a_{23}\\
a_{31} & a_{32} & a_{33}\\
a_{41} & a_{42} & a_{43}
\end{bmatrix}
\;\longrightarrow\;
A^t = \begin{bmatrix}
a_{11} & a_{21} & a_{31} & a_{41}\\
a_{12} & a_{22} & a_{32} & a_{42}\\
a_{13} & a_{23} & a_{33} & a_{43}
\end{bmatrix}
\]

\begin{minipage}[t]{0.485\textwidth}
\textbf{12. Matriz simétrica:} Es toda matriz tal que $A = A^t$.
\[
A = \begin{bmatrix} 1 & 0 & 1\\ 0 & 2 & 4\\ 1 & 4 & 3\end{bmatrix}
\rightarrow
A^t = \begin{bmatrix} 1 & 0 & 1\\ 0 & 2 & 4\\ 1 & 4 & 3\end{bmatrix}
\]
\end{minipage}\hfill
\begin{minipage}[t]{0.485\textwidth}
\textbf{13. Matriz antisimétrica:} Es toda matriz tal que $A = -A^t$.
\[
A = \begin{pmatrix} 0 & -1\\ 1 & 0\end{pmatrix}
\rightarrow
A^t = \begin{pmatrix} 0 & 1\\ -1 & 0\end{pmatrix}
\]
\end{minipage}

\subsection*{Operaciones con matrices}

\begin{minipage}[t]{0.485\textwidth}
\textbf{Adición de matrices:} Si $A = (a_{ij})_{m\times n}$ y $B = (b_{ij})_{m\times n}$, se
define entonces la suma $A + B$ como la matriz de orden $m\times n$,
$C = (c_{ij})_{m\times n}$ tal que: \hl{$c_{ij} = a_{ij} + b_{ij}$}.
\[
A = \begin{pmatrix} 3 & -2 & 1\\ -4 & 7 & 5\end{pmatrix}
\quad
B = \begin{pmatrix} -9 & 8 & 6\\ -7 & 1 & 2\end{pmatrix}
\]
\[
A + B = C = \begin{pmatrix} -6 & 6 & 7\\ -11 & 8 & 7\end{pmatrix}
\]
\end{minipage}\hfill
\begin{minipage}[t]{0.485\textwidth}
\textbf{Diferencia de matrices:} La diferencia de matrices $A$ y $B$, de orden $m\times n$,
se define como la matriz $D = A + (-B)$, es decir, $D = (d_{ij})_{m\times n}$ tal que
\hl{$d_{ij} = a_{ij} - b_{ij}$}.

\textbf{Propiedades:} Si $A$, $B$, $C$ son matrices de orden $m\times n$, se cumple:
\begin{itemize}[label=\textcolor{naranja}{\checkmark}]
  \item $A + B = B + A$
  \item $A + (B + C) = (A + B) + C$
  \item $A + O = O + A$
  \item Existe la matriz opuesta de la matriz $A$, que se obtiene cambiando los signos de
  todos los elementos de $A$, tal que $A + (-A) = O$.
\end{itemize}
\end{minipage}

\medskip
\textbf{Producto de un escalar por una matriz:} Dado un número real $k$ y la matriz
$A_{m\times n}$, el producto $k\cdot A$ es otra matriz del mismo orden, que resulta de
multiplicar cada elemento de $A$ por $k$.
\[
k\cdot A = k\cdot\begin{bmatrix}
a_{11} & a_{12} & \cdots & a_{1n}\\
a_{21} & a_{22} & \cdots & a_{2n}\\
a_{31} & a_{32} & \cdots & a_{3n}\\
\vdots & \vdots & \ddots & \vdots\\
a_{m1} & a_{m2} & \cdots & a_{mn}
\end{bmatrix}
=
\begin{bmatrix}
k a_{11} & k a_{12} & \cdots & k a_{1n}\\
k a_{21} & k a_{22} & \cdots & k a_{2n}\\
k a_{31} & k a_{32} & \cdots & k a_{3n}\\
\vdots & \vdots & \ddots & \vdots\\
k a_{m1} & k a_{m2} & \cdots & k a_{mn}
\end{bmatrix}
\]
\textbf{Propiedades:} $\forall \alpha, \beta \in \mathbb{R}$ y $\forall A = (a_{ij})_{m\times n}$,
$B = (b_{ij})_{m\times n}$, se cumple:
\begin{multicols}{2}
\begin{itemize}[label=\textcolor{naranja}{\checkmark}]
  \item $(\alpha\cdot\beta)A = \alpha(\beta\cdot A)$
  \item $\alpha\cdot A = A\cdot\alpha$
  \item $\alpha(A + B) = \alpha A + \alpha B$
  \item $A(\alpha + \beta) = A\alpha + A\beta$
\end{itemize}
\end{multicols}

\textbf{Producto de matrices:} Dadas dos matrices $A_{m\times n}$ y $B_{n\times p}$ ellas son
compatibles para la multiplicación de $A$ por $B$, si el número de columnas de $A$ es igual
al número de filas de $B$. El producto $A\cdot B$ es la matriz $C$ de orden $m\times p$, tal
que los elementos $c_{ij}$ de $C$ es \hl{$c_{ij} = \displaystyle\sum_{k} a_{ik}b_{kj}$} para
cada $i, j$.

\begin{minipage}[t]{0.533\textwidth}
\[
A = \underbrace{\begin{bmatrix} 6 & 4\\ -1 & 3\\ 5 & -2\end{bmatrix}}_{3\times2}
\qquad
B = \underbrace{\begin{bmatrix} 1 & 6 & 3\\ 5 & -10 & -7\end{bmatrix}}_{2\times3}
\]
\end{minipage}\hfill
\begin{minipage}[t]{0.436\textwidth}
\textbf{Propiedades:}
\begin{itemize}[label=\textcolor{naranja}{\checkmark}]
  \item $A(BC) = (A\cdot B)C$
  \item $A(B + C) = A\cdot B + A\cdot C$
  \item $(\alpha\cdot A)^t = \alpha\cdot A^t;\ \alpha \in \mathbb{R}$
  \item $(A\cdot B)^t = B^t\cdot A^t$
  \item Si $A\cdot B = 0$ no implica que $A = 0$ ó $B = 0$.
  \item Si $A\cdot B = A\cdot C$ no implica necesariamente que $B = C$.
  \item El producto de matrices no siempre es conmutativo: $A\cdot B \neq B\cdot A$.
\end{itemize}
\end{minipage}
\[
A\cdot B = \underbrace{\begin{pmatrix}
6\cdot1 + 4\cdot5 & 6\cdot6 + 4(-10) & 6\cdot3 + 4(-7)\\
(-1)1 + 3\cdot5 & (-1)6 + 3(-10) & (-1)3 + 3(-7)\\
5\cdot1 + (-2)5 & 5\cdot6 + (-2)(-10) & 5\cdot3 + (-2)(-7)
\end{pmatrix}}_{3\times3}
\]


% ============================================================
% Páginas 28-30
% ============================================================
\subsection*{Otros tipos de matrices}
\begin{itemize}[label={}]
  \item \textbf{Matriz regular:} Es una matriz cuadrada que tiene inversa.
  \item \textbf{Matriz singular:} No tiene matriz inversa.
  \item \textbf{Matriz idempotente:} Si $A^2 = A$
  \item \textbf{Matriz nilpotente:} Si $A^k = 0 \;\wedge\; A^{k-1} \neq 0$
  \item \textbf{Matriz involutiva:} Si $A^2 = I$
  \item \textbf{Matriz ortogonal:} Si $A\cdot A^t = I$
  \item \textbf{Matrices normales:} Si conmutan con su traspuesta, esto es, si
  $A\cdot A^t = A^t\cdot A$. Obviamente, si $A$ es simétrica, antisimétrica u ortogonal,
  es necesariamente normal.
\end{itemize}

\subsection*{Determinantes}
El determinante es un número real o un escalar que se asocia a cada matriz cuadrada.

\textbf{Notación:} Sea $A$ una matriz cuadrada, el determinante de $A$, se denota por $|A|$,
$\operatorname{Det}(A)$, $\Delta_A$.

\begin{minipage}[t]{0.339\textwidth}
\textbf{Matriz de orden uno:}

Sea $A = [a_{11}]$ entonces:
\[|A| = a_{11}\]
\end{minipage}\hfill
\begin{minipage}[t]{0.630\textwidth}
\textbf{Matriz de orden dos:}

Sea $A = \begin{bmatrix} a_{11} & a_{12}\\ a_{21} & a_{22}\end{bmatrix}$ entonces:
\[|A| = a_{11}\cdot a_{22} - a_{21}\cdot a_{12}\]
\end{minipage}

\medskip
\textbf{Matriz de orden tres} (regla de Sarrus):

\begin{minipage}{0.339\textwidth}
\[
A = \begin{bmatrix}
a_{11} & a_{12} & a_{13}\\
a_{21} & a_{22} & a_{23}\\
a_{31} & a_{32} & a_{33}
\end{bmatrix}
\]
\[
\begin{array}{ccc}
a_{11} & a_{12} & a_{13}\\
a_{21} & a_{22} & a_{23}\\
a_{31} & a_{32} & a_{33}\\ \hline
a_{11} & a_{12} & a_{13}\\
a_{21} & a_{22} & a_{23}
\end{array}
\]
\end{minipage}\hfill
\begin{minipage}{0.630\textwidth}
entonces:
\begin{align*}
|A| = {}& a_{11}\cdot a_{22}\cdot a_{33} + a_{21}\cdot a_{32}\cdot a_{13} + a_{31}\cdot a_{12}\cdot a_{23}\\
& - a_{13}\cdot a_{22}\cdot a_{31} - a_{23}\cdot a_{32}\cdot a_{11} - a_{33}\cdot a_{12}\cdot a_{21}
\end{align*}
\end{minipage}

\medskip
\textbf{Teoremas:} Dadas las matrices cuadradas $A$, $B$ de orden $n$ y el escalar $k$:

\begin{minipage}[t]{0.291\textwidth}
\begin{itemize}[label=\textcolor{naranja}{\checkmark}]
  \item $|AB| = |A||B|$
  \item $|A^t| = |A|$
  \item $|A + B| \neq |A| + |B|$
\end{itemize}
\end{minipage}\hfill
\begin{minipage}[t]{0.679\textwidth}
\begin{itemize}[label=\textcolor{naranja}{\checkmark}]
  \item El determinante no varía si a todos los elementos de una de sus filas (o
  columnas) se suman o restan un múltiplo de otra fila (o columna).
  \item Si dos filas (o columnas) son proporcionales el valor del determinante es
  \hl{cero}.
\end{itemize}
\end{minipage}
\begin{itemize}[label=\textcolor{naranja}{\checkmark}]
  \item Si los elementos de una fila (o columna) son \hl{ceros} el valor del determinante
  es \hl{cero}.
  \item Cuando los elementos de una fila (o columna) tienen factor común, este factor
  puede salir a multiplicar al determinante.
  \item Cuando se permutan dos filas (o columnas) el determinante cambia de signo.
  \item El determinante de una matriz diagonal, triangular superior o inferior, es igual
  al producto de los elementos de la diagonal principal.
  \item El determinante de una matriz antisimétrica \hl{de orden impar}, es igual a
  \hl{cero}.
  \item $|k\cdot A| = k^n|A|$
  \item Determinante de VANDERMONDE:
  \[
  \begin{vmatrix} 1 & 1 & 1\\ a & b & c\\ a^2 & b^2 & c^2\end{vmatrix} = (c-a)(c-b)(b-a)
  \]
\end{itemize}

\textbf{Menor de un elemento:} Se llama menor $M_{ij}$ de un elemento $a_{ij}$ al
determinante de orden ``$n-1$'' que resulta al suprimir en la matriz $A$ de orden $n$, todos
los elementos de la fila $i$ y la columna $j$. Dada la matriz:
\[
A = \begin{pmatrix} -1 & 2 & 1\\ 1 & 0 & -1\\ 2 & 1 & 0\end{pmatrix}
\qquad\text{Luego:}\quad
M_{21} = \begin{vmatrix} 2 & 1\\ 1 & 0\end{vmatrix} = -1
\qquad
M_{33} = \begin{vmatrix} -1 & 2\\ 1 & 0\end{vmatrix} = -2
\]

\textbf{Cofactor de un elemento:} Se llama cofactor de un elemento \hl{$a_{ij}$} a su
\hl{menor} con su signo o con signo contrario según sea par o impar la suma $i + j$. El
cofactor del elemento \hl{$a_{ij}$} denotamos por \hl{$A_{ij}$}. Del ejemplo anterior:
\[
\fcolorbox{verde}{verde!15}{$A_{ij} = (-1)^{i+j}M_{ij}$}
\qquad
A_{21} = (-1)^{2+1}M_{21} = (-1)(-1) = 1
\qquad
A_{33} = (-1)^{3+3}M_{33} = 1(-2) = -2
\]

\textbf{Teorema de Laplace:} Sea: $A = [a_{ij}]_n$
\[
\left.
\begin{array}{l}
\checkmark\ \text{Respecto a la fila } i\text{:}\quad |A| = \displaystyle\sum_{j=1}^{n} a_{ij}A_{ij}\\[14pt]
\checkmark\ \text{Respecto a la columna } j\text{:}\quad |A| = \displaystyle\sum_{i=1}^{n} a_{ij}A_{ij}
\end{array}
\right\}\; A_{ij} \text{ es el cofactor de } a_{ij}
\]

Del ejemplo anterior:
\[
A = \begin{pmatrix}
\overset{+}{-1} & \overset{-}{2} & \overset{+}{1}\\
1 & 0 & -1\\
2 & 1 & 0
\end{pmatrix}_{3\times3}
\quad n = 3
\]
\textcolor{magenta}{$\ast$} Respecto a la fila 1:
\[
|A| = \underset{-1}{a_{11}}\underset{(1)}{A_{11}} + \underset{2}{a_{12}}\underset{(-2)}{A_{12}} + \underset{1}{a_{13}}\underset{(1)}{A_{13}} = -1 - 4 + 1 = \boxed{-4}
\]
\[
A_{11} = \begin{vmatrix} 0 & -1\\ 1 & 0\end{vmatrix} = 1
\qquad
A_{12} = \begin{vmatrix} 1 & -1\\ 2 & 0\end{vmatrix}(-1) = -2
\qquad
A_{13} = \begin{vmatrix} 1 & 0\\ 2 & 1\end{vmatrix} = 1
\]

\textbf{Matriz de cofactores:}
\[
\fcolorbox{verde}{verde!15}{$\operatorname{Cof}(A) = \begin{bmatrix}
A_{11} & A_{12} & A_{13}\\
A_{21} & A_{22} & A_{23}\\
A_{31} & A_{32} & A_{33}
\end{bmatrix}$}
\]
Del ejemplo anterior:
\[
\operatorname{Cof}(A) = \begin{pmatrix}
\begin{vmatrix} 0 & -1\\ 1 & 0\end{vmatrix} & -\begin{vmatrix} 1 & -1\\ 2 & 0\end{vmatrix} & \begin{vmatrix} 1 & 0\\ 2 & 1\end{vmatrix}\\[12pt]
-\begin{vmatrix} 2 & 1\\ 1 & 0\end{vmatrix} & \begin{vmatrix} -1 & 1\\ 2 & 0\end{vmatrix} & -\begin{vmatrix} -1 & 2\\ 2 & 1\end{vmatrix}\\[12pt]
\begin{vmatrix} 2 & 1\\ 0 & -1\end{vmatrix} & -\begin{vmatrix} -1 & 1\\ 1 & -1\end{vmatrix} & \begin{vmatrix} -1 & 2\\ 1 & 0\end{vmatrix}
\end{pmatrix}
\;\Rightarrow\;
\operatorname{Cof}(A) = \begin{pmatrix} 1 & -2 & 1\\ 1 & -2 & 5\\ -2 & 0 & -2\end{pmatrix}
\]

\textbf{Adjunta de una matriz:}
\[
\fcolorbox{verde}{verde!15}{$\operatorname{Adj}(A) = (\operatorname{Cof}(A))^t$}
\qquad \text{Del ejemplo anterior:}\quad
\operatorname{Adj}(A) = \begin{pmatrix} 1 & 1 & -2\\ -2 & -2 & 0\\ 1 & 5 & -2\end{pmatrix}
\]

\textbf{Matriz inversa:} Sea $A$ una matriz cuadrada regular o no singular, se llama matriz
inversa de $A$ denotada por $A^{-1}$ (que) verifica la relación:
\[\fbox{$A\cdot A^{-1} = A^{-1}\cdot A = I$} \qquad \textcolor{rojo}{I}\text{: Matriz identidad.}\]
En general para toda matriz $A$ no singular, la inversa se puede obtener:
\[\fbox{$A^{-1} = \dfrac{1}{|A|}\cdot\operatorname{Adj}(A)$}\]
Del ejemplo anterior:
\[
A^{-1} = \frac{1}{(-4)}\begin{bmatrix} 1 & 1 & -2\\ -2 & -2 & 0\\ 1 & 5 & -2\end{bmatrix}
\longrightarrow
A^{-1} = \begin{bmatrix} -1/4 & -1/4 & 1/2\\ 1/2 & 1/2 & 0\\ -1/4 & -5/4 & 1/2\end{bmatrix}
\]
Para una matriz de \hl{orden 2} podemos utilizar:
\[
A = \begin{bmatrix} a & b\\ c & d\end{bmatrix}
\longrightarrow
A^{-1} = \frac{1}{|A|}\begin{bmatrix} d & -b\\ -c & a\end{bmatrix}
\]


% ============================================================
% Páginas 31-33
% ============================================================
\subsection*{Método de Gauss--Jordan}
Se trata de plantearla como una ecuación matricial donde el objetivo es encontrar
$A^{-1}$. Para ello seguiremos los siguientes pasos:
\begin{enumerate}
  \item Escribimos una matriz doble que contiene a la matriz $A$ en un lado y a la matriz
  identidad $I$ en el otro.
  \item Realizamos \hl{operaciones elementales fila} para transformar la matriz $A$ en la
  identidad.
  \item El procedimiento termina cuando la matriz identidad está al lado izquierdo y el
  lado derecho es la matriz $B$, que precisamente, es la inversa de $A$: $B = A^{-1}$.
\end{enumerate}

Llamaremos operaciones elementales fila (sobre una matriz) a las siguientes operaciones:
\begin{enumerate}
  \item Multiplicar una fila por un número distinto de $0$.
  \item Sumar (o restar) a una fila, el múltiplo de otra fila.
  \item Intercambiar el orden de las filas.
\end{enumerate}

\textbf{Ejemplos:}

\textbf{1.} Hallar la matriz inversa de: \quad $A = \begin{pmatrix} 8 & 2\\ 7 & 2\end{pmatrix}$
\[
\left(\begin{array}{cc|cc} 8 & 2 & 1 & 0\\ 7 & 2 & 0 & 1\end{array}\right)
\xrightarrow[\;8f_2 - 7f_1\;]{\;f_1 - f_2\;}
\left(\begin{array}{cc|cc} 1 & 0 & 1 & -1\\ 0 & 2 & -7 & 8\end{array}\right)
\xrightarrow{\;f_2/2\;}
\left(\begin{array}{cc|cc} 1 & 0 & 1 & -1\\ 0 & 1 & -\frac{7}{2} & 4\end{array}\right)
\]
\[A^{-1} = \begin{pmatrix} 1 & -1\\ -\frac{7}{2} & 4\end{pmatrix}\]

\textbf{2.} Hallar la matriz inversa de: \quad
$A = \begin{pmatrix} -1 & 2 & 1\\ 1 & 0 & -1\\ 2 & 1 & 0\end{pmatrix}$
\[
\left(\begin{array}{ccc|ccc}
-1 & 2 & 1 & 1 & 0 & 0\\
1 & 0 & -1 & 0 & 1 & 0\\
2 & 1 & 0 & 0 & 0 & 1
\end{array}\right)
\xrightarrow{\substack{f_1 - 2f_3\\ f_2 + f_1\\ f_3 + 2f_1}}
\left(\begin{array}{ccc|ccc}
-5 & 0 & 1 & 1 & 0 & -2\\
0 & 2 & 0 & 1 & 1 & 0\\
0 & 5 & 2 & 2 & 0 & 1
\end{array}\right)
\]
\[
\xrightarrow{\;2f_3 - 5f_2\;}
\left(\begin{array}{ccc|ccc}
-5 & 0 & 1 & 1 & 0 & -2\\
0 & 2 & 0 & 1 & 1 & 0\\
0 & 0 & 4 & -1 & -5 & 2
\end{array}\right)
\xrightarrow{\;4f_1 - f_3\;}
\left(\begin{array}{ccc|ccc}
-20 & 0 & 0 & 5 & 5 & -10\\
0 & 2 & 0 & 1 & 1 & 0\\
0 & 0 & 4 & -1 & -5 & 2
\end{array}\right)
\]
\[
\xrightarrow{\substack{f_1/(-20)\\ f_2/2\\ f_3/4}}
\left(\begin{array}{ccc|ccc}
1 & 0 & 0 & -1/4 & -1/4 & 1/2\\
0 & 1 & 0 & 1/2 & 1/2 & 0\\
0 & 0 & 1 & -1/4 & -5/4 & 1/2
\end{array}\right)
\quad\Rightarrow\quad
A^{-1} = \begin{pmatrix} -1/4 & -1/4 & 1/2\\ 1/2 & 1/2 & 0\\ -1/4 & -5/4 & 1/2\end{pmatrix}
\]

\textbf{Teoremas:} $A$, $B$ de orden ``$n$'', no singulares.
\begin{multicols}{2}
\begin{itemize}[label=\textcolor{naranja}{\checkmark}]
  \item $|A^{-1}| = |A|^{-1}$
  \item $(A + B)^{-1} \neq A^{-1} + B^{-1}$
  \item $(A^{-1})^T = (A^T)^{-1}$
  \item $|\operatorname{Adj}(\operatorname{Adj}(A))| = |A|^{(n-1)^2}$
  \item $(A^{-1})^{-1} = A$
  \item $(A\cdot B)^{-1} = B^{-1}\cdot A^{-1}$
  \item $(\lambda A)^{-1} = \lambda^{-1}\cdot A^{-1}$
  \item Si: $A\cdot X\cdot B = C \rightarrow X = A^{-1}\cdot C\cdot B^{-1}$
\end{itemize}
\end{multicols}

\newpage
% ------------------------------------------------------------
\section{Teoría de ecuaciones I}

\textbf{Ecuación:} Es una igualdad que se establece entre dos expresiones matemáticas de
una o más variables.
\[\fbox{$A(x,y) = B(x,y)$} \qquad \textcolor{magenta}{x, y\text{: incógnitas}}\]

\textbf{Solución de una ecuación:} Es aquel valor que toma la incógnita y convierte a la
ecuación en una proposición verdadera.

\textbf{Conjunto solución:} Es un conjunto formado por todas las soluciones de la ecuación.

Ejemplos: Resuelva la ecuación $x^2 - 4 = 3x$
\[
\begin{array}{lll}
x^2 - 3x - 4 = 0 \qquad & x_1 = 4 \qquad & \therefore\ C.S. = \{-1; 4\}\\
(x-4)(x+1) = 0 & x_2 = -1 &
\end{array}
\]

\textbf{Clasificación de las ecuaciones:} De acuerdo al número de soluciones, las
ecuaciones se clasifican en:

\textcolor{azul}{\textbf{Ecuación compatible:}} Es aquella que admite al menos una solución,
a su vez pueden ser:
\begin{itemize}[label={}]
  \item \textbf{E.C. Determinada:} Es aquella que admite un número finito de soluciones.

  Ejemplo: \quad $(x-4)(x-7)(x+11) = 0 \qquad \therefore\ C.S. = \{-11; 4; 7\}$

  \item \textbf{E.C. Indeterminada:} Es aquella que admite infinitas soluciones.

  \begin{minipage}{0.5\linewidth}
  Ejemplo:
  \begin{gather*}
  (x+5)(x+1) + 4 = (x+3)^2\\
  x^2 + 6x + 9 = x^2 + 6x + 9
  \end{gather*}
  \end{minipage}%
  \begin{minipage}{0.45\linewidth}
  Esta igualdad es correcta para todo valor real de $x$, por tanto:
  \[\therefore\ C.S. = \mathbb{R}\]
  \end{minipage}
\end{itemize}

\textcolor{azul}{\textbf{Ecuación incompatible:}} Es aquella que no admite solución, es
decir, su conjunto solución es el vacío.

\begin{minipage}{0.388\textwidth}
Ejemplo: \quad $\dfrac{3}{x+1} = 0$
\end{minipage}\hfill
\begin{minipage}{0.582\textwidth}
Esta igualdad no es correcta para cualquier valor de $x$, por tanto, su conjunto solución
no tiene elementos, es decir:
\[\therefore\ C.S. = \varnothing\]
\end{minipage}


% ============================================================
% Páginas 34-36
% ============================================================
\subsection*{Ecuaciones polinomiales}

\textcolor{azul}{\textbf{Ecuación lineal:}} Es aquella ecuación que tiene la siguiente forma
general:
\[\ovalbox{$P(x) = ax + b = 0$} \qquad \forall a \neq 0;\ a \wedge b \in \mathbb{R}\]
Cuya raíz o solución es: \ \fcolorbox{azul}{azul!10}{$x = \dfrac{-b}{a}$}
\begin{itemize}[label=\checkmark]
  \item Si: $a \neq 0 \longrightarrow$ E. Lineal compatible determinada.
  \item Si: $a = 0 \wedge b = 0 \longrightarrow$ E. compatible indeterminada.
  \item Si: $a = 0 \wedge b \neq 0 \longrightarrow$ E. incompatible.
\end{itemize}

\textcolor{azul}{\textbf{Ecuación cuadrática:}} Es aquella ecuación que tiene la siguiente
forma general:
\[\ovalbox{$P(x) = ax^2 + bx + c = 0$} \qquad \forall a \neq 0;\ a \wedge b \wedge c \in \mathbb{R}\]
Presenta dos raíces ($x_1; x_2$) pero una o dos soluciones; para hallar las raíces de una
ecuación cuadrática se recomienda primero intentar factorizar por aspa simple, agrupación
o diferencia de cuadrados.

Cuando no se puede factorizando, las raíces se obtienen por la siguiente
\hl{fórmula general}:

\begin{minipage}{0.582\textwidth}
\[
\text{\textcolor{naranja}{(Fórmula de Bhaskara)}}\quad
\ovalbox{$x_{1,2} = \dfrac{-b \pm \sqrt{b^2 - 4ac}}{2a}$}
\]
\end{minipage}\hfill
\begin{minipage}{0.388\textwidth}
donde: \ $\Delta = b^2 - 4ac$

\textcolor{rojo}{$\Delta$}: Discriminante
\end{minipage}

\medskip
\textbf{Naturaleza de las raíces:} Dada la ecuación cuadrática de coeficientes reales:
$ax^2 + bx + c = 0$
\begin{enumerate}[label=\Roman*.]
  \item Si $\Delta > 0$ las raíces son reales y diferentes.
  \item Si $\Delta = 0$ las raíces son reales e iguales.
  \item Si $\Delta < 0$ las raíces son imaginarias y conjugadas.
\end{enumerate}

\textbf{Teoremas de las raíces:} Sean: $x_1; x_2$ las raíces de la cuadrática:
$ax^2 + bx + c = 0$

\begin{minipage}[t]{0.504\textwidth}
\begin{enumerate}
  \item Suma de raíces: \ \ovalbox{$x_1 + x_2 = -\dfrac{b}{a}$}
  \item Producto de raíces: \ \ovalbox{$x_1\cdot x_2 = \dfrac{c}{a}$}
  \item Diferencia de raíces: \ \ovalbox{$|x_1 - x_2| = \dfrac{\sqrt{\Delta}}{|a|}$}

  \ovalbox{$(x_1 + x_2)^2 - (x_1 - x_2)^2 = 4x_1\cdot x_2$}
  \item Raíces simétricas: \ \ovalbox{$x_1 + x_2 = 0$}

  \textcolor{rojo}{$\ast$} $b = 0$
\end{enumerate}
\end{minipage}\hfill
\begin{minipage}[t]{0.466\textwidth}
\begin{enumerate}\setcounter{enumi}{4}
  \item Raíces recíprocas: \ \ovalbox{$x_1\cdot x_2 = 1$}

  \textcolor{rojo}{$\ast$} $a = c$
  \item $x_1^2 + x_2^2 = (x_1 + x_2)^2 - 2x_1x_2$
  \item $x_1^3 + x_2^3 = (x_1 + x_2)^3 - 3x_1x_2(x_1 + x_2)$
  \item Reconstrucción de una ecuación cuadrática:
  \[\ovalbox{$x^2 - (x_1 + x_2)x + x_1x_2 = 0$}\]
\end{enumerate}
\end{minipage}

\subsection*{Teoremas}

\begin{minipage}[t]{0.436\textwidth}
$\ast$ \textbf{Ecuaciones equivalentes:}
\[
\begin{cases}
a_1x^2 + b_1x + c_1 = 0\\
a_2x^2 + b_2x + c_2 = 0
\end{cases}
\]
Se cumple: \ \ovalbox{$\dfrac{a_1}{a_2} = \dfrac{b_1}{b_2} = \dfrac{c_1}{c_2}$}

\medskip
$\ast$ Cuando poseen el mismo conjunto solución.
\end{minipage}\hfill
\begin{minipage}[t]{0.533\textwidth}
$\ast$ \textbf{Teorema de Bezout:}
\[
\left.\begin{array}{l}
a_1x^2 + b_1x + c_1 = 0\\
a_2x^2 + b_2x + c_2 = 0
\end{array}\right\}
\text{\small Ambas presentan una raíz común}
\]
Se cumple:
\[\ovalbox{$(a_1c_2 - a_2c_1)^2 = (a_1b_2 - a_2b_1)(b_1c_2 - b_2c_1)$}\]

\textbf{Obs.:} Supongamos que ``$m$'' es una raíz común.
\begin{gather*}
x^2 + \frac{b_1}{a_1}x + \frac{c_1}{a_1} = 0 \qquad x^2 + \frac{b_2}{a_2}x + \frac{c_2}{a_2} = 0\\
m^2 + \frac{b_1}{a_1}m + \frac{c_1}{a_1} = 0 \quad (-) \quad m^2 + \frac{b_2}{a_2}m + \frac{c_2}{a_2} = 0\\
m\left(\frac{b_1}{a_1} - \frac{b_2}{a_2}\right) + \frac{c_1}{a_1} - \frac{c_2}{a_2} = 0\\
m(a_2b_1 - a_1b_2) + a_2c_1 - a_1c_2 = 0\\
\therefore\ m = \frac{a_1c_2 - a_2c_1}{a_2b_1 - a_1b_2}
\end{gather*}
\end{minipage}

\newpage
% ------------------------------------------------------------
\section{Teoría de ecuaciones II}

\subsection*{Algunas definiciones previas}

\textbf{Definición 1:} El número ``$\alpha$'' es una raíz o cero del polinomio
$P(x) \longleftrightarrow P(\alpha) = 0$

\textbf{Teorema del factor:} ``$\alpha$'' es una raíz o cero del polinomio
$P(x) \longleftrightarrow (x - \alpha)$ es factor de $P(x)$.

\textbf{Definición 2:} El número ``$\alpha$'' es una raíz de multiplicidad ``$k$'' de un
polinomio $P(x) \longleftrightarrow P(x)$ es divisible por $(x-\alpha)^k$ pero no por
$(x-\alpha)^{k+1}$, es decir:
\[\ovalbox{$P(x) = (x-\alpha)^k\cdot Q(x)\;;\quad Q(\alpha) \neq 0$}\]

\textbf{Teorema fundamental del álgebra:} Un polinomio de grado ``$n$'', con coeficientes
complejos tiene por lo menos una raíz generalmente compleja.

\textbf{Corolario:} Todo polinomio de grado ``$n$'' tiene exactamente ``$n$'' raíces
contadas las multiplicidades.

\textbf{Ecuación polinomial de grado superior:} Sea:
\ \ovalbox{$P(x) = a_0x^n + a_1x^{n-1} + a_2x^{n-2} + \cdots + a_n = 0$}

$n \in \mathbb{N} \wedge$ (Coeficiente principal $a_0 \neq 0$) $\wedge$ Sus raíces son:
$x_1, x_2, x_3, \ldots, x_n$

Luego: \quad \ovalbox{$P(x) = a_0(x - x_1)(x - x_2)(x - x_3)\cdots(x - x_n) = 0$}

En una ecuación polinomial de grado ``$n$'', que tiene exactamente ``$n$'' raíces:
$x_1; x_2; x_3; \ldots; x_n$ como se trata de una ecuación también tiene soluciones, en este
tipo de ecuaciones se cumple:
\[\ovalbox{N.\textsuperscript{o} raíces $\geq$ N.\textsuperscript{o} soluciones}\]

\textbf{Ejemplo:} Calcule las raíces y soluciones de la ecuación:
\ $P(x) = (x-2)(x-9)^2(x+6)^3 = 0$

Como $^{\circ}[P(x)]$ es 6, entonces tiene 6 raíces. Igualando cada factor a cero:

\begin{minipage}{0.630\textwidth}
\begin{itemize}
  \item $x - 2 = 0 \rightarrow x = 2$: raíz simple
  \item $(x-9)^2 = 0 \rightarrow x = 9$: raíz doble o de multiplicidad 2.
  \item $(x+6)^3 = 0 \rightarrow x = -6$: raíz triple o de multiplicidad 3.
\end{itemize}
\end{minipage}\hfill
\begin{minipage}{0.339\textwidth}
$\left.\right\}$ $C.S. = \{2; 9; -6\}$ \\ tiene 3 soluciones

\smallskip
{\small $\ast$ Para calcular las raíces de un polinomio o ecuación se factoriza y se iguala
cada factor a cero.}
\end{minipage}


% ============================================================
% Páginas 37-39
% ============================================================
\subsection*{Teoremas de Cardano--Viette}

\begin{minipage}[t]{0.436\textwidth}
\cn{1} \textcolor{magenta}{Suma de raíces:}
\[\ovalbox{$x_1 + x_2 + x_3 + \cdots + x_n = -\dfrac{a_1}{a_0}$}\]
\end{minipage}\hfill
\begin{minipage}[t]{0.533\textwidth}
\cn{2} \textcolor{magenta}{Suma de productos binarios de las raíces:}
\[\ovalbox{$x_1x_2 + x_1x_3 + x_2x_3 + \cdots + x_{n-1}x_n = \dfrac{a_2}{a_0}$}\]
\end{minipage}

\medskip
\begin{minipage}[t]{0.582\textwidth}
\cn{3} \textcolor{magenta}{Suma de productos ternarios de las raíces:}
\[\ovalbox{$x_1x_2x_3 + x_1x_2x_4 + \cdots + x_{n-2}x_{n-1}x_n = -\dfrac{a_3}{a_0}$}\]
\[\vdots\]
\end{minipage}\hfill
\begin{minipage}[t]{0.388\textwidth}
\fbox{\parbox{0.9\linewidth}{\textbf{En general:} Si ``$S_k$'' representa la suma de los
productos de las raíces tomadas de ``$k$'' en ``$k$'' se cumple:
\[\ovalbox{$S_k = (-1)^k\cdot\dfrac{a_k}{a_0}$}\]}}
\end{minipage}

\cn{4} \textcolor{magenta}{Producto de raíces:}
\[\ovalbox{$x_1\cdot x_2\cdot x_3\cdot x_4\cdots x_n = (-1)^n\dfrac{a_n}{a_0}$} \qquad
\textcolor{magenta}{n\text{: depende del grado de la ecuación.}}\]

\textcolor{naranja}{\textbf{Ecuación cúbica:}} Es aquella que tiene la siguiente forma general:
\[\ovalbox{$ax^3 + bx^2 + cx + d = 0$} \qquad \forall a \neq 0;\ a, b, c \wedge d \in \mathbb{R}\]
Sus raíces son: $x_1, x_2, x_3$
\[
\cn{1}\ x_1 + x_2 + x_3 = -\frac{b}{a} \qquad
\cn{2}\ x_1x_2 + x_1x_3 + x_2x_3 = \frac{c}{a} \qquad
\cn{3}\ x_1\cdot x_2\cdot x_3 = -\frac{d}{a}
\]

\textcolor{verde}{\textbf{Obs.:}}

\begin{minipage}{0.388\textwidth}
\begin{gather*}
x^3 + px + q = 0\\
D = \left(\frac{p}{3}\right)^3 + \left(\frac{q}{2}\right)^2
\end{gather*}
\textcolor{rojo}{$D$}: Discriminante
\end{minipage}\hfill
\begin{minipage}{0.582\textwidth}
\begin{itemize}[label=\checkmark]
  \item Si $D > 0$: 1 raíz real y 2 imaginarias.
  \item Si $D = 0$: 3 raíces reales (2 raíces iguales).
  \item Si $D < 0$: 3 raíces reales distintas.
\end{itemize}
\end{minipage}

\medskip
\textcolor{naranja}{\textbf{Ecuación bicuadrada:}} Ecuación de cuarto grado, sin términos de
grado impar y presenta dos pares de raíces simétricas. Tiene la siguiente forma general:
\[\ovalbox{$ax^4 + bx^2 + c = 0$} \qquad \forall a \neq 0;\ a, b \wedge c \in \mathbb{R}\]
La resolución se realiza factorizando o mediante un cambio de variable para transformarla
en una ecuación cuadrática.

También presenta una fórmula general:
\[x_{1,2,3,4} = \pm\sqrt{\frac{-b \pm \sqrt{b^2 - 4ac}}{2a}}\]
Sus raíces son: \ $\underset{\alpha}{x_1},\ \underset{-\alpha}{x_2},\ \underset{\beta}{x_3},\ \underset{-\beta}{x_4}$

\begin{minipage}[t]{0.388\textwidth}
\cn{1} $x_1 + x_2 + x_3 + x_4 = 0$
\[\bullet\ \alpha - \alpha + \beta - \beta = 0\]

\cn{3} $x_1\cdot x_2\cdot x_3\cdot x_4 = \dfrac{c}{a}$
\[\bullet\ (\alpha)(-\alpha)(\beta)(-\beta) = \frac{c}{a}\]
\[\ovalbox{$\alpha^2\cdot\beta^2 = \dfrac{c}{a}$}\]
\end{minipage}\hfill
\begin{minipage}[t]{0.582\textwidth}
\cn{2} $x_1x_2 + x_1x_3 + x_1x_4 + x_2x_3 + x_2x_4 + x_3x_4 = \dfrac{b}{a}$
\begin{gather*}
\bullet\ -\alpha^2 + \alpha\beta - \alpha\beta - \alpha\beta + \alpha\beta - \beta^2 = \frac{b}{a}\\
-\alpha^2 - \beta^2 = \frac{b}{a} \;\rightarrow\; \ovalbox{$\alpha^2 + \beta^2 = -\dfrac{b}{a}$}
\end{gather*}

\cn{4} \textcolor{magenta}{Reconstrucción de una ecuación bicuadrada:}
\begin{gather*}
x^4 + [x_1\cdot x_2 + x_3\cdot x_4]x^2 + [x_1\cdot x_2\cdot x_3\cdot x_4] = 0\\
\bullet\ x^4 - (\alpha^2 + \beta^2)x^2 + \alpha^2\cdot\beta^2 = 0;\quad \alpha + \beta \neq 0
\end{gather*}
\end{minipage}

\medskip
$\ast$ Si las raíces de una ecuación bicuadrada están en P.A., sus raíces tienen la forma:
\[\underbrace{-3m,\ -m}_{2m},\ \underbrace{m,}_{2m}\ \underbrace{3m}_{2m}\,; \quad
\text{También: } ax^4 + bx^2 + c = 0 \;\rightarrow\; \ovalbox{$9b^2 = 100ac$}\]

\textcolor{naranja}{\textbf{Ecuaciones recíprocas:}} Se caracterizan porque los coeficientes
de los términos equidistantes de los extremos son iguales. Su forma general es:
\[\ovalbox{$Ax^n + Bx^{n-1} + Cx^{n-2} + \cdots + Cx^2 + Bx + A = 0$}\]
Ejemplo: \quad $3x^5 - 2x^4 + 6x^3 + 6x^2 - 2x + 3 = 0$
\[\text{Sus raíces son de la forma: } \ovalbox{$\alpha, \dfrac{1}{\alpha}, \beta, \dfrac{1}{\beta}, \theta$}\]
Cuando es de grado impar, una raíz es $1$ ó $-1$; si es de grado par, aplicar la
factorización recíproca.

\textcolor{magenta}{\textbf{Propiedad:}} Si $P(x)$ es un polinomio recíproco:
\[\ovalbox{$P(x) \equiv x^n\cdot P\!\left(\dfrac{1}{x}\right)$}\]

\subsection*{Teoremas de paridad de raíces}

\textcolor{verde}{\textbf{Paridad de raíces imaginarias:}} En una ecuación polinomial de
coeficientes reales, si una raíz imaginaria es ``$a + bi$'' entonces la otra raíz es la
conjugada ``$a - bi$''.
\begin{itemize}
  \item Si la ecuación polinomial de coeficientes reales es de grado impar, tendrá por lo
  menos una raíz real.
\end{itemize}

\textcolor{verde}{\textbf{Paridad de raíces irracionales:}} En una ecuación polinomial de
coeficientes racionales, si una raíz irracional es ``$a + \sqrt{b}$'' entonces la otra raíz
es la conjugada ``$a - \sqrt{b}$'', con $a, b \in \mathbb{Q}$ y $\sqrt{b} \in$ irracionales.
\begin{itemize}
  \item En una ecuación polinomial de coeficientes racionales, si una raíz es
  ``$\sqrt{a} + \sqrt{b}$'' entonces las otras raíces son la conjugada, la opuesta y la
  conjugada de la opuesta: $\sqrt{a} - \sqrt{b};\ -\sqrt{a} - \sqrt{b};\ -\sqrt{a} + \sqrt{b}$,
  con $a, b \in \mathbb{Q}^+$; $\sqrt{a}$ y $\sqrt{b} \in$ irracionales.
\end{itemize}

\subsection*{Transformadas de las raíces}
Sea la ecuación polinomial:
$a_0x^n + a_1x^{n-1} + a_2x^{n-2} + \cdots + a_{n-1}x + a_n = 0;\ a_0 \neq 0$ con raíces
$\{x_1; x_2; x_3; \ldots; x_n\}$ entonces:

\begin{enumerate}[label=\cn{\arabic*}]
  \item La ecuación de raíces aumentadas o disminuidas en un valor ``$k$'', es decir con
  raíces: $\{x_1 \pm k; x_2 \pm k; x_3 \pm k; \ldots; x_n \pm k\}$ es:
  \[\ovalbox{$a_0(x \mp k)^n + a_1(x \mp k)^{n-1} + a_2(x \mp k)^{n-2} + \cdots + a_{n-1}(x \mp k) + a_n = 0$}\]
  \item La ecuación de raíces multiplicadas por un valor ``$k$'' $(k \neq 0)$; es decir con
  raíces: $\{x_1\cdot k; x_2\cdot k; x_3\cdot k; \ldots; x_n\cdot k\}$ es:
  \[\ovalbox{$a_0\left(\dfrac{x}{k}\right)^n + a_1\left(\dfrac{x}{k}\right)^{n-1} + \cdots + a_{n-1}\left(\dfrac{x}{k}\right) + a_n = 0$}\]
  \item La ecuación de raíces invertidas, es decir con raíces:
  $\left\{\dfrac{1}{x_1}; \dfrac{1}{x_2}; \dfrac{1}{x_3}; \ldots; \dfrac{1}{x_n}\right\}$ es:
  \[\ovalbox{$a_nx^n + a_{n-1}x^{n-1} + \cdots + a_1x + a_0 = 0$}\]
\end{enumerate}


% ============================================================
% Páginas 40-42
% ============================================================
\newpage
\section{Desigualdades I}

\textbf{Desigualdad:} Es una comparación de números reales mediante los símbolos de orden:
$<;\ >;\ \geq;\ \leq$.

\textbf{Definiciones:} Siendo $a \in \mathbb{R}$, se establece:

\begin{minipage}[t]{0.485\textwidth}
\begin{itemize}[label=\textcolor{naranja}{\checkmark}]
  \item $a$ es positivo $\longleftrightarrow a > 0$
  \item $a$ es negativo $\longleftrightarrow a < 0$
  \item $a$ es no positivo $\longleftrightarrow a \leq 0$
  \item $a$ es no negativo $\longleftrightarrow a \geq 0$
\end{itemize}
\end{minipage}\hfill
\begin{minipage}[t]{0.485\textwidth}
\begin{itemize}[label=\textcolor{naranja}{\checkmark}]
  \item $a > b \longleftrightarrow a - b > 0$
  \item $a < b \longleftrightarrow a - b < 0$
  \item $a \geq b \longleftrightarrow a > b \vee a = b$
\end{itemize}
\end{minipage}

\subsection*{Axiomas de orden}

\begin{minipage}[t]{0.485\textwidth}
\cn{1} \textcolor{verde}{\textbf{Ley de tricotomía:}}
\[\forall a \wedge b \in \mathbb{R}:\ a > b \vee a = b \vee a < b\]

\cn{2} \textcolor{verde}{\textbf{Ley aditiva:}}
\[\forall a, b \wedge c \in \mathbb{R}: \text{ Si } a > b \rightarrow a + c > b + c\]
\end{minipage}\hfill
\begin{minipage}[t]{0.485\textwidth}
\cn{3} \textcolor{verde}{\textbf{Ley multiplicativa:}}
\[\forall a, b \in \mathbb{R} \wedge c \in \mathbb{R}^+: \text{ Si } a > b \leftrightarrow ac > bc\]

\cn{4} \textcolor{verde}{\textbf{Ley de transitividad:}}
\[\forall a, b \wedge c \in \mathbb{R}: \text{ Si } a > b \wedge b > c \leftrightarrow a > c\]
\end{minipage}

\subsection*{Intervalos}
Es un subconjunto de los números reales definido entre dos extremos, uno inferior y otro
superior. Existen 2 clases de intervalos.

\textcolor{verde}{\textbf{Acotados:}} Los extremos del intervalo son finitos, a su vez estos
pueden ser:

\textcolor{naranja}{\textbf{I. Intervalo abierto:}}
\[\intervalo{o}{o}{f} \qquad \ovalbox{$]a;b[ \;=\; \langle a;b\rangle = \{x \in \mathbb{R} \,/\, a < x < b\}$}\]

\textcolor{naranja}{\textbf{II. Intervalo cerrado:}}
\[\intervalo{c}{c}{f} \qquad \ovalbox{$[a;b] = \{x \in \mathbb{R} \,/\, a \leq x \leq b\}$}\]

\textcolor{naranja}{\textbf{III. Intervalo semiabierto:}}
\[\intervalo{o}{c}{f} \qquad \ovalbox{$]a;b] = \langle a;b] = \{x \in \mathbb{R} \,/\, a < x \leq b\}$}\]

\textcolor{verde}{\textbf{No acotados:}} Al menos un extremo del intervalo no es finito, también
se denomina intervalo al infinito.
\[\intervalo{o}{o}{i} \qquad \ovalbox{$]a;+\infty[ \;=\; \langle a;+\infty\rangle = \{x \in \mathbb{R} \,/\, x > a\}$}\]

\begin{itemize}
  \item En el caso del conjunto de los números reales se tiene:
  \[\ovalbox{$\mathbb{R} = \;]-\infty;+\infty[\; = \langle -\infty;+\infty\rangle$}\]
\end{itemize}

\subsection*{Operaciones con intervalos}
Si los conjuntos $A$ y $B$ son intervalos, entonces:

\begin{minipage}[t]{0.485\textwidth}
\begin{enumerate}[label=\cn{\arabic*}]
  \item Unión: $A \cup B = \{x \,/\, x \in A \vee x \in B\}$
  \item Intersección: $A \cap B = \{x \,/\, x \in A \wedge x \in B\}$
  \item Diferencia: $A - B = \{x \,/\, x \in A \wedge x \notin B\}$
\end{enumerate}
\end{minipage}\hfill
\begin{minipage}[t]{0.485\textwidth}
\begin{enumerate}[label=\cn{\arabic*}]\setcounter{enumi}{3}
  \item Complemento: $A' = \{x \,/\, x \in \mathbb{R} \wedge x \notin A\}$
  \item Diferencia simétrica:

  $A \,\Delta\, B = \{x \,/\, x \in A - B \vee x \in B - A\}$
\end{enumerate}
\end{minipage}

\subsection*{Teoremas}
Con $a; b; c; d \in \mathbb{R}$

\begin{minipage}[t]{0.485\textwidth}
\begin{enumerate}[label=\cn{\arabic*}]
  \item $\forall x \in \mathbb{R}: x^2 \geq 0$
  \item Si $a < b \wedge c < 0 \rightarrow ac > bc$
  \item Si $a < b < c$; $n \in \mathbb{N}$; cumple:
  \[
  \begin{cases}
  a^{2n+1} < b^{2n+1} < c^{2n+1}\\[4pt]
  \sqrt[2n+1]{a} < \sqrt[2n+1]{b} < \sqrt[2n+1]{c}
  \end{cases}
  \]
\end{enumerate}
\end{minipage}\hfill
\begin{minipage}[t]{0.485\textwidth}
\begin{enumerate}[label=\cn{\arabic*}]\setcounter{enumi}{3}
  \item Si $a > b > c > 0$; $n \in \mathbb{N} \rightarrow a^{2n} > b^{2n} > c^{2n} > 0$
  \item $ab > 0 \leftrightarrow \{a > 0 \wedge b > 0\} \vee \{a < 0 \wedge b < 0\}$
  \item $ab < 0 \leftrightarrow \{a > 0 \wedge b < 0\} \vee \{a < 0 \wedge b > 0\}$
  \item Si $ab > 0 \wedge a > b \rightarrow \dfrac{1}{a} < \dfrac{1}{b}$
  \item Si: $0 < x < a \rightarrow \dfrac{1}{x} > \dfrac{1}{a}$
\end{enumerate}
\end{minipage}

\begin{enumerate}[label=\cn{\arabic*}]\setcounter{enumi}{8}
  \item Si $ab < 0 \wedge a < x < b \rightarrow 0 \leq x^2 < \operatorname{Máx}\{a^2; b^2\}$
  \item $\forall x > 0:\ x + \dfrac{1}{x} \geq 2$
  \item $\forall x < 0:\ x + \dfrac{1}{x} \leq -2$
  \item Si $\{a; b; c\} \subset \mathbb{R}^+$, las medias cumplen:
  \ovalbox{$MA \geq MG \geq MH$}
  \[
  \bullet\ \frac{a+b}{2} \geq \sqrt{ab} \geq \frac{2}{\frac{1}{a} + \frac{1}{b}}
  \qquad\qquad
  \bullet\ \frac{a+b+c}{3} \geq \sqrt[3]{abc} \geq \frac{3}{\frac{1}{a} + \frac{1}{b} + \frac{1}{c}}
  \]
  \item \textcolor{magenta}{\textbf{Desigualdad de Cauchy--Schwarz--Buniakovski}}

  Sean $a_1, a_2, \ldots, a_n, b_1, b_2, \ldots, b_n$ números reales no todos nulos, entonces:
  \[\ovalbox{$(a_1b_1 + a_2b_2 + \cdots + a_nb_n)^2 \leq (a_1^2 + a_2^2 + \cdots + a_n^2)(b_1^2 + b_2^2 + \cdots + b_n^2)$}\]
\end{enumerate}


% ============================================================
% Páginas 43-45
% ============================================================
\newpage
\section{Desigualdades II}

\textbf{Inecuación:} Son aquellas desigualdades planteadas entre expresiones matemáticas
definidas sobre $\mathbb{R}$; los valores de las incógnitas que verifican la desigualdad
forman el conjunto solución (C.S.).

\subsection*{Inecuaciones polinomiales}

\textcolor{verde}{\textbf{Inecuación lineal:}} Es aquella inecuación polinomial que se reduce
a la forma general:
\[\ovalbox{$P(x) = ax + b \gtrless 0$}\;;\ a \neq 0\]
{\small\[
\begin{array}{cccc|l}
\cn{1}\ ax + b > 0 & \cn{2}\ ax + b < 0 & \cn{3}\ ax + b \geq 0 & \cn{4}\ ax + b \leq 0 & \\[4pt]
x > -\frac{b}{a} & x < -\frac{b}{a} & x \geq -\frac{b}{a} & x \leq -\frac{b}{a} & \text{Con } a \in \mathbb{R}^+\\[8pt]
x \in \left\langle -\frac{b}{a}; +\infty\right\rangle &
x \in \left\langle -\infty; -\frac{b}{a}\right\rangle &
x \in \left[-\frac{b}{a}; +\infty\right\rangle &
x \in \left\langle -\infty; -\frac{b}{a}\right] &
\end{array}
\]}

\textcolor{verde}{\textbf{Inecuación cuadrática:}} Son aquellas inecuaciones polinomiales de
la forma \ \fbox{$P(x) = ax^2 + bx + c \gtrless 0$}

Siendo: $a, b, c \in \mathbb{R};\ a \neq 0$.

Considerando $a > 0$, tenemos:
\[
a\left(x^2 + \frac{b}{a}x + \frac{c}{a}\right) \gtrless 0
\longrightarrow
a\left(\left(x + \frac{b}{2a}\right)^2 - \left(\frac{b^2 - 4ac}{4a^2}\right)\right) \gtrless 0
\longrightarrow
a\left(\left(x + \frac{b}{2a}\right)^2 - \frac{\Delta}{4a^2}\right) \gtrless 0
\]

\textcolor{rojo}{\textbf{Caso I:}} Si $\Delta = 0$, reemplazando nos damos cuenta que como
$a > 0 \rightarrow P(x) \geq 0\ \forall x \in \mathbb{R}$, pues:
\[\left(x + \frac{b}{2a}\right)^2 \geq 0 \;:\; \text{Siempre es positivo}\]
Ejemplo: \quad $x^2 - 6x + 9 < 0 \leftrightarrow (x-3)^2 < 0, \quad C.S. = \varnothing$

\textcolor{rojo}{\textbf{Caso II:}} Si $\Delta > 0$, tenemos:
\[
\left(x + \frac{b}{2a} + \frac{\sqrt{\Delta}}{2a}\right)\left(x + \frac{b}{2a} - \frac{\sqrt{\Delta}}{2a}\right) \gtrless 0
\]
y para resolverlo aplicaremos la \underline{regla de los signos} o de los
\underline{puntos de referencia}.

Ejemplo:

\begin{minipage}[t]{0.485\textwidth}
\centering
$x^2 - 5x + 4 \geq 0$\\
$(x-4)(x-1) \geq 0$\\[6pt]
\rectasignos{1/c,4/c}{+,-,+}{3}\\[6pt]
$C.S. = \langle -\infty; 1] \cup [4; +\infty\rangle$
\end{minipage}\hfill
\begin{minipage}[t]{0.485\textwidth}
\centering
$x^2 - 4x - 5 < 0$\\
$(x-5)(x+1) < 0$\\[6pt]
\rectasignos{-1/o,5/o}{+,-,+}{3}\\[6pt]
$C.S. = \langle -1; 5\rangle$
\end{minipage}

\medskip
\textcolor{rojo}{\textbf{Caso III:}} Si $\Delta < 0$, observamos que como
$a > 0 \rightarrow P(x) > 0\ \forall x \in \mathbb{R}$, pues:
\[\underbrace{\left(x + \frac{b}{2a}\right)^2}_{+}\underbrace{-\frac{\Delta}{4a^2}}_{+} > 0\]

\textcolor{naranja}{\textbf{Teorema (Trinomio positivo):}} El polinomio cuadrático
$P(x) = ax^2 + bx + c$ con coeficientes reales \underline{es positivo} para todo valor real
de $x \leftrightarrow a > 0 \wedge \Delta < 0$.

Ejemplo:
\begin{align*}
&x^2 + 2x + 3 > 0 \quad \text{como } \Delta < 0 \rightarrow C.S. = \mathbb{R}\\
&x^2 + 4x + 7 < 0 \quad \text{como } \Delta < 0 \rightarrow C.S. = \varnothing
\end{align*}

\textcolor{naranja}{\textbf{Teorema (Trinomio no negativo):}} El polinomio cuadrático
$P(x) = ax^2 + bx + c$ con coeficientes reales es \underline{no negativo} para todo valor
real de $x \leftrightarrow a > 0 \wedge \Delta \leq 0$.

\textcolor{verde}{\textbf{Inecuación polinómica de grado superior:}} Son aquellas que se
reducen a la forma:
\[\ovalbox{$P(x) = a_0x^n + a_1x^{n-1} + a_2x^{n-2} + \cdots + a_{n-1}x + a_n \gtrless 0$}\]
\[a_0; a_1; a_2; a_3; \ldots; a_n \in \mathbb{R} \wedge a_0 \neq 0 \qquad
\text{Generalmente: } a_0 > 0,\ n \in \mathbb{Z}^+;\ n \geq 2\]

Tenemos que factorizar el polinomio hasta obtener factores lineales o sus potencias y
factores de signo conocido, para luego aplicar la regla de los
\underline{puntos de referencia}.
\begin{itemize}
  \item Cuando se presentan factores de la forma $(x-\alpha)^{2n+1}$ se toman como si
  tuviesen exponente 1, es decir $(x - \alpha)$.
  \item $(x-\alpha)^{2n}$: el punto crítico no será colocado en la recta numérica; se
  analiza por separado, reemplazando en la inecuación para saber si forma parte o no del C.S.
  \item Si aparecen factores positivos o $\Delta < 0$, estos serán eliminados y no alteran
  la ecuación.
\end{itemize}

Ejemplo:

\begin{minipage}{0.436\textwidth}
\begin{gather*}
x^6 - x^4 - 2x^2 + 2 \leq 0\\
x^4(x^2 - 1) - 2(x^2 - 1) \leq 0\\
(x^4 - 2)(x^2 - 1) \leq 0\\
\underbrace{(x^2 + \sqrt{2})}_{+}(x^2 - \sqrt{2})(x - 1)(x + 1) \leq 0
\end{gather*}
\end{minipage}\hfill
\begin{minipage}{0.533\textwidth}
\centering
$(x - \sqrt[4]{2})(x + \sqrt[4]{2})(x - 1)(x + 1) \leq 0$\\[6pt]
\rectasignos{-\sqrt[4]{2}/c,-1/c,1/c,\sqrt[4]{2}/c}{+,-,+,-,+}{5}\\[6pt]
$\therefore\ C.S. = [-\sqrt[4]{2}; -1] \cup [1; \sqrt[4]{2}]$
\end{minipage}

\newpage
% ------------------------------------------------------------
\section{Desigualdades III}

\subsection*{Inecuaciones racionales o fraccionarias}
\[
\frac{P(x)}{F(x)} \gtrless 0 \qquad
\begin{array}{l}
\text{donde: } \textcolor{magenta}{P(x), F(x)}\text{: son polinomios}\\
\textcolor{magenta}{F(x)}\text{: al menos de primer grado.}
\end{array}
\]
Para resolver las inecuaciones polinómicas y fraccionarias utilizaremos el método de los
puntos de referencia.

En el caso de una inecuación polinómica lineal tan sólo despejamos la incógnita
considerando los teoremas sobre desigualdades ya estudiados. Los puntos de referencia son
las raíces del polinomio.

\textbf{Procedimiento:}
\begin{enumerate}
  \item Transponer la expresión a un solo miembro, luego factorizar tratando de obtener
  factores lineales y/o cuadráticos con coeficientes principales positivos.
  \item Igualamos cada factor a cero para encontrar los puntos de referencia, estos se
  disponen en la recta real y se analizan (abiertos o cerrados).
  \item Colocamos los signos $+$ y $-$ en las zonas, de derecha a izquierda. Cuando un
  factor lineal tiene exponente par la zona positiva se refleja.
  \item Si la expresión es \underline{mayor que cero} se tomarán \underline{las zonas
  positivas} y si es \underline{menor que cero} se tomarán las \underline{zonas negativas}.
  \item En la inecuación fraccionaria si tenemos $\geq;\ \leq$, los puntos de referencia son
  cerrados en el numerador y siempre abiertos en el denominador.
\end{enumerate}

\textbf{Ejemplos:}

\begin{minipage}[t]{0.466\textwidth}
\begin{gather*}
\frac{\overset{+}{2006^7}(x-2)^7(1-x)(3x-1)^8}{\underset{+}{2005^7}} \leq 0\\
(x-2)^7(1-x)(3x-1)^8 \leq 0\\
(x-2)^7(x-1)(3x-1)^8 \geq 0
\end{gather*}
\centering
\rectasignos{\tfrac{1}{3}/c,1/c,2/c}{+,+,-,+}{4}
\begin{gather*}
C.S. = \langle -\infty; 1] \cup [2; +\infty\rangle\\
C.S. = \mathbb{R} - \langle 1; 2\rangle
\end{gather*}
\end{minipage}\hfill
\begin{minipage}[t]{0.485\textwidth}
\[\frac{(2x-3)^{31}(5x-9)^{15}}{(4x-7)^{20}(3x-4)^{42}} \geq 0\]
\centering
\rectasignos{\tfrac{4}{3}/o,\tfrac{3}{2}/c,\tfrac{7}{4}/o,\tfrac{9}{5}/c}{+,+,-,-,+}{5}
\begin{gather*}
C.S. = \left\langle -\infty; \tfrac{4}{3}\right\rangle \cup \left\langle \tfrac{4}{3}; \tfrac{3}{2}\right] \cup \left[\tfrac{9}{5}; +\infty\right\rangle\\
C.S. = \mathbb{R} - \left\{\tfrac{4}{3}\right\} - \left\langle \tfrac{3}{2}; \tfrac{9}{5}\right\rangle
\end{gather*}
\end{minipage}


% ============================================================
% Páginas 46-48
% ============================================================
\subsection*{Inecuaciones irracionales}

\textbf{I. Forma:} \ \ovalbox{$\sqrt[2n]{A} > B$}$,\ n \in \mathbb{Z}^+$
\[
\left.
\begin{array}{l}
S_1: A \geq 0 \wedge B \geq 0 \wedge A > B^{2n}\\
S_2: A \geq 0 \wedge B < 0
\end{array}
\right\}\; C.S. = S_1 \cup S_2
\]

\textbf{II. Forma:} \ \ovalbox{$\sqrt[2n]{A} < B$}$;\ n \in \mathbb{Z}^+$
\[C.S. = A \geq 0 \wedge B > 0 \wedge A < B^{2n}\]

\textbf{III. Forma:} \ \ovalbox{$\sqrt[2n]{A} \gtrless \sqrt[2n]{B}$}$;\ m \wedge n \in \mathbb{Z}^+$
\[C.S. = A \geq 0 \wedge B \geq 0 \wedge A^n \gtrless B^n\]

\subsection*{Ecuaciones irracionales}
\begin{enumerate}[label=\Roman*.]
  \item $\sqrt[2n]{f(x)} = g(x) \longleftrightarrow f(x) \geq 0 \wedge g(x) \geq 0 \wedge f(x) = [g(x)]^{2n}$
  \item $\sqrt[2n]{f(x)} = \sqrt[2n]{g(x)} \longleftrightarrow f(x) \geq 0 \wedge g(x) \geq 0 \wedge f(x) = g(x)$
  \item $\sqrt{f_1(x)} + \sqrt{f_2(x)} + \sqrt{f_3(x)} + \cdots + \sqrt{f_n(x)} = 0$
  \[\longleftrightarrow f_1(x) = f_2(x) = f_3(x) = \cdots = f_n(x) = 0\]
\end{enumerate}

\subsection*{Existencia de una expresión irracional sobre $\mathbb{R}$}
Si $n \in \mathbb{Z}^+$:
\[
\ovalbox{$\sqrt[2n]{f(x)} \in \mathbb{R} \longleftrightarrow f(x) \geq 0$}
\qquad
\ovalbox{$\sqrt[2n+1]{g(x)} \in \mathbb{R} \longleftrightarrow g(x) \in \mathbb{R}$}
\]

\newpage
% ------------------------------------------------------------
\section{Valor absoluto}

\textbf{Definición:}
\[
|x| = \begin{cases}
x; & \text{si } x \geq 0\\
\operatorname{op}(x) = -x; & \text{si } x < 0
\end{cases}
\]

\subsection*{Teoremas}
\begin{minipage}[t]{0.320\textwidth}
\begin{enumerate}[label=\cn{\arabic*}]
  \item $|x| = 0 \longleftrightarrow x = 0$
  \item $|x| \geq 0$
  \item $|x| = |-x|$
  \item $|x| \geq x,\ |x| \geq -x$
\end{enumerate}
\end{minipage}\hfill
\begin{minipage}[t]{0.291\textwidth}
\begin{enumerate}[label=\cn{\arabic*}]\setcounter{enumi}{5}
  \item $|xy| = |x||y|$
  \item $\left|\dfrac{x}{y}\right| = \dfrac{|x|}{|y|}$
  \item $|x^2| = |x|^2 = x^2$
  \item $\sqrt[2n]{x^{2n}} = |x|$
\end{enumerate}
\end{minipage}\hfill
\begin{minipage}[t]{0.359\textwidth}
\begin{enumerate}[label=\cn{\arabic*}]\setcounter{enumi}{9}
  \item Si $a < x < b \wedge ab < 0$ entonces:
  \[0 \leq |x| < \operatorname{Máx}\{|a|, |b|\}\]
\end{enumerate}
\end{minipage}
\begin{enumerate}[label=\cn{\arabic*}]\setcounter{enumi}{4}
  \item $|x + y| \leq |x| + |y|$ \ \textcolor{magenta}{(Desigualdad triangular)}
\end{enumerate}

\textbf{Ecuaciones con valor absoluto:} Forma general \fbox{$F(x) = 0$}. Donde $F(x)$ es una
expresión con valor absoluto.

\textcolor{naranja}{\textbf{Teoremas:}} Sean $f(x)$ y $g(x)$ expresiones definidas en $\mathbb{R}$.
\begin{enumerate}[label=\cn{\arabic*}]
  \item $|f(x)| = g(x) \longleftrightarrow g(x) \geq 0 \wedge [f(x) = g(x) \vee f(x) = -g(x)]$
  \item $|f(x)| = |g(x)| \longleftrightarrow [f(x) = g(x) \vee f(x) = -g(x)]$
  \item $|f(x) + g(x)| = |f(x)| + |g(x)| \longleftrightarrow f(x)\cdot g(x) \geq 0$
\end{enumerate}

\textbf{Ejemplos:}

\begin{minipage}[t]{0.504\textwidth}
\[|2x - 4| = 5x - x^2 - 6\]
\begin{minipage}[t]{0.5\linewidth}
\begin{gather*}
-x^2 - 6 + 5x \geq 0\\
x^2 - 5x + 6 \leq 0\\
(x-3)(x-2) \leq 0\\
x \in [2; 3]
\end{gather*}
\end{minipage}\hfill
\begin{minipage}[t]{0.5\linewidth}
\begin{gather*}
2x - 4 = 5x - x^2 - 6\\
x^2 - 3x + 2 = 0\\
(x-2)(x-1) = 0\\
x \in \{1; 2\}
\end{gather*}
\end{minipage}
\begin{gather*}
2x - 4 = -5x + x^2 + 6\\
x^2 - 7x + 10 = 0\\
(x-5)(x-2) = 0\\
x \in \{2; 5\}\\
\therefore\ x = 2
\end{gather*}
\end{minipage}\hfill
\begin{minipage}[t]{0.466\textwidth}
\[|x^2 + 3| = |2x + 2|\]
\begin{minipage}[t]{0.5\linewidth}
\begin{gather*}
x^2 + 3 = 2x + 2\\
x^2 - 2x + 1 = 0\\
x = 1
\end{gather*}
\end{minipage}\hfill
\begin{minipage}[t]{0.5\linewidth}
\begin{gather*}
x^2 + 3 = -2x - 2\\
x^2 + 2x + 5 = 0\\
\Delta < 0\\
x \in \varnothing
\end{gather*}
\end{minipage}
\[\therefore\ x = 1\]
\end{minipage}

\subsection*{Inecuaciones con valor absoluto}
Sean $f(x)$ y $g(x)$ expresiones definidas en $\mathbb{R}$.
\begin{enumerate}[label=\cn{\arabic*}]
  \item $|f(x)| \leq g(x) \longleftrightarrow g(x) \geq 0 \wedge [-g(x) \leq f(x) \leq g(x)]$
  \item $|f(x)| \geq g(x) \longleftrightarrow [f(x) \geq g(x) \vee f(x) \leq -g(x)]$
  \item $|f(x)| \lesseqgtr |g(x)| \longleftrightarrow (f(x))^2 \lesseqgtr (g(x))^2$
  \item $|f(x) + g(x)| < |f(x)| + |g(x)| \longleftrightarrow f(x)\cdot g(x) < 0$
\end{enumerate}

\textbf{Ejemplos:}

\begin{minipage}[t]{0.291\textwidth}
\begin{gather*}
\left|x + \frac{1}{x}\right| \geq 2\\
x + \frac{1}{x} \geq 2 \;\vee\; x + \frac{1}{x} \leq -2\\
x \in \mathbb{R}^+ - \{0\} \;\vee\; x \in \mathbb{R}^- - \{0\}\\
\therefore\ x \in \mathbb{R} - \{0\}
\end{gather*}
\end{minipage}\hfill
\begin{minipage}[t]{0.330\textwidth}
\begin{gather*}
|x^2 + 3| > |x + 3|\\
(x^2 + 3)^2 > (x + 3)^2\\
(x^2 + 3)^2 - (x + 3)^2 > 0\\
(x^2 - x)\underbrace{(x^2 + x + 6)}_{+} > 0\\
x(x - 1) > 0\\
\therefore\ x \in \mathbb{R} - [0; 1]
\end{gather*}
\end{minipage}\hfill
\begin{minipage}[t]{0.349\textwidth}\small
\begin{gather*}
\left|\frac{x^2}{x - 2}\right| \leq x - 4\\
\text{C.V.A.} \rightarrow x \geq 4;\ x \neq 2\\
\frac{x^2}{x-2} \leq x - 4 \;\wedge\; 4 - x \leq \frac{x^2}{x-2}\\
x^2 \leq x^2 - 6x + 8 \;\wedge\; -x^2 + 6x - 8 \leq x^2\\
x \leq \frac{4}{3} \;\wedge\; x^2 - 3x + 4 \geq 0\\
\Delta < 0 \rightarrow x \in \mathbb{R}\\
x \leq \frac{4}{3} \;\wedge\; x \geq 4\\
\therefore\ C.S. = \varnothing
\end{gather*}
\end{minipage}


\medskip
Hay otros casos donde no es tan simple aplicar y para estos hay otro método (por zonas):

\begin{minipage}[t]{0.466\textwidth}
$\bullet\ |x - 2| + 3|x - 1| < 4 \rightarrow C.S. = S_1 \cup S_2 \cup S_3 = \left\langle \frac{1}{4}; \frac{9}{4}\right\rangle$

\begin{center}
\begin{tikzpicture}[x=1.2cm]
  \draw[<->] (0,0) -- (4,0);
  \foreach \p/\l in {1.3/1, 2.6/2} {\draw (\p,0.1) -- (\p,-0.1); \node[below] at (\p,-0.1) {\small $\l$};}
  \node[below,naranja] at (0.65,-0.1) {\small I};
  \node[below,naranja] at (1.95,-0.1) {\small II};
  \node[below,naranja] at (3.3,-0.1) {\small III};
  \node[above,magenta] at (0.65,0.05) {\scriptsize $\langle -\infty;1\rangle$};
  \node[above,magenta] at (1.95,0.05) {\scriptsize $[1;2\rangle$};
  \node[above,magenta] at (3.3,0.05) {\scriptsize $[2;+\infty\rangle$};
\end{tikzpicture}
\end{center}

\textcolor{naranja}{En I:}
\begin{gather*}
2 - x + 3(1 - x) < 4\\
2 - x + 3 - 3x < 4\\
x > 1/4 \;\cap\; x < 1\\
S_1: x \in \langle 1/4; 1\rangle
\end{gather*}
\textcolor{naranja}{En II:}
\begin{gather*}
2 - x + 3(x - 1) < 4\\
2 - x + 3x - 3 < 4\\
x < 5/2 \;\cap\; 1 \leq x < 2\\
S_2: x \in [1; 2\rangle
\end{gather*}
\textcolor{naranja}{En III:}
\begin{gather*}
x - 2 + 3(x - 1) < 4\\
x < 9/4 \;\cap\; x \geq 2\\
S_3: x \in [2; 9/4\rangle
\end{gather*}
\end{minipage}\hfill
\begin{minipage}[t]{0.485\textwidth}
$\bullet\ |x + 6| > |x + 9| + |x + 2| \rightarrow C.S. = S_1 \cup S_2 \cup S_3 \cup S_4 = \varnothing$

\begin{center}
\begin{tikzpicture}[x=1.1cm]
  \draw[<->] (0,0) -- (5,0);
  \foreach \p/\l in {1.2/-9, 2.4/-6, 3.6/-2} {\draw (\p,0.1) -- (\p,-0.1); \node[below] at (\p,-0.1) {\small $\l$};}
  \node[below,naranja] at (0.6,-0.1) {\small I};
  \node[below,naranja] at (1.8,-0.1) {\small II};
  \node[below,naranja] at (3.0,-0.1) {\small III};
  \node[below,naranja] at (4.3,-0.1) {\small IV};
  \node[above,magenta] at (0.6,0.05) {\scriptsize $\langle -\infty;-9\rangle$};
  \node[above,magenta] at (1.8,0.05) {\scriptsize $[-9;-6\rangle$};
  \node[above,magenta] at (3.0,0.05) {\scriptsize $[-6;-2\rangle$};
  \node[above,magenta] at (4.3,0.05) {\scriptsize $[-2;+\infty\rangle$};
\end{tikzpicture}
\end{center}

\begin{minipage}[t]{0.5\linewidth}
\textcolor{naranja}{En I:}
\begin{gather*}
-x - 6 > -x - 9 - x - 2\\
x > -5 \;\cap\; x < -9\\
S_1: x = \varnothing
\end{gather*}
\textcolor{naranja}{En III:}
\begin{gather*}
x + 6 > x + 9 - x - 2\\
x > 1 \;\cap\; -6 \leq x < -2\\
S_3: x = \varnothing
\end{gather*}
\end{minipage}\hfill
\begin{minipage}[t]{0.5\linewidth}
\textcolor{naranja}{En II:}
\begin{gather*}
-x - 6 > x + 9 - x - 2\\
x < -13 \;\cap\; -9 \leq x < -6\\
S_2: x = \varnothing
\end{gather*}
\textcolor{naranja}{En IV:}
\begin{gather*}
x + 6 > x + 9 + x + 2\\
x < -5 \;\cap\; x \geq -2\\
S_4: x = \varnothing
\end{gather*}
\end{minipage}

\medskip
\textcolor{rojo}{$\ast$} También se utiliza en ecuaciones con valor absoluto.
\end{minipage}


% ============================================================
% Páginas 49-51
% ============================================================
\newpage
\section{Relaciones}

\subsection*{Conceptos previos}

\textcolor{magenta}{\textbf{Pares ordenados:}} Son entes matemáticos que consisten de dos
elementos ``$a$'' y ``$b$'', a los cuales se les denomina primera componente y segunda
componente respectivamente, y se les denota por $(a;b)$.

\textbf{Igualdad de pares ordenados:} Dados dos pares ordenados $(a;b)$ y $(c;d)$, entonces:
\[\ovalbox{$(c;d) = (a;b) \longleftrightarrow a = c \text{ y } b = d$}\]
Ejemplo:
\begin{itemize}
  \item $(2;3) \neq (3;2)$
  \item $(2x + y; 1) = (3; 2x - y) \longleftrightarrow
  \begin{array}{l} 2x + y = 3\\ 1 = 2x - y\end{array}
  \longleftrightarrow
  \begin{array}{l} x = 1\\ y = 1\end{array}$
\end{itemize}

\textcolor{magenta}{\textbf{Producto cartesiano:}} Dados dos conjuntos no vacíos $A$ y $B$ se
define el producto cartesiano $A\times B$ como el conjunto de pares ordenados.
\[\ovalbox{$A\times B = \{(a;b) \,/\, a \in A \wedge b \in B\}$}\]
Tales que su 1\textsuperscript{ra} componente está en $A$ y su 2\textsuperscript{da} componente
está en $B$.

Ejemplo: \quad Sea $A = \{1; 2; 3\}$ y $B = \{a; b\}$, entonces:
\[A\times B = \{(1;a); (1;b); (2;a); (2;b); (3;a); (3;b)\}\]

\textcolor{magenta}{\textbf{Nota:}} En general, si los conjuntos $A$ y $B$ son finitos con
$m$ y $n$ elementos respectivamente, entonces el producto cartesiano $A\times B$ tiene
$m\times n$ elementos. De aquí su nombre y notación. En general, el producto cartesiano no
es conmutativo; es decir $A\times B \neq B\times A$, a menos que $A = B$.

\textbf{Relaciones:} Dados dos conjuntos no vacíos $A$ y $B$, a un conjunto $R$ de pares
ordenados se le denomina ``Relación de $A$ en $B$'' si $R$ es un subconjunto cualquiera de
$A\times B$. También se le llama ``Relación Binaria''.
\[\ovalbox{$R$ es una relación de $A$ en $B$ si y sólo si $R \subset A\times B$}\]
Puesto que, en general, si $A\times B$ tiene $n$ elementos entonces $A\times B$ tiene $2^n$
subconjuntos, entonces existen $2^n$ relaciones de $A$ en $B$.

Cuando un par ordenado $(a;b)$ pertenece a una relación $R$, también se denota: $a\,R\,b$,
es decir, $a\,R\,b$ si y sólo si $(a;b) \in R$ y en tal caso se lee: ``$a$ está relacionado
con $b$ según la relación $R$'' y si $(a;b) \notin R$ entonces se denota $a \not\!R\, b$.

\textbf{Definición:} Se dice que $R$ es una relación en el conjunto $A$ si
$R \subset A\times A$ ó $R \subset A^2$.

Ejemplo: Si $R$ es una relación en $A$, $A = \{2; 3; 4\}$ tal que:
\[R = \{(x;y) \,/\, y + 1 < x^2\}\]
\[R = \{(2;2); (2;3); (3;2); (3;3); (3;4); (4;2); (4;3); (4;4)\}\]
pues para $(x;y) \in A\times A$, con $x \in A$ y $y \in A$:
\begin{align*}
x = 2:\ y + 1 < 2^2 &\longrightarrow y \in \{2; 3\} \longrightarrow (2;2), (2;3) \in R\\
x = 3:\ y + 1 < 3^2 &\longrightarrow y \in \{2; 3; 4\} \longrightarrow (3;2), (3;3), (3;4) \in R\\
x = 4:\ y + 1 < 4^2 &\longrightarrow y \in \{2; 3; 4\} \longrightarrow (4;2), (4;3), (4;4) \in R
\end{align*}

\subsection*{Tipos de relaciones}

\textcolor{magenta}{\textbf{Relaciones reflexivas:}} Una relación $R$ es reflexiva en $A$,
$[R \subset A^2]$ si $\forall a \in A,\ (a;a) \in R$. Es decir, $R$ es reflexiva en $A$ si todo
elemento de $A$ está relacionado consigo mismo mediante la relación $R$.

Ejemplo: Sean $A = \{1; 2; 3; 4\}$, y las relaciones en $A$:
\[
R_1 = \{(1;2); \underline{(3;3)}; (3;4); \underline{(4;4)}; (4;1); \underline{(2;2)}; \underline{(1;1)}\}
\qquad
R_2 = \{\underline{(1;1)}; \underline{(2;2)}; (3;4); (4;3); \underline{(4;4)}\}
\]
Entonces $R_1$ es reflexiva pues $(a;a) \in R_1,\ \forall a \in A$, además de otros puntos; en
cambio a $R_2$ le falta $(3;3)$ para serlo.

\textcolor{magenta}{\textbf{Relaciones simétricas:}} Una relación $R$ en un conjunto $A$ es
simétrica si se cumple la siguiente implicación:
\[\ovalbox{$(a;b) \in R \longrightarrow (b;a) \in R$}\]
Es decir, si $(a;b)$ está en $R$ entonces el elemento $(b;a)$ también debe estar en $R$ para
que $R$ sea simétrica.

Ejemplo: Sea $A = \{1; 2; 3; 4\}$, y las relaciones en $A$:
\begin{gather*}
R_1 = \{(1;2); (2;3); (4;2); (3;2); (2;1); (2;4)\} \qquad R_2 = \{(1;1); (2;2); (3;3)\}\\
R_3 = \{(1;1); (3;3); (4;1); (2;3); (1;4)\}
\end{gather*}
Entonces $R_1$ y $R_2$ son simétricas pero $R_3$ no lo es pues le falta el elemento $(3;2)$
para serlo.

\textcolor{magenta}{\textbf{Relaciones transitivas:}} Una relación $R$ en un conjunto $A$ es
transitiva si se cumple la siguiente implicación:
\[\ovalbox{$[(a;b) \in R \wedge (b;c) \in R] \longrightarrow (a;c) \in R$}\]
Ejemplo: Sea $A = \{1; 2; 3; 4\}$, la relación en $A$:
\[R_1 = \{(1;2); (2;3); (1;3); (3;1); (1;1)\}\]
No es transitiva, pues si bien se cumplen las implicaciones:
\begin{align*}
(1;2) \in R_1 \wedge (2;3) \in R_1 &\longrightarrow (1;3) \in R_1\\
(1;3) \in R_1 \wedge (3;1) \in R_1 &\longrightarrow (1;1) \in R_1
\end{align*}
en cambio falla en:
\[(2;3) \in R_1 \wedge (3;1) \in R_1 \longrightarrow (2;1) \notin R_1\]
\[R_2 = \{(1;2); (2;1); (2;2); (1;1)\} \quad \text{sí es transitiva.}\]

\textcolor{magenta}{\textbf{Relaciones de equivalencia:}} Una relación $R$ en $A$ es una
relación de equivalencia si satisface simultáneamente las tres condiciones:
\begin{itemize}[label=\textcolor{rojo}{\checkmark}]
  \item \textcolor{rojo}{Reflexiva:} $\forall a \in A,\ (a;a) \in R$.
  \item \textcolor{rojo}{Simétrica:} Si $(a;b) \in R$ entonces $(b;a) \in R$.
  \item \textcolor{rojo}{Transitiva:} Si $(a;b) \in R$ y $(b;c) \in R$ entonces $(a;c) \in R$.
\end{itemize}
Ejemplos: Sea $A = \{1; 2; 3; 4\}$ entonces la relación
$R = \{(1;1); (2;1); (1;2); (2;2); (3;3); (4;4)\}$ es una relación de equivalencia en $A$.

\textbf{Dominio y rango de una relación:} Se llama dominio al conjunto de todas las primeras
componentes de los pares ordenados de $R$, se llama rango al conjunto de todas las segundas
componentes.


% ============================================================
% Páginas 52-54
% ============================================================
\[
\ovalbox{$\operatorname{Dom}(R) = \{x \,/\, (x;y) \in R\}$}
\qquad
\ovalbox{$\operatorname{Ran}(R) = \{y \,/\, (x;y) \in R\}$}
\]
Ejemplo: Dada la relación en $A = \{1; 2; 3; 4; 5\}$
\[R = \{(1;1); (2;1); (2;2); (3;1); (3;2); (4;2); (5;3)\}\]
\[\operatorname{Dom}(R) = \{1; 2; 3; 4; 5\} \qquad \operatorname{Ran}(R) = \{1; 2; 3\}\]

\textbf{Relación inversa:} Dado un conjunto no vacío $A$ y la relación $R: A \to A$ tal que:
\[\ovalbox{$R = \{(x;y) \in A^2 \,/\, P_{(x;y)}\}$}\]
Se define la relación INVERSA de $A$ como:
\[\ovalbox{$R^* = R^{-1} = \{(y;x) \in A^2 \,/\, P_{(x;y)}\}$}\]
donde: \ \ovalbox{$\operatorname{Dom}R^* = \operatorname{Ran}R \;\wedge\; \operatorname{Ran}R^* = \operatorname{Dom}R$}

\textcolor{magenta}{\textbf{Propiedad:}} Dada la relación $R: A \to A$, donde $A$ es un
conjunto no vacío:
\[\ovalbox{$R$ es simétrica $\longleftrightarrow R = R^*$}\]

\newpage
% ------------------------------------------------------------
\section{Funciones I}

\textbf{Definición:} Dados dos conjuntos no vacíos $A$ y $B$ se define una función de $A$ en
$B$ como un conjunto de pares ordenados $(x;y)$ tal que a cada elemento $x \in A$ existe un
único elemento $y \in B$.

\textcolor{magenta}{\textbf{Notación:}} Si $f$ es una función de $A$ en $B$, luego:
\ \ovalbox{$f: A \to B$}

\textcolor{magenta}{\textbf{Condición de existencia y unicidad:}} Sea $f: A \to B$, se debe
cumplir:

\begin{minipage}[t]{0.485\textwidth}
\begin{enumerate}
  \item $\forall x \in A;\ \exists!\, y \in B \,/\, (x;y) \in f$

  \textcolor{magenta}{$\exists!$}: Se lee ``Existe y es único''
\end{enumerate}
\end{minipage}\hfill
\begin{minipage}[t]{0.485\textwidth}
\begin{enumerate}\setcounter{enumi}{1}
  \item Si $(x;y) \in f \wedge (x;z) \in f \rightarrow y = z$
\end{enumerate}
\end{minipage}

\smallskip
De la definición de una función se deduce que dos pares ordenados diferentes no pueden tener
la misma primera componente.

\textbf{Ejemplos:}

\begin{center}
\begin{tikzpicture}[>=Stealth, every node/.style={font=\small}]
  % f1
  \begin{scope}
    \draw (0,0) ellipse (0.7 and 1.3); \draw (2.2,0) ellipse (0.7 and 1.3);
    \node at (0,1.6) {$A$}; \node at (2.2,1.6) {$B$};
    \draw[->,verde,thick] (0.2,1.6) to[bend left=20] node[above] {$f_1$} (2.0,1.6);
    \node (a1) at (0,0.7) {1}; \node (a2) at (0,0) {2}; \node (a3) at (0,-0.7) {3};
    \node (b1) at (2.2,0.35) {4}; \node (b2) at (2.2,-0.35) {5};
    \draw[->] (a1) -- (b1); \draw[->] (a2) -- (b2); \draw[->] (a3) -- (b1);
    \node at (1.1,-1.9) {\underline{Son funciones}};
  \end{scope}
  % f2
  \begin{scope}[xshift=4.2cm]
    \draw (0,0) ellipse (0.7 and 1.3); \draw (2.2,0) ellipse (0.7 and 1.3);
    \node at (0,1.6) {$A$}; \node at (2.2,1.6) {$B$};
    \draw[->,verde,thick] (0.2,1.6) to[bend left=20] node[above] {$f_2$} (2.0,1.6);
    \node (a1) at (0,0.9) {1}; \node (a2) at (0,0.3) {2}; \node (a3) at (0,-0.3) {3}; \node (a4) at (0,-0.9) {4};
    \node (b1) at (2.2,0) {5};
    \foreach \a in {a1,a2,a3,a4} \draw[->] (\a) -- (b1);
  \end{scope}
  % f3
  \begin{scope}[xshift=8.4cm]
    \draw (0,0) ellipse (0.7 and 1.3); \draw (2.2,0) ellipse (0.7 and 1.3);
    \node at (0,1.6) {$A$}; \node at (2.2,1.6) {$B$};
    \draw[->,verde,thick] (0.2,1.6) to[bend left=20] node[above] {$f_3$} (2.0,1.6);
    \node (a1) at (0,0.5) {2}; \node (a2) at (0,-0.5) {3};
    \node (b1) at (2.2,0.7) {8}; \node (b2) at (2.2,0) {7}; \node (b3) at (2.2,-0.7) {6};
    \draw[->] (a1) -- (b1); \draw[->] (a1) -- (b2); \draw[->] (a2) -- (b3);
    \node at (1.1,-1.9) {\underline{No es función}};
  \end{scope}
\end{tikzpicture}
\end{center}

\textbf{Dominio y rango de una función:} El dominio es el conjunto de todas las primeras
componentes de los pares ordenados de $f$, se llama rango al conjunto de todas las segundas
componentes.
\[
\ovalbox{$\operatorname{Dom}(f) = \{x \in A \,/\, \exists!\, y \in B \wedge (x;y) \in f\} \subset A$}
\quad
\ovalbox{$\operatorname{Ran}(f) = \{y \in B \,/\, \exists x \in A \wedge (x;y) \in f\} \subset B$}
\]
Ejemplo: Sea $f = \{(1;2); (3;4); (5;6); (7;8)\}$, determinar el dominio y el rango.
\[D_f = \{1; 3; 5; 7\} \qquad R_f = \{2; 4; 6; 8\}\]

\textbf{Regla de correspondencia:} Es aquella relación que se establece entre la primera y
la segunda componente de una función. Esta relación se establece mediante una fórmula
matemática.

\begin{minipage}{0.388\textwidth}
\centering
\begin{tikzpicture}[>=Stealth]
  \draw (0,0) ellipse (0.8 and 1.3); \draw (2.6,0) ellipse (0.8 and 1.3);
  \node at (0,1.6) {$A$}; \node at (2.6,1.6) {$B$};
  \draw[->,verde,thick] (0.2,1.6) to[bend left=20] node[above] {$f: A \to B$} (2.4,1.6);
  \draw[->] (0,0) node[left] {$x$} -- (2.6,0) node[right] {$y$};
  \node at (1.3,-1.8) {\ovalbox{$y = f(x)$}};
\end{tikzpicture}
\end{minipage}\hfill
\begin{minipage}{0.582\textwidth}
\begin{itemize}
  \item Toda función queda completamente definida si se conoce el dominio y la regla de
  correspondencia.
  \item Toda función es una relación pero no toda relación es función.
\end{itemize}
\end{minipage}

\medskip
\textbf{Función real de variable real:} Sea $f: A \to B$, si
\ovalbox{$A \subset \mathbb{R} \wedge B \subset \mathbb{R}$} diremos que $f$ es una función
real de variable real.

\textcolor{magenta}{\textbf{Teorema fundamental de las funciones reales:}} Una función real
es de variable real, si toda recta vertical corta a su gráfica en un solo punto.

\begin{center}
\begin{tikzpicture}[scale=0.8]
  \begin{scope}
    \draw[->] (-2.5,0) -- (3,0) node[below] {$x$};
    \draw[->] (0,-2) -- (0,3) node[left] {$y$};
    \node[below left] at (0,0) {$O$};
    \draw[thick,rojo!70!black] (-2.3,-1.8) .. controls (-2,-0.5) and (-1,0.5) .. (0,0.8)
      .. controls (1,1) and (1.2,2.2) .. (2.6,2.6);
    \node at (-2.2,-0.8) {$f_1$};
    \draw[thick,rojo] (1.5,-2) -- (1.5,3);
    \fill (1.5,1.75) circle (2pt);
    \node at (0.8,-1.6) {\small $f_1$ es función};
  \end{scope}
  \begin{scope}[xshift=7.5cm]
    \draw[->] (-2.5,0) -- (3,0) node[below] {$x$};
    \draw[->] (0,-2) -- (0,3) node[left] {$y$};
    \node[below left] at (0,0) {$O$};
    \draw[thick] (2.8,2.5) -- (0.4,0.6) -- (2.8,-1.4);
    \node at (2.6,1.5) {$f_2$};
    \draw[thick,rojo] (1.5,-2) -- (1.5,3);
    \fill (1.5,1.47) circle (2pt); \fill (1.5,-0.27) circle (2pt);
    \node at (-0.6,-1.6) {\small $f_2$ no es función};
  \end{scope}
\end{tikzpicture}
\end{center}

\begin{itemize}
  \item El dominio se obtiene planteando la existencia de la regla de correspondencia
  sobre $\mathbb{R}$.
  \item No existe una regla específica para el cálculo del rango, sin embargo se recomienda
  despejar $x$ en función de $y$, para luego analizar para qué valores de $y$ la función
  está definida.
\end{itemize}

Ejemplo: Halle el rango de: \ $f(x) = \dfrac{3x - 1}{x - 2} = y$
\begin{enumerate}[label=\roman*)]
  \item $x - 2 \neq 0 \rightarrow x \neq 2$; luego: $\operatorname{Dom}f = \mathbb{R} - \{2\}$
  \item $yx - 2y = 3x - 1$
  \[x = \frac{2y - 1}{y - 3}; \text{ luego: } y - 3 \neq 0 \rightarrow y \neq 3, \text{ luego: } \operatorname{Ran}f = \mathbb{R} - \{3\}\]
\end{enumerate}


% ============================================================
% Páginas 55-57
% ============================================================
\newpage
\section{Funciones II}

\textbf{Gráfica de una función:} Es el conjunto de pares ordenados $(x;y)$ en donde ``$x$''
está en el dominio de la función y además ``$y = f(x)$''. Es decir:
\[\ovalbox{$\operatorname{Graf}(f) = \{(x;y) \in \mathbb{R}^2 \,/\, y = f(x);\ x \in \operatorname{Dom}(f)\}$}\]

\textbf{Funciones elementales:} Sirven de base para estudiar otras funciones.

\begin{minipage}[t]{0.466\textwidth}
\textbf{Función constante:}
\begin{center}
\begin{tikzpicture}[scale=0.7]
  \ejes{-0.5}{4}{-0.3}{3}
  \draw[<->,thick,graf] (-0.5,1.8) -- (4,1.8);
  \node[above] at (0.3,1.8) {$c$};
  \node at (2.8,2.6) {\ovalbox{$f(x) = c$}};
\end{tikzpicture}

\ovalbox{$\operatorname{Dom}f = \mathbb{R}$} \ \ovalbox{$\operatorname{Ran}f = \{c\}$}
\end{center}
\end{minipage}\hfill
\begin{minipage}[t]{0.485\textwidth}
\textbf{Función escalón unitario:} Su regla de correspondencia es:
\[
\ovalbox{$f(x) = \mu_a(x) = \begin{cases} 0; & x < a\\ 1; & x \geq a\end{cases}\quad a \in \mathbb{R}$}
\]
\begin{center}
\begin{tikzpicture}[scale=0.7]
  \ejes{-1}{4.5}{-0.5}{1.8}
  \draw[<-,thick,graf] (-1,0) -- (1.5,0); \draw[graf,fill=white] (1.5,0) circle (2.5pt);
  \draw[->,thick,graf] (1.5,1) -- (4.3,1); \fill[graf] (1.5,1) circle (2.5pt);
  \node[below] at (1.5,-0.1) {$a$}; \node[left] at (0,1) {$1$};
\end{tikzpicture}

\ovalbox{$\operatorname{Dom}f = \mathbb{R}$} \ \ovalbox{$\operatorname{Ran}f = \{0; 1\}$}
\end{center}
\end{minipage}

\medskip
\textbf{Función signo (Sgn):}

\begin{minipage}{0.436\textwidth}
\[
\ovalbox{$f(x) = \operatorname{Sgn}(x) = \begin{cases} -1; & x < 0\\ 0; & x = 0\\ 1; & x > 0\end{cases}$}
\]
\end{minipage}\hfill
\begin{minipage}{0.533\textwidth}
\centering
\begin{tikzpicture}[scale=0.7]
  \ejes{-3}{3}{-1.5}{1.6}
  \draw[thick,graf] (0,1) -- (2.8,1); \draw[graf,fill=white] (0,1) circle (2.5pt);
  \draw[thick,graf] (-3,-1) -- (0,-1); \draw[graf,fill=white] (0,-1) circle (2.5pt);
  \fill[graf] (0,0) circle (2.5pt);
  \node[left] at (0,1) {$1$}; \node[right] at (0,-1) {$-1$}; \node[below left] at (0,0) {$0$};
\end{tikzpicture}

\ovalbox{$\operatorname{Dom}f = \mathbb{R}$} \ \ovalbox{$\operatorname{Ran}f = \{-1; 0; 1\}$}
\end{minipage}

\medskip
\textbf{Función máximo entero:} Veamos algunas propiedades, definimos $x \in \mathbb{R}$ y
$n \in \mathbb{Z}$.

\begin{minipage}[t]{0.388\textwidth}
\begin{itemize}
  \item $[\![x]\!] \in \mathbb{Z}$
  \item $[\![x]\!] = x \leftrightarrow x \in \mathbb{Z}$
  \item $[\![x]\!] = n \leftrightarrow n \leq x < n + 1$
  \item $[\![x + n]\!] = [\![x]\!] + n$
  \item $[\![\,[\![x]\!]\,]\!] = [\![x]\!]$
\end{itemize}
\end{minipage}\hfill
\begin{minipage}[t]{0.582\textwidth}
Ejemplos:
\[
\begin{array}{llll}
[\![2{,}75]\!] = 2 & [\![0]\!] = 0 & [\![-4]\!] = -4 & [\![\sqrt{2}]\!] = 1\\[4pt]
[\![1{,}99]\!] = 1 & [\![5]\!] = 5 & [\![-3{,}02]\!] = -4 & [\![-\sqrt{2}]\!] = -2
\end{array}
\]
\end{minipage}

\begin{minipage}{0.291\textwidth}
\centering
\ovalbox{$f(x) = [\![x]\!]$}\\[8pt]
\ovalbox{$\operatorname{Dom}f = \mathbb{R}$}\\[4pt]
\ovalbox{$\operatorname{Ran}f = \mathbb{Z}$}
\end{minipage}\hfill
\begin{minipage}{0.407\textwidth}
\centering
\begin{tikzpicture}[scale=0.55]
  \ejes{-2.5}{5}{-2.5}{3.8}
  \foreach \n in {-2,...,3} {
    \draw[thick,graf] (\n,\n) -- (\n+1,\n);
    \fill[graf] (\n,\n) circle (3pt);
    \draw[graf,fill=white] (\n+1,\n) circle (3pt);
  }
  \foreach \n in {1,2,3,4} \node[below] at (\n,0) {\scriptsize $\n$};
  \foreach \n in {1,2,3} \node[left] at (0,\n) {\scriptsize $\n$};
  \node[below] at (-1,0) {\scriptsize $-1$}; \node[below] at (-2,0) {\scriptsize $-2$};
\end{tikzpicture}
\end{minipage}\hfill
\begin{minipage}{0.272\textwidth}
\[
\begin{array}{r|l}
-2 & \text{Si } -2 \leq x < -1\\
-1 & \text{Si } -1 \leq x < 0\\
0 & \text{Si } 0 \leq x < 1\\
1 & \text{Si } 1 \leq x < 2\\
2 & \text{Si } 2 \leq x < 3\\
3 & \text{Si } 3 \leq x < 4
\end{array}
\]
\end{minipage}

\begin{multicols}{2}
\begin{itemize}
  \item $[\![x]\!] < n \leftrightarrow x < n$
  \item $[\![x]\!] \geq n \leftrightarrow x \geq n$
  \item $[\![x]\!] \leq n \leftrightarrow x < n + 1$
  \item $[\![x]\!] > n \leftrightarrow x \geq n + 1$
\end{itemize}
\end{multicols}

\begin{minipage}[t]{0.485\textwidth}
\textbf{Función identidad:}
\begin{center}
\begin{tikzpicture}[scale=0.6]
  \ejes{-3}{4}{-3}{4}
  \draw[<->,thick,graf] (-2.8,-2.8) -- (3.5,3.5);
  \draw[dashed] (3,0) -- (3,3) -- (0,3); \draw[dashed] (-2,0) -- (-2,-2) -- (0,-2);
  \node[below] at (3,0) {\scriptsize 4}; \node[left] at (0,3) {\scriptsize 4};
  \node[above] at (-2,0) {\scriptsize $-2$}; \node[right] at (0,-2) {\scriptsize $-2$};
  \draw[rojo] (0.6,0) arc (0:45:0.6); \node[rojo] at (1.1,0.3) {\scriptsize $45^\circ$};
  \node at (2,3.6) {\ovalbox{$f(x) = x$}};
\end{tikzpicture}

\ovalbox{$\operatorname{Dom}f = \mathbb{R}$} \ \ovalbox{$\operatorname{Ran}f = \mathbb{R}$}
\end{center}
\end{minipage}\hfill
\begin{minipage}[t]{0.485\textwidth}
\textbf{Función valor absoluto:}
\begin{center}
\begin{tikzpicture}[scale=0.6]
  \ejes{-4}{4}{-0.5}{4}
  \draw[<->,thick,graf] (-3.5,3.5) -- (0,0) -- (3.5,3.5);
  \draw[rojo] (0.6,0) arc (0:45:0.6); \node[rojo] at (1.1,0.3) {\scriptsize $45^\circ$};
  \draw[rojo] (-0.6,0) arc (180:135:0.6); \node[rojo] at (-1.1,0.3) {\scriptsize $45^\circ$};
  \node at (2.5,4) {\ovalbox{$f(x) = |x|$}};
\end{tikzpicture}

\ovalbox{$f(x) = |x| = \begin{cases} x; & x \geq 0\\ -x; & x < 0\end{cases}$}\\[4pt]
\ovalbox{$\operatorname{Dom}f = \mathbb{R}$} \ \ovalbox{$\operatorname{Ran}f = \mathbb{R}^+_0$}
\end{center}
\end{minipage}

\medskip
\begin{minipage}[t]{0.485\textwidth}
\textbf{Función raíz cuadrada:}
\begin{center}
\begin{tikzpicture}[scale=0.6]
  \ejes{-0.5}{5}{-0.5}{3}
  \draw[->,thick,graf,domain=0:4.8,samples=60] plot (\x,{sqrt(\x)});
  \node at (3,2.7) {\ovalbox{$f(x) = \sqrt{x}$}};
\end{tikzpicture}

\ovalbox{$\operatorname{Dom}f = \mathbb{R}^+_0$} \ \ovalbox{$\operatorname{Ran}f = \mathbb{R}^+_0$}
\end{center}
\end{minipage}\hfill
\begin{minipage}[t]{0.485\textwidth}
\textbf{Función inverso multiplicativo:}
\begin{center}
\begin{tikzpicture}[scale=0.6]
  \ejes{-3.5}{3.5}{-3.5}{3.5}
  \draw[thick,graf,domain=0.3:3.3,samples=60] plot (\x,{1/\x});
  \draw[thick,graf,domain=-3.3:-0.3,samples=60] plot (\x,{1/\x});
  \draw[dashed] (1,0) -- (1,1) -- (0,1); \draw[dashed] (-1,0) -- (-1,-1) -- (0,-1);
  \node[below] at (1,0) {\scriptsize 1}; \node[left] at (0,1) {\scriptsize 1};
  \node[above] at (-1,0) {\scriptsize $-1$}; \node[right] at (0,-1) {\scriptsize $-1$};
  \node at (2.3,2.8) {\ovalbox{$f(x) = \frac{1}{x}$}};
\end{tikzpicture}

\ovalbox{$\operatorname{Dom}f = \mathbb{R} - \{0\}$}\\[2pt]
\ovalbox{$\operatorname{Ran}f = \mathbb{R} - \{0\}$}
\end{center}
\end{minipage}

\subsection*{Funciones polinomiales}
Son aquellas cuya regla de correspondencia está asociada a un polinomio. Entre las
principales tenemos:

\textbf{Función lineal:}
\begin{itemize}
  \item Regla de correspondencia: \ovalbox{$f(x) = mx + b;\ m \neq 0$}
  \item Dominio: \ovalbox{$\operatorname{Dom}f = \mathbb{R}$} $\wedge$ Rango: \ovalbox{$\operatorname{Ran}f = \mathbb{R}$}
\end{itemize}
\begin{center}
\begin{tikzpicture}[scale=0.6]
  \begin{scope}
    \ejes{-3.5}{3.5}{-2.5}{3.5}
    \draw[<->,thick,graf] (-3.2,-2.3) -- (2.5,3.4);
    \fill (0,0.9) circle (2pt) node[right] {$b$};
    \node[below] at (-0.9,0) {\scriptsize $-\frac{b}{m}$};
    \draw[rojo] (-0.4,0) arc (0:45:0.5); \node[rojo] at (-0.1,0.35) {\scriptsize $\theta$};
    \node[rojo] at (-2.5,2) {$m > 0$};
  \end{scope}
  \begin{scope}[xshift=9cm]
    \ejes{-3.5}{3.5}{-2.5}{3.5}
    \draw[<->,thick,graf] (-2.5,3.4) -- (3.2,-2.3);
    \fill (0,0.9) circle (2pt) node[right] {$b$};
    \node[below] at (0.9,-0.1) {\scriptsize $-\frac{b}{m}$};
    \draw[rojo] (1.4,0) arc (0:135:0.5); \node[rojo] at (1.6,0.5) {\scriptsize $\theta$};
    \node[rojo] at (-2,1.5) {$m < 0$};
  \end{scope}
\end{tikzpicture}
\end{center}

\textbf{Función cuadrática:}
\begin{itemize}
  \item Regla de correspondencia: \ovalbox{$f(x) = ax^2 + bx + c,\ a \neq 0$}
  \item Dominio: \ovalbox{$\operatorname{Dom}f = \mathbb{R}$}
  \item Gráfica: Su gráfica $G_f$ es una parábola simétrica respecto a una recta vertical
  (llamada eje de simetría) que pasa por el vértice de la parábola.

  Si $a > 0$, la parábola abre hacia arriba: $\cup$\\
  Si $a < 0$, la parábola abre hacia abajo: $\cap$
\end{itemize}
\begin{center}
\begin{tikzpicture}[scale=0.55]
  \begin{scope}
    \ejes{-3}{3.5}{-0.5}{4.5}
    \draw[thick,graf,domain=-2:2,samples=50] plot (\x,{\x*\x});
    \node at (3,3.5) {\ovalbox{$f(x) = x^2$}};
  \end{scope}
  \begin{scope}[xshift=10cm,yshift=4cm]
    \ejes{-3}{3.5}{-4.5}{0.8}
    \draw[thick,graf,domain=-2:2,samples=50] plot (\x,{-\x*\x});
    \node at (3.6,-2) {\ovalbox{$f(x) = -x^2$}};
  \end{scope}
\end{tikzpicture}
\end{center}
\begin{itemize}
  \item Vértice de la parábola:
  \ $V(h;k) = V\left(-\dfrac{b}{2a}; -\dfrac{b^2 - 4ac}{4a}\right)$
\end{itemize}

\textcolor{verde}{\textbf{Análisis de la gráfica de la función cuadrática:}} Sea
$f(x) = ax^2 + bx + c,\ a \neq 0$, de raíces $x_1, x_2$ y discriminante $\Delta = b^2 - 4ac$.

\begin{center}
\begin{tabular}{|c|c|c|c|}
\hline
$f(x)$ & \textcolor{naranja}{$\Delta > 0$} & \textcolor{naranja}{$\Delta = 0$} & \textcolor{naranja}{$\Delta < 0$}\\ \hline
\textcolor{verde}{$a > 0$} &
\parabola{1}{-1}{0} & \parabola{1}{0}{0} & \parabola{1}{1}{0}\\ \hline
\textcolor{verde}{$a < 0$} &
\parabola{-1}{1}{0} & \parabola{-1}{0}{0} & \parabola{-1}{-1}{0}\\ \hline
\end{tabular}
\end{center}

\textbf{Función cúbica:}
\begin{itemize}
  \item Regla de correspondencia: \ovalbox{$f(x) = ax^3 + bx^2 + cx + d$}
  \item Dominio: \ovalbox{$\operatorname{Dom}f = \mathbb{R}$} $\wedge$ Rango: \ovalbox{$\operatorname{Ran}f = \mathbb{R}$}
  \item Gráfica: Su gráfica $G_f$ es una curva que corta en al menos un punto al eje $X$, y
  al eje $Y$ en el punto $(0;d)$.
\end{itemize}


% ============================================================
% Páginas 58-60
% ============================================================
\begin{center}
\begin{tikzpicture}[scale=0.55]
  \begin{scope}
    \ejes{-3}{3.5}{-3.5}{3.5}
    \draw[thick,graf,domain=-1.5:1.5,samples=60] plot (\x,{\x*\x*\x});
    \node at (2.5,3) {\ovalbox{$f(x) = x^3$}};
  \end{scope}
  \begin{scope}[xshift=10cm]
    \ejes{-3}{3.5}{-3.5}{3.5}
    \draw[thick,graf,domain=-1.5:1.5,samples=60] plot (\x,{-\x*\x*\x});
    \node at (-2.3,3) {\ovalbox{$f(x) = -x^3$}};
  \end{scope}
\end{tikzpicture}
\end{center}

\textcolor{verde}{\textbf{Análisis de la gráfica de la función cúbica reducida:}} Sea
$f(x) = a(x^3 + px + q);\ a \neq 0$ una función cúbica reducida de raíces: $x_1, x_2$ y
$x_3$. La gráfica $f$ depende de $a$ y la expresión
$\Delta = \left(\dfrac{p}{3}\right)^3 + \left(\dfrac{q}{2}\right)^2$, tal como se observa en
los siguientes casos.

\begin{center}
\begin{tabular}{|c|c|c|c|}
\hline
$f(x)$ & \textcolor{naranja}{$\Delta < 0$} & \textcolor{naranja}{$\Delta = 0$} & \textcolor{naranja}{$\Delta > 0$}\\ \hline
\textcolor{verde}{$a > 0$} &
\begin{tikzpicture}[scale=0.45]
  \ejes{-2.6}{2.8}{-2}{2}
  \draw[thick,graf,domain=-2.15:2.15,samples=60] plot (\x,{0.4*(\x*\x*\x-3*\x)});
  \node[below left] at (-1.73,0) {\tiny $x_1$}; \node[below left] at (0,0) {\tiny $x_2$}; \node[below right] at (1.73,0) {\tiny $x_3$};
\end{tikzpicture} &
\begin{tikzpicture}[scale=0.45]
  \ejes{-2.6}{2.8}{-2}{2}
  \draw[thick,graf,domain=-2.3:1.9,samples=60] plot (\x,{0.4*(\x*\x*\x-3*\x+2)});
  \node[below] at (-2,0) {\tiny $x_1$}; \node[below] at (1,0) {\tiny $x_2 = x_3$};
\end{tikzpicture} &
\begin{tikzpicture}[scale=0.45]
  \ejes{-2.6}{2.8}{-2}{2}
  \draw[thick,graf,domain=-1.8:1.6,samples=60] plot (\x,{0.4*(\x*\x*\x+\x+1)});
  \node[below] at (-0.7,0) {\tiny $x_1$};
\end{tikzpicture}\\ \hline
\textcolor{verde}{$a < 0$} &
\begin{tikzpicture}[scale=0.45]
  \ejes{-2.6}{2.8}{-2}{2}
  \draw[thick,graf,domain=-2.15:2.15,samples=60] plot (\x,{-0.4*(\x*\x*\x-3*\x)});
  \node[below left] at (-1.73,0) {\tiny $x_1$}; \node[below right] at (0,0) {\tiny $x_2$}; \node[below right] at (1.73,0) {\tiny $x_3$};
\end{tikzpicture} &
\begin{tikzpicture}[scale=0.45]
  \ejes{-2.6}{2.8}{-2}{2}
  \draw[thick,graf,domain=-1.9:2.3,samples=60] plot (\x,{-0.4*(\x*\x*\x-3*\x-2)});
  \node[below] at (-1,0) {\tiny $x_2 = x_3$}; \node[below] at (2,0) {\tiny $x_1$};
\end{tikzpicture} &
\begin{tikzpicture}[scale=0.45]
  \ejes{-2.6}{2.8}{-2}{2}
  \draw[thick,graf,domain=-1.6:1.8,samples=60] plot (\x,{-0.4*(\x*\x*\x+\x-1)});
  \node[below] at (0.7,0) {\tiny $x_1$};
\end{tikzpicture}\\ \hline
\end{tabular}
\end{center}

\begin{itemize}
  \item Si $\Delta < 0$, $x_1, x_2$ y $x_3 \in \mathbb{R}$ y $x_1 \neq x_2 \neq x_3$
  \item Si $\Delta = 0$, $x_1, x_2$ y $x_3 \in \mathbb{R}$ y $x_2 = x_3$ y $x_1 = -2x_2$
  \item Si $\Delta > 0$, $x_1 \in \mathbb{R} \wedge x_2, x_3 \in \mathbb{C}$
  \textcolor{magenta}{(Imaginarias conjugadas)}
\end{itemize}

\textbf{Función de grado superior:} Siempre y cuando sea factorizable.
\[
\ovalbox{$f(x) = a_0(x - r_1)^{\alpha_1}(x - r_2)^{\alpha_2}(x - r_3)^{\alpha_3}\cdots(x - r_k)^{\alpha_k}$}
\qquad \ovalbox{$\operatorname{Dom}f = \mathbb{R}$}
\]

\textcolor{verde}{\textbf{Considere:}}

\textcolor{magenta}{\textbf{i)}} Comportamiento en el extremo derecho:
\begin{center}
\begin{tikzpicture}[scale=0.5]
  \begin{scope}
    \ejes{-3}{3.5}{-2.5}{3}
    \draw[->,thick,graf] (1.5,1) -- (2.8,2.8);
    \node at (0,-3.2) {Si $a_0 > 0$};
  \end{scope}
  \begin{scope}[xshift=9cm]
    \ejes{-3}{3.5}{-2.5}{3}
    \draw[->,thick,graf] (1.5,-1) -- (2.8,-2.8);
    \node at (0,-3.2) {Si $a_0 < 0$};
  \end{scope}
\end{tikzpicture}
\end{center}

\textcolor{magenta}{\textbf{ii)}} Comportamiento cerca de cada raíz $r_i$:

\begin{minipage}{0.533\textwidth}
\centering
\begin{tikzpicture}[scale=0.5]
  \begin{scope}
    \ejes{-1}{2.5}{-1}{1.5}
    \draw[thick,graf] (0.4,-0.6) -- (1.6,0.8); \node[below] at (0.9,0) {\scriptsize $r_i$};
  \end{scope}
  \begin{scope}[xshift=4.5cm]
    \ejes{-1}{2.5}{-1}{1.5}
    \draw[thick,naranja] (0.4,0.8) -- (1.6,-0.6); \node[below] at (1.1,0) {\scriptsize $r_i$};
  \end{scope}
\end{tikzpicture}
\end{minipage}\hfill
\begin{minipage}{0.436\textwidth}
Si $r_i$ es raíz simple o de multiplicidad impar.
\end{minipage}

\medskip
\begin{minipage}{0.533\textwidth}
\centering
\begin{tikzpicture}[scale=0.5]
  \begin{scope}
    \ejes{-1}{2.5}{-1}{1.5}
    \draw[thick,graf,domain=0.3:1.7,samples=30] plot (\x,{1.6*(\x-1)*(\x-1)}); \node[below] at (1,0) {\scriptsize $r_i$};
  \end{scope}
  \begin{scope}[xshift=4.5cm]
    \ejes{-1}{2.5}{-1}{1.5}
    \draw[thick,naranja,domain=0.3:1.7,samples=30] plot (\x,{-1.6*(\x-1)*(\x-1)}); \node[above] at (1,0) {\scriptsize $r_i$};
  \end{scope}
\end{tikzpicture}
\end{minipage}\hfill
\begin{minipage}{0.436\textwidth}
Si $r_i$ es raíz de multiplicidad par.
\end{minipage}

\medskip
\textbf{Ejemplo:} \quad
$f(x) = 3\underset{\text{Simple}}{(x+1)}\underset{\text{M. par}}{(x-3)^2}\underset{\text{M. impar}}{(x-5)^3} \longrightarrow a_0 = 3 > 0$

\begin{center}
\begin{tikzpicture}[scale=0.8]
  \draw[->] (-3,0) -- (7,0) node[below] {$x$};
  \draw[->] (0,-1.5) -- (0,3.5) node[left] {$y$};
  \draw[thick] plot[smooth,tension=0.7] coordinates {(-2.2,3.2) (-1,0) (1,-1) (3,0) (4,-1) (5,0) (5.8,3.2)};
  \foreach \p in {-1,3,5} {\fill (\p,0) circle (1.5pt); \node[below] at (\p,-0.05) {\small $\p$};}
\end{tikzpicture}
\end{center}

\textbf{Función potencial:}
\begin{itemize}
  \item Regla de correspondencia: \ovalbox{$f(x) = x^n$}$;\ n \in \mathbb{Z}^+ \wedge n \geq 2$
  \item Dominio: \ovalbox{$\operatorname{Dom}f = \mathbb{R}$}
\end{itemize}
\begin{center}
\begin{tikzpicture}[scale=1.1]
  \begin{scope}
    \ejes{-2}{2.2}{-0.4}{2.6}
    \draw[thick,graf,domain=-1.5:1.5,samples=50] plot (\x,{\x*\x}) node[right] {\small $x^2$};
    \draw[thick,magenta,domain=-1.25:1.25,samples=50] plot (\x,{\x*\x*\x*\x}) node[left] {\small $x^4$};
    \draw[dashed] (-1,0) -- (-1,1) -- (1,1) -- (1,0);
    \node[below] at (-1,0) {\scriptsize $-1$}; \node[below] at (1,0) {\scriptsize $1$};
    \node[above left] at (0,1) {\scriptsize $1$};
    \node at (-1.6,2.3) {\small $n$ par};
  \end{scope}
  \begin{scope}[xshift=5cm,yshift=1cm]
    \ejes{-2}{2.2}{-1.8}{1.8}
    \draw[thick,graf,domain=-1.2:1.2,samples=50] plot (\x,{\x*\x*\x}) node[right] {\small $x^3$};
    \draw[thick,magenta,domain=-1.12:1.12,samples=50] plot (\x,{\x*\x*\x*\x*\x}) node[left] {\small $x^5$};
    \draw[dashed] (-1,0) -- (-1,-1) -- (0,-1); \draw[dashed] (1,0) -- (1,1) -- (0,1);
    \node[above] at (-1,0) {\scriptsize $-1$}; \node[below] at (1,0) {\scriptsize $1$};
    \node[left] at (0,1) {\scriptsize $1$}; \node[right] at (0,-1) {\scriptsize $-1$};
    \node at (-1.6,1.5) {\small $n$ impar};
  \end{scope}
\end{tikzpicture}
\end{center}

\subsection*{Desplazamiento y reflejos de la gráfica de una función}
Conociendo la gráfica donde: $F: y = f(x)$
\begin{center}
\begin{tikzpicture}[scale=0.7]
  \ejes{-1.5}{3}{-1.2}{1.8}
  \rampa{0}{0}{graf}
\end{tikzpicture}
\end{center}
y considerando un número positivo $h$, entonces:

\begin{minipage}[t]{0.485\textwidth}
\textbf{Desplazamiento horizontal:}
\begin{center}
\begin{tikzpicture}[scale=0.6]
  \ejes{-4}{4.5}{-1.5}{2.3}
  \rampa{0}{0}{graf}
  \rampa{-2.2}{0}{magenta,dashed}
  \rampa{2.2}{0}{magenta,dashed}
  \node[below] at (-2.2,0) {\small $-h$}; \node[below] at (2.2,0) {\small $h$};
  \node[above] at (-1.5,1.1) {\scriptsize $f(x+h)$};
  \node[above] at (1.1,1.1) {\scriptsize $f(x)$};
  \node[above] at (3.4,1.1) {\scriptsize $f(x-h)$};
\end{tikzpicture}

{\small ``$h$'' unidades a la izquierda \quad ``$h$'' unidades a la derecha}
\end{center}
\end{minipage}\hfill
\begin{minipage}[t]{0.485\textwidth}
\textbf{Desplazamiento vertical:}
\begin{center}
\begin{tikzpicture}[scale=0.6]
  \ejes{-2}{4}{-3}{3.5}
  \rampa{0}{0}{graf}
  \rampa{0}{1.6}{magenta,dashed}
  \rampa{0}{-1.6}{magenta,dashed}
  \node[left] at (0,1.6) {\small $h$}; \node[left] at (0,-1.6) {\small $-h$};
  \node[right] at (2.3,2.6) {\scriptsize $f(x)+h$};
  \node[right] at (2.3,1) {\scriptsize $f(x)$};
  \node[right] at (2.3,-0.6) {\scriptsize $f(x)-h$};
\end{tikzpicture}

{\small ``$h$'' unidades hacia arriba / ``$h$'' unidades hacia abajo}
\end{center}
\end{minipage}


% ============================================================
% Páginas 61-63
% ============================================================
\begin{minipage}[t]{0.485\textwidth}
\textbf{Reflejo con respecto al eje ``$x$'':}
\begin{center}
\begin{tikzpicture}[scale=0.6]
  \ejes{-2}{4}{-2.5}{2.5}
  \draw[thick,graf] (-1.2,-1.2) -- (1.2,1.2) -- (2.5,1.2);
  \draw[thick,magenta,dashed] (-1.2,1.2) -- (1.2,-1.2) -- (2.5,-1.2);
  \draw[dotted] (-1.2,-1.2) -- (-1.2,1.2);
  \node[above] at (2.2,1.2) {\small $f(x)$};
  \node[below] at (2.2,-1.2) {\small $-f(x)$};
\end{tikzpicture}

{\small El eje ``$x$'' se comporta como un espejo.}
\end{center}
\end{minipage}\hfill
\begin{minipage}[t]{0.485\textwidth}
\textbf{Reflejo con respecto al eje ``$y$'':}
\begin{center}
\begin{tikzpicture}[scale=0.6]
  \ejes{-3.5}{3.5}{-2}{2.5}
  \draw[thick,graf] (-1,-1) -- (1,1) -- (2.5,1);
  \draw[thick,magenta,dashed] (1,-1) -- (-1,1) -- (-2.5,1);
  \draw[dotted] (-1,-1) -- (1,-1);
  \node[above] at (2.2,1) {\small $f(x)$};
  \node[above] at (-2.2,1) {\small $f(-x)$};
\end{tikzpicture}

{\small El eje ``$y$'' se comporta como un espejo.}
\end{center}
\end{minipage}

\medskip
\textbf{Reflejo por valor absoluto:}

\begin{minipage}{0.291\textwidth}
Si parte de una función se encuentra en el IIIC o IVC.
\end{minipage}\hfill
\begin{minipage}{0.388\textwidth}
\centering
\begin{tikzpicture}[scale=0.6]
  \ejes{-2.5}{4}{-2}{2.5}
  \draw[thick,graf] (-1.5,-1.5) -- (0,0) -- (1.2,1.2) -- (2.5,1.2);
  \draw[thick,magenta,dashed] (-1.5,1.5) -- (0,0);
  \draw[dotted] (-1.5,-1.5) -- (-1.5,1.5);
  \node[above] at (2.2,1.2) {\small $f(x)$};
  \node[above] at (-1.2,1.6) {\small $|f(x)|$};
\end{tikzpicture}
\end{minipage}\hfill
\begin{minipage}{0.291\textwidth}
El valor absoluto refleja esa parte.
\end{minipage}

\newpage
% ------------------------------------------------------------
\section{Funciones III}

\subsection*{Álgebra de funciones}

\textbf{Igualdad de funciones:} Dos funciones $F$ y $G$ son iguales si cumplen las condiciones:
\begin{enumerate}[label=\roman*)]
  \item $\operatorname{Dom}F = \operatorname{Dom}G$
  \item $F(x) = G(x),\ \forall x \in \operatorname{Dom}F = \operatorname{Dom}G$
\end{enumerate}

Si $F$ y $G$ son dos funciones que cumplen la condición
$A = \operatorname{Dom}F \cap \operatorname{Dom}G \neq \varnothing$, se definen operaciones
entre estas, como sigue:

\begin{minipage}[t]{0.320\textwidth}
\textbf{Función adición:}
\begin{gather*}
(F + G)(x) = F(x) + G(x)\\
\operatorname{Dom}(F + G) = A
\end{gather*}
\end{minipage}\hfill
\begin{minipage}[t]{0.320\textwidth}
\textbf{Función diferencia:}
\begin{gather*}
(F - G)(x) = F(x) - G(x)\\
\operatorname{Dom}(F - G) = A
\end{gather*}
\end{minipage}\hfill
\begin{minipage}[t]{0.320\textwidth}
\textbf{Función producto:}
\begin{gather*}
(F\cdot G)(x) = F(x)\cdot G(x)\\
\operatorname{Dom}(F\cdot G) = A
\end{gather*}
\end{minipage}

\begin{minipage}[t]{0.388\textwidth}
\textbf{Función cociente:}
\begin{gather*}
(F/G)(x) = F(x)/G(x)\\
\operatorname{Dom}(F/G) = A - \{x \,/\, G(x) = 0\}
\end{gather*}
\end{minipage}\hfill
\begin{minipage}[t]{0.582\textwidth}
\textbf{Potencia de una función:}
\begin{gather*}
F(x)^n = \underbrace{F(x)\cdot F(x)\cdot F(x)\cdots F(x)}_{\text{``$n$'' veces}};\ \forall n \in \mathbb{Z}^+\\
\operatorname{Dom}(F^n) = \operatorname{Dom}(F)
\end{gather*}
\end{minipage}

\medskip
\textbf{Ejemplo:}
\[
\left.
\begin{array}{l}
f = \{(0;2), (1;-2), (2;3), (3;0)\}\\
g = \{(0;-1), (1;0), (2;1), (3;4)\}
\end{array}
\right\}
\ovalbox{$\operatorname{Dom}f \cap \operatorname{Dom}g = \{0; 1; 2; 3\}$}
\]
\begin{minipage}[t]{0.485\textwidth}
\begin{itemize}
  \item $f + g = \{(0;1), (1;-2), (2;4), (3;4)\}$
  \item $f - g = \{(0;3), (1;-2), (2;2), (3;-4)\}$
  \item $f\cdot g = \{(0;-2), (1;0), (2;3), (3;0)\}$
  \item $f^2 = \{(0;4), (1;4), (2;9), (3;0)\}$
  \item $2f = \{(0;4), (1;-4), (2;6), (3;0)\}$
\end{itemize}
\end{minipage}\hfill
\begin{minipage}[t]{0.485\textwidth}
\begin{itemize}
  \item $\dfrac{f}{g} = \{(0;-2), (2;3), (3;0)\}$
  \item $\dfrac{g}{f} = \left\{\left(0;-\tfrac{1}{2}\right), (1;0), \left(2;\tfrac{1}{3}\right)\right\}$
  \item $g^3 = \{(0;-1), (1;0), (2;1), (3;64)\}$
  \item $3g = \{(0;-3), (1;0), (2;3), (3;12)\}$
\end{itemize}
\end{minipage}

\medskip
\textbf{Composición de funciones:} Dadas las funciones $F$ y $G$ definimos la función
compuesta $F$ con $G$ (en ese orden y denotada $F\circ G$) a aquella que cumple:
\[
\ovalbox{$\operatorname{Dom}(F\circ G) = \{x \,/\, x \in \operatorname{Dom}G \wedge G(x) \in \operatorname{Dom}F\}$}
\qquad
\ovalbox{$F\circ G(x) = F(G(x))$}
\]

\begin{center}
\begin{tikzpicture}[>=Stealth, scale=0.9]
  \draw[graf] (0,0) ellipse (0.9 and 1.5); \node[rojo] at (0,-1.8) {\small $\operatorname{Dom}G$};
  \draw[graf] (3,0) ellipse (0.9 and 1.5); \node[rojo] at (2.8,-1.8) {\small $\operatorname{Ran}G$};
  \draw[magenta] (4.1,0) ellipse (0.9 and 1.5); \node[rojo] at (4.3,-1.8) {\small $\operatorname{Dom}F$};
  \draw[magenta] (7,0) ellipse (0.9 and 1.5); \node[rojo] at (7,-1.8) {\small $\operatorname{Ran}F$};
  \node (x) at (0,0.3) {$x$};
  \node (g) at (3.6,0.3) {\small $G(x)$};
  \node (f) at (7,0.3) {\small $F(G(x))$};
  \draw[->] (x) -- (g); \draw[->] (g) -- (f);
  \draw[->,rojo!60!black] (0.2,1.6) to[bend left=20] node[above] {$G$} (2.8,1.6);
  \draw[->,rojo!60!black] (4.2,1.6) to[bend left=20] node[above] {$F$} (6.8,1.6);
  \draw[->,verde,thick] (0.3,-1.4) to[bend right=25] node[below] {$F\circ G$} (6.6,-1.2);
\end{tikzpicture}
\end{center}

\textbf{Teoremas:}

\begin{minipage}[t]{0.485\textwidth}
\begin{itemize}[label=$\ast$]
  \item $\exists\, F\circ G \leftrightarrow \operatorname{Ran}G \cap \operatorname{Dom}F \neq \varnothing$
  \item $F\circ G \neq G\circ F$
  \item $(F\circ G)\circ H = F\circ(G\circ H)$
\end{itemize}
\end{minipage}\hfill
\begin{minipage}[t]{0.485\textwidth}
\begin{itemize}[label=$\ast$]
  \item $F\circ I = I\circ F$
  \item $(F + G)\circ H = F\circ H + G\circ H$
  \item $(F\cdot G)\circ H = (F\circ H)(G\circ H)$
\end{itemize}
\end{minipage}

\medskip
\textbf{Ejemplo:}

\begin{minipage}[t]{0.436\textwidth}
\begin{gather*}
f(x) = x^2 - 3;\ x \in [-2; 7\rangle\\
g(x) = 2x - 1;\ x \in [-1; 5\rangle\\[6pt]
\textcolor{naranja}{\text{¿}\operatorname{Dom}(f\circ g)\text{?}}
\end{gather*}
\end{minipage}\hfill
\begin{minipage}[t]{0.533\textwidth}
\begin{gather*}
x \in \operatorname{Dom}g \;\wedge\; g(x) \in \operatorname{Dom}f\\
x \in [-1; 5\rangle \;\wedge\; 2x - 1 \in [-2; 7\rangle\\
2x \in [-1; 8\rangle\\
x \in [-1/2; 4\rangle
\end{gather*}
\end{minipage}

\begin{center}
\begin{tikzpicture}[x=1cm]
  \draw[<->] (-0.5,0) -- (7,0);
  \node[below] at (-0.3,0) {\small $-\infty$}; \node[below] at (6.8,0) {\small $+\infty$};
  \foreach \p/\l in {1/-1, 2.2/-1/2, 4.2/4, 5.6/5} \node[below] at (\p,-0.1) {\small $\l$};
  \draw[thick,graf] (1,0) -- (1,0.6) -- (5.6,0.6) -- (5.6,0);
  \draw[thick,magenta] (2.2,0) -- (2.2,0.8) -- (4.2,0.8) -- (4.2,0);
  \fill[pattern=north east lines] (2.2,0) rectangle (4.2,0.6);
  \fill[graf] (1,0) circle (2pt); \draw[graf,fill=white] (5.6,0) circle (2pt);
  \fill[magenta] (2.2,0) circle (2pt); \draw[magenta,fill=white] (4.2,0) circle (2pt);
\end{tikzpicture}
\end{center}
\[\therefore\ \operatorname{Dom}f\circ g = [-1/2; 4\rangle\]
\begin{gather*}
f\circ g(x) = f(g(x)) = g^2(x) - 3\\
f\circ g(x) = (2x - 1)^2 - 3\\
\therefore\ f\circ g(x) = 4x^2 - 4x - 2
\end{gather*}

\subsection*{Función inversa}

\textbf{Función inyectiva (Univalente):}
\[
\ovalbox{$x_1 \neq x_2 \in \operatorname{Dom}f \longrightarrow f(x_1) \neq f(x_2)$}
\;\ldots\text{negando}\ldots\;
\ovalbox{$f(x_1) = f(x_2) \longrightarrow x_1 = x_2$}
\]
Ejemplo:

\begin{minipage}[t]{0.320\textwidth}
\begin{gather*}
f(x) = x^2 - 4x - 2\\
f(x) = (x-2)^2 - 6\\
f(x_1) = f(x_2)
\end{gather*}
\end{minipage}\hfill
\begin{minipage}[t]{0.359\textwidth}
\begin{gather*}
(x_1 - 2)^2 - 6 = (x_2 - 2)^2 - 6\\
\sqrt{\phantom{x}}:\ |x_1 - 2| = |x_2 - 2|\\
x_1 - 2 = x_2 - 2 \;\vee\; x_1 - 2 = 2 - x_2
\end{gather*}
\end{minipage}\hfill
\begin{minipage}[t]{0.291\textwidth}
\begin{gather*}
x_1 = x_2 \;\vee\; \boxed{x_1 + x_2 = 4}\\
\therefore\ f \text{ no es inyectiva}
\end{gather*}
\end{minipage}


% ============================================================
% Páginas 64-66
% ============================================================
\textcolor{naranja}{\textbf{Gráficamente:}}

\begin{minipage}{0.436\textwidth}
\centering
\begin{tikzpicture}[scale=0.55]
  \draw[<->] (-1.5,0) -- (6,0) node[right] {$x$};
  \draw[->] (0,-4) -- (0,4.5) node[left] {$y$};
  \draw[thick,graf,domain=-0.5:4.5,samples=60] plot (\x,{0.9*((\x-2)*(\x-2)-6)*0.5});
  \draw[<->,thick] (-1,2.8) -- (5.5,2.8);
  \fill (-0.27,2.8) circle (2pt) node[above left] {\scriptsize $x_1$};
  \fill (4.27,2.8) circle (2pt) node[above right] {\scriptsize $x_2$};
  \node[rojo] at (2,-3.6) {\small No es inyectiva};
\end{tikzpicture}
\end{minipage}\hfill
\begin{minipage}{0.533\textwidth}
\begin{gather*}
f(x) = 0\\
(x-2)^2 - 6 = 0\\
x_1 = 2 + \sqrt{6} \;\vee\; x_2 = 2 - \sqrt{6}
\end{gather*}
$\ast$ Toda recta horizontal corta al gráfico de una función inyectiva en un punto.
\end{minipage}

\medskip
\textbf{Función suryectiva (Sobreyectiva, Epiyectiva):}
\[f: A \longrightarrow B \text{ es suryectiva cuando } \ovalbox{$B = \operatorname{Ran}f$}\]

\textbf{Función biyectiva:} $f$ es biyectiva cuando es inyectiva y suryectiva.

\textbf{Función inversa:} Sea la función $f$, si es biyectiva entonces posee inversa
simbolizada por $f^*$ y está definida por:
\[
\ovalbox{$f: A \to B$} \qquad\qquad \ovalbox{$f^*: B \to A$}
\]
\[
\ovalbox{$f = \{(x;y) \,/\, y = f(x);\ x \in \operatorname{Dom}f\}$}
\quad
\ovalbox{$f^{-1} = f^* = \{(y;x) \,/\, x = f^*(y);\ y \in \operatorname{Dom}f^* = \operatorname{Ran}f\}$}
\]
\begin{center}
\begin{tikzpicture}[>=Stealth]
  \draw[graf] (0,0) ellipse (0.9 and 1.5); \node at (0,1.8) {$A$};
  \draw[graf] (5,0) ellipse (0.9 and 1.5); \node at (5,1.8) {$B$};
  \node at (0,0) {$x$}; \node at (5,0) {$y = f(x)$};
  \draw[->] (0.6,0.6) -- node[above] {$f$} (4.6,0.6);
  \draw[<-] (0.6,-0.6) -- node[below] {$f^{-1}$} (4.6,-0.6);
\end{tikzpicture}
\end{center}
\[
\ovalbox{$\operatorname{Dom}f^{-1} = \operatorname{Ran}f$}
\qquad
\ovalbox{$\operatorname{Ran}f^{-1} = \operatorname{Dom}f$}
\]

\textcolor{rojo}{\textbf{Teoremas:}} La gráfica de una función $f$ y la gráfica de su inversa
son simétricas con respecto a la recta identidad $y = x$.

\subsection*{\textcolor{magenta}{Método para calcular la inversa de una función}}
\begin{enumerate}
  \item Sustituye $f(x)$ por $y$.
  \item Intercambia entre ellas todas las $x$ por las $y$.
  \item Despeja para $y$.
  \item Sustituye $y$ por $f^{-1}(x)$.
\end{enumerate}

\textcolor{naranja}{\textbf{Teoremas:}} Sean $f$ y $g$ biyectivas, también $f\circ g$ existe.

\begin{minipage}[t]{0.533\textwidth}
\begin{enumerate}
  \item $(f^{-1})^{-1} = f$
  \item $(f^{-1}\circ f)(x) = x;\ x \in \operatorname{Dom}f$
  \item $(f\circ f^{-1})(x) = x;\ x \in \operatorname{Dom}f^{-1}$
  \item $(f\circ g)^{-1} = g^{-1}\circ f^{-1}$
  \item $h = f\circ g \rightarrow f = h\circ g^{-1}$
  \item $h = f\circ g \rightarrow g = f^{-1}\circ h$
\end{enumerate}
\end{minipage}\hfill
\begin{minipage}[t]{0.436\textwidth}
\begin{enumerate}\setcounter{enumi}{6}
  \item $f^*\circ f = I$
  \item $g\circ g^* = I$
\end{enumerate}
\end{minipage}

\medskip
\textbf{Ejemplo:}

\begin{minipage}[t]{0.380\textwidth}\small
\begin{gather*}
f(x) = x^2 - 4x - 2;\ x \in [3; 6\rangle\\
f(x) = (x-2)^2 - 6
\end{gather*}
\begin{itemize}
  \item $f(x_1) = f(x_2) \rightarrow x_1 = x_2$
  \begin{gather*}
  (x_1 - 2)^2 - 6 = (x_2 - 2)^2 - 6\\
  \sqrt{\phantom{x}}:\ |\overset{+}{x_1 - 2}| = |\overset{+}{x_2 - 2}|\\
  x_1 = x_2 \;:\; f \text{ es inyectiva}
  \end{gather*}
\end{itemize}
\begin{center}
\begin{tikzpicture}[scale=0.55]
  \draw[->] (-1,0) -- (5,0) node[right] {$x$};
  \draw[->] (0,-1.5) -- (0,5) node[left] {$y$};
  \draw[thick] (-1,-1) -- (4.8,4.8) node[right] {\small $y = x$};
  \draw[thick,rojo!70!black,domain=1.8:3.2,samples=40] plot (\x,{1.8*(\x-1.8)*(\x-1.8)+1.8*0}) node[right] {\small $f$};
  \draw[thick,naranja,domain=1.8:3.2,samples=40] plot ({1.8*(\x-1.8)*(\x-1.8)},{\x}) ;
  \node[naranja] at (-0.5,2.6) {\small $f^*$};
\end{tikzpicture}
\end{center}
\end{minipage}\hfill
\begin{minipage}[t]{0.330\textwidth}\small
\begin{itemize}
  \item $y = (x-2)^2 - 6$
  \begin{gather*}
  3 \leq x < 6 \rightarrow 1 \leq x - 2 < 4\\
  1 \leq (x-2)^2 < 16 \rightarrow -5 \leq y < 10\\
  \operatorname{Ran}f = [-5; 10\rangle : \operatorname{Dom}f^*
  \end{gather*}
  \item $y + 6 = (x - 2)^2$
  \[\sqrt{\phantom{x}}:\ \sqrt{y + 6} = |\overset{+}{x - 2}|\]
\end{itemize}
\end{minipage}\hfill
\begin{minipage}[t]{0.250\textwidth}
\begin{gather*}
x = 2 + \sqrt{y + 6}\\
\downarrow\\
y = 2 + \sqrt{x + 6}\\[6pt]
f^*(x) = 2 + \sqrt{x + 6}
\end{gather*}
\end{minipage}

\[
f\circ f^* = f(2 + \sqrt{x + 6}) = (2 + \sqrt{x + 6} - 2)^2 - 6 = x
\]


\newpage
% ------------------------------------------------------------
\section{Logaritmos en $\mathbb{R}$}

\textbf{Definición:} El logaritmo de un número ``$x$'' en una base dada ``$b$'', es el exponente
real ``$y$'' al que se debe elevarse la base para obtener el número dado.
\[\fbox{$\log_b x = y \longleftrightarrow b^y = x$}\]
\begin{center}
\fbox{\begin{tabular}{l}
Número: $x > 0$\\
Base: $b > 0,\ b \neq 1$\\
Logaritmo: $y \in \mathbb{R}$
\end{tabular}}
\end{center}

\textbf{Teoremas:} Los números del logaritmo son positivos y las bases positivas y diferentes
de 1.

\begin{minipage}[t]{0.485\textwidth}
\begin{enumerate}
  \item \textcolor{magenta}{Identidad fundamental:} \ $b^{\log_b N} = N$
  \item $\log_b \dfrac{M}{N} = \log_b M - \log_b N$
  \item $\log_b (M\cdot N) = \log_b M + \log_b N$
  \item $\log_b N^M = M\log_b N$
  \[\log_b \sqrt[M]{N} = \frac{1}{M}\log_b N\]
\end{enumerate}
\end{minipage}\hfill
\begin{minipage}[t]{0.485\textwidth}
\begin{enumerate}\setcounter{enumi}{4}
  \item $\log_b N = \log_{b^a} N^a = \log_{\sqrt[a]{b}} \sqrt[a]{N}$
  \item $\log_{b^q} N^p = \dfrac{p}{q}\log_b N;\ q \neq 0$
  \item $p^{\log_b N} = N^{\log_b p}$
  \item \textcolor{magenta}{Cambio de base:}
  \[\log_b N = \frac{\log_k N}{\log_k b}\;; \quad \log_b N = \frac{1}{\log_N b}\]
\end{enumerate}
\end{minipage}
\begin{enumerate}\setcounter{enumi}{8}
  \item \textcolor{magenta}{Regla de cadena:} \ $\log_b a\cdot\log_c b\cdot\log_d c = \log_d a$
\end{enumerate}

\textbf{Sistemas de logaritmos:} Dependen del valor de la base, por tanto, existen infinitos
sistemas de logaritmos, siendo los principales:

\textbf{I. Logaritmos decimales:}

Base: $10$, \ Notación: \ \fbox{$\log_{10} N = \log N$}
\[\log 2 = 0{,}30103\ldots \qquad \log 3 = 0{,}47712\ldots \qquad \log 5 = 1 - \log 2\]

\textbf{II. Logaritmos neperianos o naturales:}

Base: $e = 2{,}71828\ldots$, \ Notación: \ \fbox{$\log_e N = \ln N = L\,N$}


% ============================================================
% Páginas 67-69
% ============================================================
\textcolor{naranja}{\textbf{Cologaritmo:}} \quad
\fbox{$\operatorname{Colog}_b N = \log_b\left(\dfrac{1}{N}\right) = -\log_b N$}

\quad Donde: $N > 0;\ b > 0;\ b \neq 1$
\[\ast\ \operatorname{Colog}_2\left(\tfrac{1}{16}\right) = -\log_2\left(\tfrac{1}{16}\right) = 4\]

\textcolor{naranja}{\textbf{Antilogaritmo:}} \quad \fbox{$\operatorname{Antilog}_b x = b^x$}

\quad Donde: $x \in \mathbb{R};\ b > 0;\ b \neq 1$
\[\ast\ \operatorname{Antilog}_4(-3) = 4^{-3} = \frac{1}{64}\]

\subsection*{Función logarítmica y exponencial}

\textbf{Función exponencial:} \ \fbox{$f(x) = b^x$}$,\ \operatorname{Dom}f = \mathbb{R},\
b \in \mathbb{R}^+ - \{1\},\ \operatorname{Ran}f = \mathbb{R}^+$

\begin{minipage}[t]{0.485\textwidth}
\textbf{I.} $b > 1$
\begin{center}
\begin{tikzpicture}[scale=0.6]
  \ejes{-3.5}{3.5}{-0.5}{3.5}
  \draw[thick,rojo!70!black,<->,domain=-3.3:1.6,samples=60] plot (\x,{pow(2,\x)});
  \fill (0,1) circle (2pt); \node[above left,rojo] at (0,1) {\small 1};
  \draw[dashed,gray] (-0.5,-0.5) -- (3.3,3.3);
\end{tikzpicture}

\fbox{$b^{x_1} > b^{x_2} \longrightarrow x_1 > x_2$}
\end{center}
\end{minipage}\hfill
\begin{minipage}[t]{0.485\textwidth}
\textbf{II.} $0 < b < 1$
\begin{center}
\begin{tikzpicture}[scale=0.6]
  \ejes{-3.5}{3.5}{-0.5}{3.5}
  \draw[thick,rojo!70!black,<->,domain=-1.6:3.3,samples=60] plot (\x,{pow(0.5,\x)});
  \fill (0,1) circle (2pt); \node[above right,rojo] at (0,1) {\small 1};
\end{tikzpicture}

\fbox{$b^{x_1} > b^{x_2} \longrightarrow x_1 < x_2$}
\end{center}
\end{minipage}

\medskip
\textbf{Función logarítmica:} \ \fbox{$f(x) = \log_b x$}$,\ x \in \mathbb{R}^+,\
b \in \mathbb{R}^+ - \{1\},\ \operatorname{Ran}f = \mathbb{R}$

\begin{minipage}[t]{0.485\textwidth}
\textbf{I.} $b > 1$
\begin{center}
\begin{tikzpicture}[scale=0.6]
  \ejes{-0.5}{4}{-3}{2.5}
  \draw[thick,rojo!70!black,->,domain=0.12:3.6,samples=60] plot (\x,{ln(\x)/ln(2)});
  \fill (1,0) circle (2pt); \node[below right,rojo] at (1,0) {\small 1};
  \draw[dashed,gray] (-0.5,-0.5) -- (2.5,2.5) node[above] {\scriptsize $y = x$};
\end{tikzpicture}

\fbox{$\log_b x_1 > \log_b x_2 \longrightarrow x_1 > x_2$}
\end{center}
\end{minipage}\hfill
\begin{minipage}[t]{0.485\textwidth}
\textbf{II.} $0 < b < 1$
\begin{center}
\begin{tikzpicture}[scale=0.6]
  \ejes{-0.5}{4}{-3}{2.5}
  \draw[thick,rojo!70!black,->,domain=0.12:3.6,samples=60] plot (\x,{-ln(\x)/ln(2)});
  \fill (1,0) circle (2pt); \node[above right,rojo] at (1,0) {\small 1};
  \draw[dashed,gray] (-0.5,-0.5) -- (2.5,2.5) node[above] {\scriptsize $y = x$};
\end{tikzpicture}

\fbox{$\log_b x_1 > \log_b x_2 \longrightarrow x_1 < x_2$}
\end{center}
\end{minipage}

\newpage
% ------------------------------------------------------------
\section{Números complejos}

\textbf{Unidad imaginaria:} \quad \fcolorbox{magenta}{white}{$i = \sqrt{-1} = (0;1)$}

\textbf{Cantidades imaginarias:} \quad \fcolorbox{magenta}{white}{$\sqrt{-k} = \sqrt{k}\cdot i$}$;\ k > 0$

\textbf{Potencias de la unidad imaginaria:}
\[
\left.
\begin{array}{llll}
i^1 = i & i^5 = i & \cdots & i^{\overset{\circ}{4}+1} = i\\
i^2 = -1 & i^6 = -1 & \cdots & i^{\overset{\circ}{4}+2} = -1\\
i^3 = -i & i^7 = -i & \cdots & i^{\overset{\circ}{4}+3} = -i\\
i^4 = 1 & i^8 = 1 & \cdots & \fcolorbox{magenta}{white}{$i^{\overset{\circ}{4}} = 1$}
\end{array}
\right]
\;\fcolorbox{magenta}{white}{$i^{\overset{\circ}{4}+m} = i^m$}
\]

\textbf{Número complejo:} Es un par ordenado de componentes reales;
\ \fcolorbox{magenta}{white}{$Z = (a;b)$}

\qquad \textcolor{azul}{Parte real:} $\operatorname{Re}(Z) = a$ \qquad
\textcolor{azul}{Parte imaginaria:} $\operatorname{Im}(Z) = b$

\textcolor{naranja}{\textbf{Conjunto complejo:}} \quad
\fcolorbox{magenta}{white}{$\mathbb{C} = \{(a;b) \,/\, a, b \in \mathbb{R}\}$}

\subsection*{Formas de representar un número complejo}

\textbf{I. Forma binomial:}
\[
Z = (a;b) \text{ en su forma binomial } \fcolorbox{magenta}{white}{$Z = a + bi$}
\rightarrow \fcolorbox{magenta}{white}{$Z = \operatorname{Re}(Z) + \operatorname{Im}(Z)\cdot i$}
\]
\[
Z_1 = (6;1) = 6 + 1\cdot i \qquad\qquad Z_2 = (-2;\sqrt{3}) = -2 + \sqrt{3}\,i
\]

\textcolor{rojo}{\textbf{Clases de números complejos:}}

\begin{minipage}[t]{0.485\textwidth}
\textcolor{azul}{Complejo real:} \ \fcolorbox{magenta}{white}{$Z = a + 0\cdot i$}

\textcolor{azul}{Complejo imaginario puro:} \ \fcolorbox{magenta}{white}{$Z = 0 + bi$}
\end{minipage}\hfill
\begin{minipage}[t]{0.485\textwidth}
\textcolor{azul}{Número imaginario:} \ \fcolorbox{magenta}{white}{$\operatorname{Im}(Z) \neq 0$}

\textcolor{azul}{Complejo nulo:} \ \fcolorbox{magenta}{white}{$Z = (0;0) = 0 + 0\cdot i$}
\end{minipage}

\medskip
\textcolor{rojo}{\textbf{Relaciones entre complejos:}}
\begin{itemize}[label={}]
  \item \textcolor{azul}{Complejos iguales:} $Z_1 = a + bi \quad Z_2 = c + di$ \quad
  \fcolorbox{magenta}{white}{$Z_1 = Z_2 \leftrightarrow a = c \wedge b = d$}
  \item \textcolor{azul}{Complejos conjugados:} $Z = a + bi$ \quad
  \fcolorbox{magenta}{white}{$\overline{Z} = a - bi$}
  \item \textcolor{azul}{Complejos opuestos:} $Z = a + bi$ \quad
  \fcolorbox{magenta}{white}{$Z^* = -a - bi$}
\end{itemize}

\textcolor{rojo}{\textbf{Representación gráfica de un complejo:}} Sea: $Z = a + bi$

\begin{minipage}{0.436\textwidth}
\centering
\begin{tikzpicture}[scale=0.8]
  \draw[->] (-0.3,0) -- (4.5,0) node[below] {\scriptsize $x$ (eje real)};
  \draw[->] (0,-0.3) -- (0,3.6) node[above] {\scriptsize $y$ (eje imaginario)};
  \draw[thick] (0,0) -- (3,2.5) node[above right] {$(a;b)$};
  \draw[dashed] (3,0) -- (3,2.5) -- (0,2.5);
  \node at (1.3,1.4) {$r$};
  \draw (0.5,0) arc (0:40:0.5); \node at (0.75,0.25) {\scriptsize $\theta$};
  \draw[|-|] (0,-0.4) -- node[below] {$a$} (3,-0.4);
  \draw[|-|] (-0.4,0) -- node[left] {$b$} (-0.4,2.5);
  \node[below left] at (0,0) {$O$};
\end{tikzpicture}
\end{minipage}\hfill
\begin{minipage}{0.533\textwidth}
$r$: Radio vector, módulo o norma.

$\theta$: Argumento o amplitud.

De la figura: \ \fcolorbox{magenta}{white}{$r = |Z| = \sqrt{a^2 + b^2}$}
\end{minipage}

\medskip
\textcolor{rojo}{\textbf{Teoremas:}}

\begin{minipage}[t]{0.320\textwidth}
\begin{enumerate}
  \item $|Z| \geq 0$
  \item $|Z_1 + Z_2| \leq |Z_1| + |Z_2|$
  \item $|Z| = |\overline{Z}| = |Z^*|$
  \item $\dfrac{|Z_1|}{|Z_2|} = \left|\dfrac{Z_1}{Z_2}\right|$
  \item $|Z_1\cdot Z_2| = |Z_1||Z_2|$
  \item $|Z^n| = |Z|^n$
\end{enumerate}
\end{minipage}\hfill
\begin{minipage}[t]{0.359\textwidth}
\begin{enumerate}\setcounter{enumi}{6}
  \item $\operatorname{Re}(Z_1 + Z_2) = \operatorname{Re}(Z_1) + \operatorname{Re}(Z_2)$
  \item $\overline{(Z_1 \pm Z_2)} = \overline{Z_1} \pm \overline{Z_2}$
  \item $Z + \overline{Z} = 2\operatorname{Re}(Z)$
  \item $Z - \overline{Z} = 2\operatorname{Im}(Z)\cdot i$
  \item $Z\cdot\overline{Z} = |Z|^2$
  \item $(1 + i)^2 = 2i$
\end{enumerate}
\end{minipage}\hfill
\begin{minipage}[t]{0.291\textwidth}
\begin{enumerate}\setcounter{enumi}{12}
  \item $(1 - i)^2 = -2i$
  \item $(1 + i)^3 = -2 + 2i$
  \item $(1 - i)^3 = -2 - 2i$
  \item $(1 + i)^4 = (1 - i)^4 = -4$
  \item $\dfrac{1 + i}{1 - i} = i$
  \item $\dfrac{1 - i}{1 + i} = -i$
\end{enumerate}
\end{minipage}

\medskip
\textbf{II. Forma polar o trigonométrica de un complejo:} Sea: $Z = a + bi$

\begin{minipage}{0.436\textwidth}
\centering
\begin{tikzpicture}[scale=0.8]
  \draw[->] (-0.3,0) -- (4.5,0) node[right] {\small Re};
  \draw[->] (0,-0.3) -- (0,3.6) node[left] {\small Im};
  \draw[thick] (0,0) -- node[above left] {$|Z|$} (3.3,2.5) node[above] {\small Afijo};
  \fill (3.3,2.5) circle (2pt);
  \draw[dashed] (3.3,0) -- (3.3,2.5) -- (0,2.5);
  \draw (0.5,0) arc (0:37:0.5); \node at (0.75,0.25) {\scriptsize $\theta$};
  \node[below,rojo] at (1.65,0) {\small $|Z|\cos\theta$};
  \node[right,rojo] at (3.3,1.25) {\small $|Z|\operatorname{sen}\theta$};
  \node[below left] at (0,0) {\small Polo};
\end{tikzpicture}
\end{minipage}\hfill
\begin{minipage}{0.533\textwidth}
\begin{gather*}
Z = |Z|\cdot\cos\theta + |Z|\operatorname{sen}\theta\cdot i\\
\fcolorbox{magenta}{white}{$Z = |Z|(\cos\theta + i\operatorname{sen}\theta)$}\\
\fcolorbox{magenta}{white}{$Z = |Z|\operatorname{Cis}\theta$}\\
\fcolorbox{verde}{verde!15}{$\theta = \operatorname{Arg}(Z)$} \rightarrow
\fcolorbox{verde}{verde!15}{$\tan\theta = \dfrac{b}{a}$}
\end{gather*}
\end{minipage}

\medskip
\textbf{III. Forma exponencial de un número complejo:} \quad
\fcolorbox{magenta}{white}{$e^{\theta i} = \cos\theta + i\operatorname{sen}\theta$}
\ \textcolor{azul}{\small (Teorema de Euler)}

El número complejo no nulo se puede escribir como: \quad
\fcolorbox{magenta}{white}{$Z = |Z|\cdot e^{\theta i}$}


% ============================================================
% Páginas 70-72
% ============================================================
\textcolor{magenta}{$\theta$}: Argumento en radianes. \qquad
\textcolor{magenta}{$e$}: Número de Neper: $2{,}71828\ldots$

\textcolor{rojo}{\textbf{Operaciones en forma polar o exponencial:}} Dados
$Z_1 = r_1\cdot\operatorname{Cis}\theta;\ Z_2 = r_2\cdot\operatorname{Cis}\alpha$
\begin{enumerate}
  \item \textcolor{azul}{Multiplicación:} \
  \fcolorbox{magenta}{white}{$Z_1\cdot Z_2 = r_1\cdot r_2\operatorname{Cis}(\theta + \alpha)$}
  \item \textcolor{azul}{División:} \
  \fcolorbox{magenta}{white}{$\dfrac{Z_1}{Z_2} = \dfrac{r_1}{r_2}\operatorname{Cis}(\theta - \alpha)$}
  \item \textcolor{azul}{Potenciación:} Teorema de De Moivre

  Dado: \ \fcolorbox{magenta}{white}{$Z = r\cdot\operatorname{Cis}\theta \rightarrow Z^n = r^n\operatorname{Cis}(n\cdot\theta)$}
  \item \textcolor{azul}{Raíz enésima de un complejo:} Todo complejo no nulo admite ``$n$''
  raíces distintas dadas por:
  \[
  \fcolorbox{magenta}{white}{$\sqrt[n]{Z} = \sqrt[n]{r}\operatorname{Cis}\left(\dfrac{\theta + 2\pi k}{n}\right)$}
  \quad \{k = 0; 1; 2; \ldots; (n-1)\}
  \]
  \item \textcolor{azul}{Teorema:} Si tenemos dos números complejos en forma polar:
  $Z_1 = \operatorname{Cis}\theta;\ Z_2 = \operatorname{Cis}\alpha$, entonces:
  \[
  \fcolorbox{magenta}{white}{$\operatorname{Cis}\theta + \operatorname{Cis}\alpha = 2\cos\left(\dfrac{\theta - \alpha}{2}\right)\operatorname{Cis}\left(\dfrac{\theta + \alpha}{2}\right)$}
  \]
\end{enumerate}

\textcolor{rojo}{\textbf{Raíces cúbicas de la unidad real:}}
\[Z = 1 \rightarrow \sqrt[3]{Z} = \operatorname{Cis}\frac{2\pi k}{3} = \cos\frac{2\pi k}{3} + i\operatorname{sen}\frac{2\pi k}{3}\]
\begin{align*}
k = 0 &\rightarrow \sqrt[3]{Z} = \cos 0 + i\cdot\operatorname{sen}0 = 1 & W_1 &= 1\\
k = 1 &\rightarrow \sqrt[3]{Z} = \cos\frac{2\pi}{3} + i\cdot\operatorname{sen}\frac{2\pi}{3} = -\frac{1}{2} + \frac{\sqrt{3}}{2}i & W_2 &= -\frac{1}{2} + \frac{\sqrt{3}}{2}i\\
k = 2 &\rightarrow \sqrt[3]{Z} = \cos\frac{4\pi}{3} + i\cdot\operatorname{sen}\frac{4\pi}{3} = -\frac{1}{2} - \frac{\sqrt{3}}{2}i & W_3 &= -\frac{1}{2} - \frac{\sqrt{3}}{2}i
\end{align*}

\textcolor{rojo}{\textbf{Propiedades de las raíces cúbicas de la unidad:}}
\begin{enumerate}
  \item Si una raíz compleja es ``$W$'' la otra es ``$W^2$'' y la tercera es 1, así:
  $\sqrt[3]{1} = 1, W, W^2$
\end{enumerate}
\[
\text{2. } 1 + W + W^2 = 0 \qquad
\text{3. } W^3 = 1 \rightarrow W^{3k} = 1 \qquad
\text{4. } W^{3k + r} = W^r
\]

\begin{center}
\begin{tikzpicture}[scale=1.1]
  \begin{scope}
    \draw[->] (-1.6,0) -- (1.6,0); \draw[->] (0,-1.5) -- (0,1.6);
    \draw (0,0) circle (1);
    \coordinate (A) at (0:1); \coordinate (B) at (120:1); \coordinate (C) at (240:1);
    \draw[magenta,thick] (A) -- (B) -- (C) -- cycle;
    \foreach \p in {A,B,C} \draw[->,graf,thick] (0,0) -- (\p);
    \node[right] at (A) {\scriptsize 1};
    \node[above left] at (B) {\scriptsize $-\frac{1}{2} + \frac{\sqrt{3}}{2}i$};
    \node[below left] at (C) {\scriptsize $-\frac{1}{2} - \frac{\sqrt{3}}{2}i$};
    \node[azul] at (1.2,1.6) {\small Raíces cúbicas de la unidad};
  \end{scope}
  \begin{scope}[xshift=5.5cm]
    \draw[->] (-1.6,0) -- (1.6,0); \draw[->] (0,-1.5) -- (0,1.6);
    \draw (0,0) circle (1);
    \draw[magenta,thick] (1,0) -- (0,1) -- (-1,0) -- (0,-1) -- cycle;
    \foreach \a in {0,90,180,270} \draw[->,graf,thick] (0,0) -- (\a:1);
    \node[above right] at (1,0) {\scriptsize 1}; \node[above left] at (0,1) {\scriptsize $i$};
    \node[above left] at (-1,0) {\scriptsize $-1$}; \node[below left] at (0,-1) {\scriptsize $-i$};
    \node[azul] at (1.2,1.6) {\small Raíces cuartas de la unidad};
  \end{scope}
\end{tikzpicture}
\end{center}

\textcolor{rojo}{\textbf{Teorema:}} Los afijos de las raíces $n$-ésimas de un número complejo
son los vértices de un polígono regular de $n$ lados.

\begin{minipage}{0.436\textwidth}
\centering
\begin{tikzpicture}[scale=1.4]
  \draw[->] (-1.4,0) -- (1.5,0) node[right] {$x$};
  \draw[->] (0,-1.2) -- (0,1.4) node[left] {$y$};
  \draw[graf] (0,0) circle (1);
  \foreach \k in {0,...,11} \draw (0,0) -- ({20+30*\k}:1);
  \foreach \k/\l in {0/Z_0,1/Z_1,2/Z_2,3/Z_3,4/Z_4,5/Z_5} \node at ({20+30*\k}:1.18) {\scriptsize $\l$};
  \node at (-10:1.25) {\scriptsize $Z_{n-1}$};
  \draw[->,azul] (0.3,0) arc (0:20:0.3); \node[azul] at (0.45,0.08) {\scriptsize $\theta$};
\end{tikzpicture}
\end{minipage}\hfill
\begin{minipage}{0.533\textwidth}
Del gráfico se observa:
\[\theta = \frac{360^\circ}{n} \;\Leftrightarrow\; \frac{2\pi\,\text{rad}}{n}\]
El área de un polígono de $n$ lados:
\[S = \frac{n|Z_0|^2}{2}\operatorname{sen}\left(\frac{2\pi}{n}\right)\]
$Z_0$: Es una de las raíces de $Z$.

También se cumple:
\begin{gather*}
Z_0^n = Z_1^n = Z_2^n = Z_3^n = \cdots = Z_{n-1}^n = Z\\
Z_0 + Z_1 + Z_2 + Z_3 + \cdots + Z_{n-1} = 0
\end{gather*}
\end{minipage}

\newpage
% ------------------------------------------------------------
\section{MCD, MCM y fracciones}

\textbf{Factor de un polinomio:} \quad $P(x) \equiv Q(x)\cdot H(x)$; \quad $H(x)$: polinomio no nulo.

\textbf{Factor común de dos o más polinomios:}
\[P(x) = M(x)\cdot f(x);\quad Q(x) = M(x)\cdot g(x);\quad \text{Factor común: } M(x)\]
\begin{gather*}
P(x) = (2x - 1)^2(x - 3)^3(5x + 2)\\
Q(x) = (x - 2)(2x - 1)(5x + 2)^5
\end{gather*}
\begin{center}
Factores comunes: $(2x - 1)$, $(5x + 2)$, $(2x - 1)(5x + 2)$
\end{center}

\textbf{Múltiplo de un polinomio:}
\[P(x) = (x + 2)(x - 5) \qquad \text{Múltiplos: } (x+2)(x-5),\ (x+2)^2(x-5),\ (x+2)(x-5)x, \ldots\]

\textbf{Polinomio múltiplo común:}
\begin{gather*}
P(x) = (x + 1)(x^2 + 3) \qquad Q(x) = (x - 1)(x + 1)\\
\text{Múltiplos de los dos: } (x-1)(x+1)(x^2+3),\ (x-1)^2(x+1)(x^2+3),\ (x-1)(x+1)^2(x^2+3)^3, \ldots
\end{gather*}

\textbf{Máximo común divisor (MCD):}

\begin{minipage}{0.485\textwidth}
\begin{gather*}
P(x) = (2x + 7)^4(x - 1)^2(3x - 1)\\
Q(x) = (2x - 1)^5(3x - 1)^2(x - 1)^2
\end{gather*}
\end{minipage}\hfill
\begin{minipage}{0.485\textwidth}
\[\operatorname{MCD}(P;Q) = (x - 1)^2(3x - 1)\]
\end{minipage}

\textbf{Mínimo común múltiplo (MCM):}

\begin{minipage}{0.485\textwidth}
\begin{gather*}
P(x) = (2x - 1)(x + 1)^3\\
Q(x) = (x - 1)(x + 1)^2
\end{gather*}
\end{minipage}\hfill
\begin{minipage}{0.485\textwidth}
\[\operatorname{MCM}(P;Q) = (2x - 1)(x - 1)(x + 1)^3\]
\end{minipage}

\textcolor{rojo}{\textbf{Teorema:}} \quad $P(x)\cdot Q(x) = \operatorname{MCD}(P;Q)\cdot\operatorname{MCM}(P;Q)$

\textbf{Fracciones \hl{algebraicas}:} \quad
$f(x) = \dfrac{N(x)}{D(x)}$ \quad
\begin{tabular}{l}
$N(x)$: Numerador\\
$D(x)$: Denominador (no nulo ni constante)
\end{tabular}

\begin{itemize}[label=\checkmark]
  \item $N(x)$ y $D(x)$ son expresiones algebraicas racionales enteras.
  \item NO es una división exacta.
  \item $\text{CVA}(f) = \mathbb{U} - \{D(x) = 0\}$
\end{itemize}


% ============================================================
% Páginas 73-75
% ============================================================
\textbf{Fracciones algebraicas reducibles:}
\[f(x) = \frac{N(x)}{D(x)}: \text{ reducible si } N(x) \wedge D(x) \text{ poseen factores comunes.}\]

\textcolor{naranja}{\textbf{Clasificación de fracciones algebraicas:}}
\begin{itemize}[label={}]
  \item \textcolor{verde}{Propias:} $[N(x)]^\circ < [D(x)]^\circ$ \hfill
  $\dfrac{x+1}{x^2+2};\ \dfrac{1}{x^5}$ \hspace*{3cm}
  \item \textcolor{verde}{Impropias:} $[N(x)]^\circ > [D(x)]^\circ$ \hfill
  $\dfrac{x^2-1}{x};\ \dfrac{x^2+2x+4}{x-3}$ \hspace*{3cm}
  \item \textcolor{verde}{Homogéneas:} $\dfrac{1}{x+1};\ \dfrac{x+2}{x+1};\ \dfrac{x^2}{x+1}$
  \item \textcolor{verde}{Heterogéneas:} $\dfrac{2}{x+2};\ \dfrac{5x}{x+3};\ \dfrac{6}{x+1}$
\end{itemize}

\textcolor{magenta}{\textbf{Teorema:}}
\[
\text{Si } f(x;y) = \frac{a_1x + b_1xy + c_1y + d_1}{a_2x + b_2xy + c_2y + d_2}
\text{ no depende ni de $x$ ni de $y$} \rightarrow
\frac{a_1}{a_2} = \frac{b_1}{b_2} = \frac{c_1}{c_2} = \frac{d_1}{d_2} = k
\]
\[a_2 \neq 0;\ b_2 \neq 0;\ c_2 \neq 0;\ d_2 \neq 0\]

\subsection*{Descomposición en fracciones parciales}

\textcolor{azul}{\textbf{Caso I:}} Para fracciones propias; el $D(x)$ sólo posee factores de la
forma ``$ax + b$''.
\[
\frac{3x+4}{x^2+3x+2} = \frac{3x+4}{(x+1)(x+2)} = \frac{A}{(x+1)} + \frac{B}{(x+2)}
= \frac{(A+B)x + (2A+B)}{(x+1)(x+2)} = \frac{1}{x+1} + \frac{2}{x+2}
\]
\[A + B = 3;\quad 2A + B = 4;\quad A = 1;\ B = 2\]

El $D(x)$ posee factores de la forma $(ax+b)^n \wedge a \neq 0$
\[
\frac{3x+1}{(x-1)^2} = \frac{A}{(x-1)} + \frac{B}{(x-1)^2} = \frac{Ax + (B-A)}{(x-1)^2}
= \frac{3}{x-1} + \frac{4}{(x-1)^2}
\]
\begin{itemize}[label=\checkmark]
  \item Los dos anteriores se pueden mezclar.
  \item No necesariamente es para ``$ax + b$'', también lo es para ``$ax^2 + bx + c$'' y
  ``$ax^3 + bx^2 + cx + d$'' mientras sean no factorizables.
\end{itemize}

\textcolor{azul}{\textbf{Caso II:}} Para fracciones impropias.
\[
f(x) = \frac{N(x)}{D(x)} = q(x) + \underbrace{\frac{R(x)}{D(x)}}_{\substack{\text{Fracción}\\ \text{propia}}}
\;:\; \text{Se trabaja como el caso I}
\]

\newpage
% ------------------------------------------------------------
\section{Sucesiones}

\textbf{Sucesión de números reales:} Es de la forma: \
$x: \mathbb{N} \to \mathbb{R};\ n \to x(n)$
\[\ovalbox{$x = \{(n, x(n)) \mid n \in \mathbb{N}\}$ \ además \ $x(n) = x_n$}\]

\textcolor{naranja}{\textbf{Convergencia:}} Se dice que una sucesión $\{x_n\}_{n\in\mathbb{N}}$
es una sucesión convergente si existe el siguiente límite:
\[
\lim_{n\to+\infty} x_n = L \in \mathbb{R}\;; \qquad
\forall \varepsilon > 0,\ \exists n_0 \in \mathbb{N},\ \forall n \geq n_0 \rightarrow |x_n - L| < \varepsilon
\]

\textcolor{magenta}{\textbf{Ejemplos:}}
\[
\lim_{n\to+\infty}\frac{1}{n} = 0 \qquad
\lim_{n\to+\infty}\sqrt[n]{n} = 1 \qquad
\lim_{n\to+\infty}\sqrt[n]{a} = 1 \qquad
\lim_{n\to+\infty}a^n = 0 \ \ \textcolor{rojo!60!black}{(0 \leq a < 1)}
\]

\textbf{Teorema:} Sean $p(n)$ y $q(n)$ polinomios no nulos. Se tiene que:
\begin{itemize}
  \item $\operatorname{gr}(p) = \operatorname{gr}(q)$
  \[\lim_{n\to+\infty}\frac{p(n)}{q(n)} = \frac{\text{Coeficiente principal de } p}{\text{Coeficiente principal de } q}\]
\end{itemize}
\begin{minipage}[t]{0.485\textwidth}
\begin{itemize}
  \item $\operatorname{gr}(p) < \operatorname{gr}(q)$
  \[\lim_{n\to+\infty}\frac{p(n)}{q(n)} = 0\]
\end{itemize}
\end{minipage}\hfill
\begin{minipage}[t]{0.485\textwidth}
\begin{itemize}
  \item $\operatorname{gr}(p) > \operatorname{gr}(q)$
  \[\lim_{n\to+\infty}\frac{p(n)}{q(n)} = +\infty\ (-\infty)\]
\end{itemize}
\end{minipage}

\textcolor{verde}{\textbf{Teorema del Sándwich:}} Si $x_n \leq y_n \leq z_n,\ \forall n \in \mathbb{N}$ y
\[
\lim_{n\to+\infty}x_n = \lim_{n\to+\infty}z_n = L \quad \text{entonces:} \quad \lim_{n\to+\infty}y_n = L
\]

\textcolor{naranja}{\textbf{Acotada:}} Se dice que $\{x_n\}_{n\in\mathbb{N}}$ es una sucesión
acotada si
\[
\exists K > 0,\ \forall n \in \mathbb{N},\ |x_n| \leq K \quad \text{o también:} \quad
\exists a, b \in \mathbb{R},\ \forall n \in \mathbb{N},\ a \leq x_n \leq b
\]
\begin{minipage}[t]{0.485\textwidth}
\textcolor{azul}{Acotada inferiormente:}
\[\exists a \in \mathbb{R},\ \forall n \in \mathbb{N},\ a \leq x_n\]
\end{minipage}\hfill
\begin{minipage}[t]{0.485\textwidth}
\textcolor{azul}{Acotada superiormente:}
\[\exists b \in \mathbb{R},\ \forall n \in \mathbb{N},\ x_n \leq b\]
\end{minipage}


% ============================================================
% Páginas 76-79
% ============================================================
\textbf{Teorema:} Si una sucesión es convergente, entonces la sucesión es acotada.

\textbf{Nota:} El recíproco del teorema es falso.

\textcolor{magenta}{\textbf{Sucesión monótona:}} Se dice que la sucesión
$\{x_n\}_{n\in\mathbb{N}}$ es una sucesión monótona si cumple alguno de los siguientes casos.
\[
\begin{array}{lll}
\bullet & \textcolor{rojo}{\text{Estrictamente creciente}} & \forall n \in \mathbb{N}\quad x_n < x_{n+1}\\
\bullet & \textcolor{rojo}{\text{Estrictamente decreciente}} & \forall n \in \mathbb{N}\quad x_n > x_{n+1}\\
\bullet & \textcolor{rojo}{\text{Creciente}} & \forall n \in \mathbb{N}\quad x_n \leq x_{n+1}\\
\bullet & \textcolor{rojo}{\text{Decreciente}} & \forall n \in \mathbb{N}\quad x_n \geq x_{n+1}
\end{array}
\]

\textbf{Observación \#1:} Para analizar la monotonía se calcula la siguiente expresión:
\[M = x_{n+1} - x_n\]
\[
\begin{array}{lll}
\bullet & \text{Si } M > 0, & \text{Estrictamente creciente}\\
\bullet & \text{Si } M < 0, & \text{Estrictamente decreciente}\\
\bullet & \text{Si } M \geq 0, & \text{Creciente}\\
\bullet & \text{Si } M \leq 0, & \text{Decreciente}
\end{array}
\]

\textbf{Observación \#2:} Para analizar la monotonía se calcula la siguiente expresión:
\[M = \frac{x_{n+1}}{x_n}\;; \quad \text{donde \hl{los términos de la sucesión son positivos.}}\]
\[
\begin{array}{lll}
\bullet & \text{Si } M > 1, & \text{Estrictamente creciente.}\\
\bullet & \text{Si } M < 1, & \text{Estrictamente decreciente.}\\
\bullet & \text{Si } M \geq 1, & \text{Creciente.}\\
\bullet & \text{Si } M \leq 1, & \text{Decreciente.}
\end{array}
\]

\textbf{Teorema:} Si una sucesión es acotada y monótona, entonces la sucesión es convergente.
\begin{itemize}
  \item Si $\{x_n\}_{n\in\mathbb{N}}$ es creciente (o E. creciente) y acotada superiormente,
  entonces es convergente.
  \item Si $\{x_n\}_{n\in\mathbb{N}}$ es decreciente (o E. decreciente) y acotada
  inferiormente, entonces es convergente.
\end{itemize}

\textbf{Teorema:} Si $\displaystyle\lim_{n\to+\infty}x_n = L$, entonces:
\[
\bullet\ \lim_{n\to+\infty}\frac{x_1 + x_2 + \cdots + x_n}{n} = L
\qquad\qquad
\bullet\ \lim_{n\to+\infty}\sqrt[n]{x_1\cdot x_2\cdots x_n} = L
\]

\textcolor{magenta}{\textbf{El número neperiano:}}

Se tiene que la sucesión es acotada y monótona:
\[
x_n = \left(1 + \frac{1}{n}\right)^n \qquad\qquad
\lim_{n\to+\infty}\left(1 + \frac{1}{n}\right)^n = e
\]
En general: \quad $\displaystyle\lim_{n\to+\infty}\left(1 + \frac{x}{n}\right)^n = e^x$

\textcolor{verde}{\textbf{Teorema:}} Si $\displaystyle\lim_{n\to+\infty}a_n = 1$ y
$\displaystyle\lim_{n\to+\infty}b_n = +\infty$, entonces:
\[\lim_{n\to+\infty}(a_n)^{b_n} = e^{\lim\limits_{n\to+\infty}(a_n - 1)b_n}\]

\textcolor{magenta}{\textbf{Criterio de la raíz o razón:}}

Sea \quad $\displaystyle\lim_{n\to+\infty}\left|\frac{a_{n+1}}{a_n}\right| = L$ \quad o \quad
$\displaystyle\lim_{n\to+\infty}\sqrt[n]{|a_n|} = L$
\begin{enumerate}[label=\cn{\arabic*}]
  \item Si $0 \leq L < 1$, entonces $\displaystyle\lim_{n\to+\infty}a_n = 0$
  \item Si $L > 1$, entonces la sucesión $\{a_n\}_{n\in\mathbb{N}}$ es una sucesión divergente.
\end{enumerate}

\newpage
% ------------------------------------------------------------
\section{Series}

Una serie $\{S_n\}_{n\in\mathbb{N}}$ es una sucesión con la siguiente regla de correspondencia
\[S_n = \sum_{k=1}^{n}a_k\]

\textbf{Convergencia:} Se dice que $\{S_n\}_{n\in\mathbb{N}}$ es una serie convergente si existe
el siguiente límite.
\[\lim_{n\to+\infty}S_n = L \in \mathbb{R}\;: \quad \textcolor{magenta}{L = \sum_{k=1}^{\infty}a_k}\]

\textbf{Series geométricas:} Sea $|R| < 1$, se cumple
\begin{itemize}
  \item $A_n = \displaystyle\sum_{k=0}^{n}R^k \qquad \lim_{n\to+\infty}A_n = \frac{1}{1 - R}$
  \item $A_n = \displaystyle\sum_{k=m}^{n}R^k \qquad \lim_{n\to+\infty}A_n = \frac{R^m}{1 - R}$
\end{itemize}

\textbf{Nota:} Si $\displaystyle\lim_{n\to+\infty}a_n = 0$, entonces \quad
\fbox{$\displaystyle\lim_{n\to+\infty}\sum_{k=1}^{n}(a_k - a_{k+1}) = a_1$}

\textbf{Teorema:} Si la serie $S_n = \displaystyle\sum_{k=1}^{n}a_k$ es convergente, entonces
\[\lim_{n\to+\infty}a_n = 0\]
\begin{itemize}
  \item El recíproco es falso.
\end{itemize}

\textbf{Serie $n^p$:} Sea $p \in \mathbb{R}$ y la serie:

\begin{minipage}{0.436\textwidth}
\[S_k = \sum_{n=1}^{k}\frac{1}{n^p}\]
\end{minipage}\hfill
\begin{minipage}{0.533\textwidth}
Se tiene que:
\begin{itemize}
  \item Si $p \leq 1$, entonces $S_k$ diverge.
  \item Si $p > 1$, entonces $S_k$ converge.
\end{itemize}
\end{minipage}

\textbf{Comparación:} Sean $\forall n \in \mathbb{N},\ 0 \leq a_n \leq b_n$
\begin{itemize}
  \item Si la serie $\displaystyle\sum_{k=1}^{n}b_k$ es convergente entonces
  $\displaystyle\sum_{k=1}^{n}a_k$ converge.
  \item Si la serie $\displaystyle\sum_{k=1}^{n}a_k$ es divergente entonces
  $\displaystyle\sum_{k=1}^{n}b_k$ diverge.
\end{itemize}

\textcolor{verde}{\textbf{Teorema de Leibnitz:}} Si $\displaystyle\lim_{n\to+\infty}a_n = 0$ y
$\forall n \in \mathbb{N}$ y $a_n \geq a_{n+1} \geq 0$
\[\text{entonces:} \quad \sum_{k=1}^{n}(-1)^k a_k \text{ es una serie convergente}\]

\textcolor{magenta}{\textbf{Criterio de la raíz o razón:}}

Sea la serie $S_n = \displaystyle\sum_{k=1}^{n}a_k$
\[
\lim_{k\to+\infty}\left|\frac{a_{k+1}}{a_k}\right| = L \quad \text{o} \quad
\lim_{k\to+\infty}\sqrt[k]{|a_k|} = L
\]
\begin{itemize}
  \item Si $0 \leq L < 1$, entonces la serie $S_n$ es convergente.
  \item Si $L > 1$, entonces la serie $S_n$ es divergente.
\end{itemize}

\textbf{Observación:}
\[\sum_{k=0}^{\infty}\frac{1}{k!} = e \qquad\qquad \sum_{k=0}^{\infty}\frac{x^k}{k!} = e^x\]

\textcolor{magenta}{\textbf{Sucesión de Fibonacci:}}

Sea: \quad $t_{n+2} = t_{n+1} + t_n$ \quad con $t_1 = 1,\ t_2 = 1$
\[
t_n = \frac{1}{\sqrt{5}}\left(\frac{1 + \sqrt{5}}{2}\right)^n - \frac{1}{\sqrt{5}}\left(\frac{1 - \sqrt{5}}{2}\right)^n
\]

\end{document}
